\documentclass[10pt,a4paper]{article}
\usepackage[utf8]{inputenc}
\usepackage[T1]{fontenc}
\usepackage[spanish,es-tabla,es-noquoting]{babel}
\usepackage{lmodern}
\usepackage[margin=1.9cm]{geometry}
\usepackage{longtable,booktabs,array}
\usepackage{enumitem}
\usepackage{xcolor}
\usepackage{titlesec}
\usepackage{fancyhdr}
\usepackage[hidelinks]{hyperref}
\definecolor{eje}{HTML}{7A1F2B}
\definecolor{sub}{HTML}{1F4E79}
\setlist{nosep,leftmargin=1.4em,topsep=2pt}
\setlist[itemize,1]{label=\textbullet}
\setlist[itemize,2]{label=--}
\setlist[itemize,3]{label=$\circ$}
\titleformat{\section}{\LARGE\bfseries\sffamily\color{eje}}{}{0pt}{}[\vspace{-4pt}\rule{\linewidth}{1pt}]
\titleformat{\subsection}{\Large\bfseries\sffamily\color{sub}}{}{0pt}{}
\titleformat{\subsubsection}{\large\bfseries\sffamily}{}{0pt}{}
\titlespacing*{\subsection}{0pt}{14pt}{6pt}
\titlespacing*{\subsubsection}{0pt}{10pt}{4pt}
\setcounter{secnumdepth}{0}
\setcounter{tocdepth}{2}
\newenvironment{esquema}{\par\medskip\begingroup\small\ttfamily\color{sub}}{\endgroup\medskip}
\pagestyle{fancy}\fancyhf{}
\fancyhead[L]{\small\sffamily DSI}
\fancyhead[R]{\small\sffamily\nouppercase{\leftmark}}
\fancyfoot[C]{\small\thepage}
\renewcommand{\sectionmark}[1]{\markboth{#1}{}}
\setlength{\headheight}{14pt}
\setlength{\parindent}{0pt}
\setlength{\emergencystretch}{3em}
\begin{document}
\begin{titlepage}\centering\vspace*{5cm}
{\Huge\bfseries\sffamily\color{eje} Doctrina Social de la Iglesia\par}\vspace{0.8cm}
{\LARGE\sffamily Resumen completo para el examen final\par}\vspace{0.6cm}
{\large\sffamily Ejes 1 a 5 --- Programa 2026\par}\vspace{2cm}
{\large Licenciatura en Ciencias de Datos --- UCASAL\par}
\end{titlepage}
\tableofcontents

\begin{quote}\itshape Resumen de los cinco ejes del programa 2026 (UCASAL, Licenciatura en Ciencias de Datos), armado a partir de la bibliografía de la cátedra. Sirve para el examen final: opción múltiple y preguntas a desarrollar.\end{quote}

\clearpage\section{Mapa del programa}

{\small\renewcommand{\arraystretch}{1.25}
\begin{longtable}{>{\raggedright\arraybackslash}p{0.045\linewidth}>{\raggedright\arraybackslash}p{0.320\linewidth}>{\raggedright\arraybackslash}p{0.546\linewidth}}
\toprule
\textbf{Eje} & \textbf{Tema} & \textbf{Bibliografía resumida} \\ \midrule \endhead
1 & El discurso social de la Iglesia & Dilexi te cap. IV; Dilexit nos cap. I; Sorge I, V, VI; Compendio Introducción y cap. II \\ \addlinespace[2pt]
2 & Del discurso a la propuesta social & Sorge VII, IX, X, XXIX; Compendio cap. IV y VII \\ \addlinespace[2pt]
3 & La presencia social de la Iglesia & Fratelli tutti V y VI; Sorge XIII y XXVIII; Compendio cap. VIII \\ \addlinespace[2pt]
4 & Debates sociales: diálogo con las ciencias & Compendio cap. III; Dignitas infinita; Laudato Si' cap. III; Sorge XXI y XXIII; Faggioni (transhumanismo) \\ \addlinespace[2pt]
5 & Aportes del papa Francisco & Camacho; Laudato Si' cap. IV; Fratelli tutti I y VII \\ \addlinespace[2pt]
\bottomrule\end{longtable}}

\noindent \textbf{Hilo conductor de la materia:} discurso social (qué dice la Iglesia y por qué puede decirlo) $\rightarrow$ propuesta social (principios y economía) $\rightarrow$ presencia social (política, laicos, diálogo) $\rightarrow$ debates actuales (persona, dignidad, ciencia y técnica) $\rightarrow$ aportes de Francisco (ecología integral y fraternidad).\par\smallskip



\clearpage\section{EJE N\textordmasculine{} 1 --- El discurso social de la Iglesia}

\subsection{1.1 León XIV, \textit{Dilexi te} (2025) --- Cap. IV: ``Una historia que continúa''}

\noindent \textbf{Idea central del capítulo:} repasa cómo, a lo largo de los últimos 150 años, el Magisterio pontificio y episcopal fue profundizando la relación de la Iglesia con los pobres, hasta llegar a reconocerlos no como objetos de asistencia sino como \textbf{sujetos} activos de la fe y de la evangelización.\par\smallskip

\subsubsection{A) Raíz popular de la Doctrina Social (\S{}82)}
\begin{itemize}
\item Los últimos dos siglos de cambios tecnológicos y sociales fueron vividos, pero también \textbf{pensados y afrontados por los pobres mismos}.
\item Los movimientos obreros, de mujeres, de jóvenes y la lucha antirracista generaron una nueva conciencia de la dignidad de los marginados.
\item La DSI tiene una \textbf{raíz popular}: no se entiende sin los laicos cristianos (y religiosos/as) que enfrentaron los desafíos concretos de su época.
\item Hoy es necesaria la interacción constante entre bautizados y Magisterio, ciudadanos y expertos, pueblo e instituciones. Se reafirma que \textbf{la realidad se ve mejor desde los márgenes}.
\end{itemize}

\subsubsection{B) El Magisterio pontificio sobre los pobres (\S{}83--89)}

{\small\renewcommand{\arraystretch}{1.25}
\begin{longtable}{>{\raggedright\arraybackslash}p{0.226\linewidth}>{\raggedright\arraybackslash}p{0.275\linewidth}>{\raggedright\arraybackslash}p{0.410\linewidth}}
\toprule
\textbf{Papa} & \textbf{Documento} & \textbf{Aporte clave} \\ \midrule \endhead
León XIII & \textit{Rerum novarum} (1891) & Denunció la situación intolerable de los obreros industriales; propuso un orden social justo. \\ \addlinespace[2pt]
San Juan XXIII & \textit{Mater et Magistra} (1961) & Justicia de dimensión mundial: los países ricos deben socorrer a los oprimidos por el hambre. \\ \addlinespace[2pt]
San Juan XXIII & Radiomensaje previo al Concilio (11-set-1962) & ``La Iglesia... como Iglesia de todos, en particular como la Iglesia de los pobres.'' \\ \addlinespace[2pt]
San Pablo VI & Apertura 2\textordfeminine{} sesión del Concilio / Audiencia 1964 & La Iglesia mira con interés a todo el que sufre; ``el pobre es representante de Cristo''. \\ \addlinespace[2pt]
Concilio Vaticano II (\textit{Gaudium et spes}) & Constitución pastoral & \textbf{Destino universal de los bienes}: la propiedad privada tiene función social; en necesidad extrema hay derecho a tomar de la riqueza ajena lo necesario. \\ \addlinespace[2pt]
San Pablo VI & \textit{Populorum progressio} & Nadie puede reservarse en exclusiva lo que supera su necesidad mientras a otros les falta lo necesario. \\ \addlinespace[2pt]
San Juan Pablo II & \textit{Sollicitudo rei socialis} & Consolida la \textbf{opción preferencial por los pobres} como forma especial de primacía de la caridad cristiana. \\ \addlinespace[2pt]
San Juan Pablo II & \textit{Laborem exercens} & El trabajo humano es ``la clave, quizá la clave esencial, de toda la cuestión social''. \\ \addlinespace[2pt]
Benedicto XVI & \textit{Caritas in veritate} & El hambre no es solo escasez material sino falta de instituciones económicas adecuadas; lectura más ``política'' de la crisis. \\ \addlinespace[2pt]
\bottomrule\end{longtable}}

\noindent \textbf{Concilio Vaticano II (\S{}84):} aunque el tema de los pobres era marginal en los documentos preparatorios, cobró centralidad gracias a obispos, teólogos y expertos (destaca la intervención del cardenal Lercaro, 6-dic-1962: ``el misterio de Cristo en la Iglesia es... el misterio de Cristo en los pobres''). Se perfila una Iglesia \textbf{más sencilla y sobria}, más parecida a su Señor que a las potencias mundanas.\par\smallskip

\subsubsection{C) El aporte latinoamericano (\S{}89)}
\begin{itemize}
\item El Papa Francisco reconoció el compromiso particular del episcopado latinoamericano.
\item \textbf{Martirio de san Óscar Romero} (arzobispo de San Salvador): testimonio vivo de identificación con el pueblo pobre.
\item Conferencias del Episcopado Latinoamericano: \textbf{Medellín, Puebla, Santo Domingo y Aparecida}, etapas clave para toda la Iglesia.
\item León XIV, misionero en Perú, se reconoce heredero de este camino.
\end{itemize}

\subsubsection{D) Estructuras de pecado que causan pobreza y desigualdades extremas (\S{}90--98)}
\begin{itemize}
\item \textbf{Medellín (\S{}90):} opción preferencial por los pobres; Cristo ``siendo rico se hizo pobre''. La Iglesia debe no solo compartir la condición de los pobres sino \textbf{ponerse de su lado}.
\item \textbf{Puebla:} confirma la opción de Medellín y califica las estructuras injustas como \textbf{``pecado social''}.
\item \textbf{\S{}91--92:} la caridad es fuerza histórica de cambio; hay que resolver las \textbf{causas estructurales} de la pobreza. Crítica a la ``dictadura de una economía que mata'': el mercado sin control estatal genera una ``tiranía invisible''.
\item \textbf{\S{}93 (\textit{Dilexit nos}, Francisco):} el pecado social se convierte en \textbf{``estructura de pecado''} cuando una mentalidad dominante normaliza el egoísmo y la indiferencia (``alienación social''). Se buscan excusas teóricas en lugar de resolver los problemas concretos.
\item \textbf{\S{}94:} urge resolver las causas estructurales; los planes asistenciales solo deben ser respuestas \textbf{pasajeras}, no la solución de fondo.
\item \textbf{\S{}95:} crítica al modelo ``exitista'' y ``privatista'' que no invierte en los más débiles; pregunta ética central: ¿los menos dotados valen menos como personas?
\item \textbf{\S{}96:} las cuestiones estructurales incluyen el hábitat urbano: ciudades que integran vs. ciudades que excluyen; vínculo entre degradación ambiental y social (afecta más a los más débiles --- eco de \textit{Laudato si'}).
\item \textbf{\S{}97--98:} responsabilidad de todo el pueblo de Dios de denunciar la injusticia, aun a costa de parecer ``estúpidos''; el Evangelio no es solo relación individual con Dios sino anuncio del \textbf{Reino}, con consecuencias sociales (justicia, paz, fraternidad).
\end{itemize}

\subsubsection{E) Los pobres como sujetos (\S{}99--102)}
\begin{itemize}
\item \textbf{Aparecida (\S{}99--100):} la opción preferencial por los pobres está implícita en la fe cristológica (Dios que se hizo pobre para enriquecernos). Se insiste en ver a las comunidades marginadas como \textbf{sujetos} capaces de crear su propia cultura y vivir el Evangelio a su modo --- \textbf{no como objetos de beneficencia}.
\item La experiencia de la pobreza da una capacidad de ver la realidad que otros no tienen; la Iglesia debe valorar la fe ``popular''.
\item \textbf{\S{}101:} la opción por los pobres exige una \textbf{``atención amante''}: valorar al pobre en su bondad propia, su cultura, su modo de vivir la fe --- no por lástima sino porque ``es bello''. Agradecimiento a quienes eligen vivir \textit{entre} los pobres y \textit{como} ellos.
\item \textbf{\S{}102:} hay que dejarse \textbf{evangelizar por los pobres}, reconociendo la sabiduría que Dios comunica a través de su experiencia de precariedad, confianza y ayuda mutua.
\end{itemize}

\subsubsection{Ideas clave para retener (síntesis rápida)}
\begin{enumerate}
\item La DSI tiene raíz popular: nace del contacto con los conflictos sociales reales.
\item Línea histórica de encíclicas: \textit{Rerum novarum $\rightarrow$ Mater et Magistra $\rightarrow$ Gaudium et spes $\rightarrow$ Populorum progressio $\rightarrow$ Sollicitudo rei socialis / Laborem exercens $\rightarrow$ Caritas in veritate $\rightarrow$ Dilexit nos $\rightarrow$ Dilexi te}.
\item El Vaticano II consolida la centralidad de los pobres, no como tema social aislado sino \textbf{cristológico}.
\item Aporte latinoamericano decisivo: Medellín, Puebla, Aparecida; opción preferencial por los pobres = ``pecado social'' como categoría.
\item Paso conceptual clave: de la \textbf{asistencia} (paternalismo) al reconocimiento de los pobres como \textbf{sujetos} activos, portadores de sabiduría y evangelizadores.
\item La pobreza tiene causas \textbf{estructurales} que deben transformarse, no solo paliarse.
\end{enumerate}



\subsection{1.2 Sorge --- Cap. I, V, VI}

\subsubsection{CAP. I --- Del evangelio al ``discurso social''}

\noindent \textbf{Idea central:} la Iglesia tiene derecho-deber de intervenir en lo social porque esa intervención \textbf{forma parte de su misión religiosa}, no es un añadido político.\par\smallskip

\medskip\noindent \textbf{1. La misión propia de la Iglesia es religiosa}\par\nopagebreak
\begin{itemize}
\item La misión de la Iglesia es esencialmente anunciar el Evangelio (salvación y redención), no es de orden político, económico o social. Su origen es sobrenatural y su fin último también (la salvación eterna).
\item Sin embargo, esta naturaleza sobrenatural \textbf{no la aparta de la historia}: ``el reino de Dios se realiza en la historia''. La redención de Cristo también restaura el orden temporal, por eso la Iglesia busca impregnar todo el orden temporal con el espíritu evangélico.
\item El mensaje evangélico vale para todos los hombres de buena voluntad, no solo para los cristianos.
\end{itemize}

\medskip\noindent \textbf{2. El evangelio libera al hombre}\par\nopagebreak
\begin{itemize}
\item Hay un nexo inseparable entre anunciar el Evangelio y la liberación del hombre: no se puede evangelizar sin que el hombre progrese también civilmente.
\item Es una liberación \textbf{integral}: ante todo del pecado, pero también de las múltiples esclavitudes culturales, económicas, sociales y políticas.
\item Problema clave: la relación entre \textbf{evangelio y cultura}. Si el mensaje no se traduce en valores, costumbres y lenguaje de una cultura concreta, queda incomprensible. Pablo VI: ``el drama de nuestra época es la fractura entre evangelio y cultura''. De ahí la necesidad de una nueva \textbf{inculturación} de la fe.
\end{itemize}

\medskip\noindent \textbf{3. Legitimidad del discurso social de la Iglesia}\par\nopagebreak
\noindent Dos motivos fundamentan que la Iglesia pueda intervenir en lo social:\par\smallskip
\begin{enumerate}
\item \textbf{El orden moral está conectado con el orden sobrenatural}: las cosas terrenas también se ordenan a la salvación del hombre. La Iglesia interviene en cuestiones sociales, económicas y políticas \textit{en la medida en que afectan al campo moral} --- no tiene competencia técnica, sino que se presenta como ``experta en humanidad''.
\item \textbf{La Revelación tiene una dimensión histórica}: Cristo entra en la historia y la asume; la Encarnación se cumple progresivamente a través de las épocas, por eso la Iglesia debe leer e iluminar cada época con la luz del Evangelio, ofreciendo criterios de juicio, principios de reflexión y orientaciones de acción --- un ``discurso social'' abierto y siempre vigente, iniciado hace más de cien años con \textit{Rerum novarum} (1891).
\end{enumerate}

\medskip\noindent \textbf{4. ``Doctrina social''}\par\nopagebreak
\begin{itemize}
\item El término indica un contenido orgánico y sistemático: un conjunto de principios deducidos teóricamente y orientados a la acción social y política.
\item Corresponde al uso típico entre León XIII y Pío XII: una \textbf{``ideología católica''} o ``tercera vía'' frente al socialismo y al capitalismo.
\end{itemize}

\medskip\noindent \textbf{5. ``Enseñanza social''}\par\nopagebreak
\begin{itemize}
\item Pone el acento en el aspecto histórico y práctico (método inductivo: partir de los ``signos de los tiempos'' y leerlos a la luz de los principios), no tanto en la afirmación teórica deductiva.
\item El Concilio Vaticano II evitó conscientemente la expresión ``doctrina social'' en sus documentos (salvo una excepción por malentendido en \textit{Gaudium et spes} n. 76), prefiriendo ``enseñanza social'', para subrayar el carácter \textbf{dinámico} de las intervenciones del magisterio, con elementos permanentes y otros contingentes.
\item Puebla (1979) también prefirió hablar de ``enseñanzas sociales'' en plural, aunque el texto final en Roma volvió a ``Enseñanza social'' en singular.
\item Juan Pablo II clarificó que no es una ``tercera vía'' ni una ideología, sino la reflexión atenta sobre la vida del hombre a la luz de la fe y de la tradición eclesial; pertenece al ámbito de la \textbf{teología (moral)}, no de la ideología.
\end{itemize}

\medskip\noindent \textbf{6. El ``discurso social'' de la Iglesia}\par\nopagebreak
\begin{itemize}
\item Hoy conviene hablar de \textbf{``discurso social''} de la Iglesia: integra doctrina social y enseñanza social como etapas de un mismo proceso, formado por múltiples intervenciones magisteriales (grandes encíclicas y pronunciamientos más locales y contingentes).
\item Se distingue lo \textbf{permanente} (grandes principios morales, conformes al Evangelio y a la recta razón) de lo \textbf{contingente/histórico} (juicios sobre realidades sociales, políticas y económicas cambiantes).
\item Sorge propone leer este discurso social por \textbf{fases históricas}: la ideología católica (1891-1931), la ``nueva cristiandad'' (1931-1958), el diálogo (1958-1978) y la profecía (1978-2005) --- esta última es justamente el tema del Capítulo V.
\item Recuadro del \textit{Compendio} (nn. 72-86): la doctrina social es un \textit{corpus} doctrinal único y orgánico, parte integrante de la teología moral católica; no solo ofrece valores y principios, sino también directrices de acción; toda la comunidad eclesial participa en su elaboración, pero el magisterio la promulga y unifica.
\end{itemize}



\subsubsection{CAP. V --- La profecía (1978-2005)}

\noindent \textbf{Contexto:} pontificado de Juan Pablo II (1978-2005), con el breve paréntesis de Juan Pablo I. La Iglesia entra en una etapa \textbf{profética}: confirma lo adquirido por el discurso social anterior pero lo relanza hacia el tercer milenio.\par\smallskip

\begin{itemize}
\item La cuestión social ya no es solo cuantitativa (distribución material) sino \textbf{cualitativa}: toca la vida humana y sus valores fundamentales (vida, familia, trabajo). Es cuestión de ``calidad de vida'' y de ``muerte del hombre y de su hábitat''.
\item Frente a esto, la Iglesia no analiza científicamente los cambios sociales (no es su tarea), sino que orienta esos cambios para que sea realidad un auténtico progreso del hombre.
\item El discurso social se vuelve \textbf{universal (``ecuménico'')}: puede ser comprendido y compartido por creyentes y no creyentes, del Este y del Oeste.
\end{itemize}

\noindent El capítulo desarrolla tres ``evangelios'' como ejes del anuncio profético de Juan Pablo II:\par\smallskip

\medskip\noindent \textbf{1. El evangelio de la vida}\par\nopagebreak
\begin{itemize}
\item Diagnóstico: la cultura actual da a la vida humana un valor solo \textbf{relativo} (puede ceder ante el bienestar, el interés económico, la investigación científica), en lugar de un valor \textbf{absoluto}.
\item Consecuencia: pérdida del sentido de Dios $\rightarrow$ pérdida del sentido del hombre, su dignidad y su vida (círculo vicioso). Se instala un materialismo práctico: ``los valores del ser son sustituidos por los del tener''.
\item Paradoja: las declaraciones universales de derechos humanos afirman la dignidad de toda persona, pero en la realidad se la niega con frecuencia --- ``sociedades de excluidos'' en vez de ``sociedades de convivientes''.
\item El derecho a la vida es el \textbf{fundamento de todos los demás derechos}: sin él, de nada sirven la libertad de conciencia, palabra o asociación.
\item La \textbf{familia} aparece como ``santuario de la vida'': ámbito privilegiado donde la vida (don de Dios) puede ser acogida y protegida.
\item Recuadro \textit{Compendio} n. 231: la familia fundada en el matrimonio es el santuario de la vida.
\end{itemize}

\medskip\noindent \textbf{2. El evangelio del trabajo}\par\nopagebreak
\begin{itemize}
\item El trabajo también sufre una crisis de sentido: muchos trabajadores lo ven solo como fatalidad o medio de subsistencia, no como bien para el crecimiento personal y social.
\item Ni el capitalismo ni el socialismo son la solución. Clave (\textit{Laborem exercens}): el fundamento del valor del trabajo humano \textbf{no es el tipo de trabajo, sino el hecho de que quien lo realiza es una persona} $\rightarrow$ dimensión subjetiva por encima de la objetiva: ``el trabajo está en función del hombre, no el hombre en función del trabajo''.
\item Tres aspectos del verdadero sentido del trabajo:
\begin{enumerate}
\item El trabajo es siempre acto de una \textbf{persona}.
\item El trabajo es \textbf{acto creativo}: el capital es solo instrumento, nace del trabajo humano (prioridad del trabajo sobre el capital).
\item El trabajo es intrínsecamente un \textbf{acto de solidaridad} --- una ``solidaridad para'' (positiva), no ``en contra de'' alguien.
\end{enumerate}
\item A través del trabajo el hombre participa en la obra de la creación y colabora con Cristo en su obra de liberación y redención.
\end{itemize}

\medskip\noindent \textbf{3. El evangelio de la caridad}\par\nopagebreak
\begin{itemize}
\item La caridad cristiana \textbf{excede} la mera filantropía y la regla de oro: es el mandamiento nuevo (``amaos los unos a los otros como yo os he amado''), con el mismo amor sobrenatural de Dios.
\item La caridad es central para la credibilidad del anuncio cristiano (``la caridad está en el centro del evangelio'') y para la construcción de una sociedad a la medida del hombre.
\item Juan Pablo II sitúa en primer lugar de su doctrina social la \textbf{opción u opción preferencial por los pobres}, abarcando a hambrientos, mendigos, sin techo, sin asistencia sanitaria y sin esperanza.
\item La opción por los pobres conecta con el tema de la \textbf{liberación}: de todas las formas de esclavitud, de las ``nuevas pobrezas'' y de las \textbf{``estructuras de pecado''} (situaciones estructurales de injusticia).
\item Relación justicia-caridad: la justicia es ``la medida mínima de la caridad'' (Pablo VI), pero la justicia sola no basta --- puede llevar a la negación de la persona si no se completa con el amor.
\item \textit{Deus caritas est} (Benedicto XVI, 2005) confirma esta síntesis: afirmar que las estructuras justas harían superfluas las obras de caridad esconde una concepción materialista del hombre (``no solo de pan vive el hombre'').
\end{itemize}

\medskip\noindent \textbf{4. Un ``discurso'' ecuménico}\par\nopagebreak
\begin{itemize}
\item El discurso social de la Iglesia alcanza su madurez: ya no es reflexión filosófico-teológica sobre ideologías ni ``tercera vía'', sino reflexión en contacto permanente con las situaciones cambiantes del mundo, bajo el impulso del Evangelio (Pablo VI).
\item Es \textbf{ecuménico} en dos sentidos: (a) de valor universal, válido para todos y todas las situaciones; (b) en sentido técnico, implicando el diálogo con otras Iglesias cristianas y confesiones religiosas (ej. encuentro de Asís, 1986, sobre la paz).
\item Causas de la injusticia = causas culturales, ligadas a una determinada visión del hombre, la sociedad y el mundo; en el centro de la cuestión cultural está el sentido moral, que se funda en el sentido religioso.
\end{itemize}



\subsubsection{CAP. VI --- El giro de Benedicto XVI: \textit{Caritas in veritate} (2009)}

\noindent \textbf{Contexto histórico:} tras la caída del muro de Berlín (1989, fin del ``socialismo real'') y la crisis financiera de 2008 (crisis del ``capitalismo liberal''), también entra en crisis el modelo intermedio de la ``nueva cristiandad'' (tercera vía de inspiración cristiana, elaborada sobre todo por Jacques Maritain). El vacío ideológico es ocupado por una nueva \textbf{``ideología tecnocrática''}, que se convierte en el ``pensamiento único'' dominante.\par\smallskip

\begin{itemize}
\item \textit{Caritas in veritate} representa un giro en la historia de las encíclicas sociales: Benedicto XVI la presenta como una especie de brújula/\textit{carta magna} para el siglo XXI, tomando como referencia no tanto \textit{Rerum novarum} sino la \textit{Populorum progressio} de Pablo VI, a la que llama ``la \textit{Rerum novarum} de la época contemporánea''.
\item La cuestión social ya no es la lucha de clases ni el enfrentamiento este-oeste, sino que se ha convertido en \textbf{``cuestión antropológica''}: el desafío central es cómo se concibe la vida humana, manipulable hoy mediante biotecnologías (fecundación in vitro, embriones, clonación, hibridación).
\end{itemize}

\medskip\noindent \textbf{1. Crítica de la ``ideología tecnocrática'' dominante}\par\nopagebreak
\begin{itemize}
\item El fin del socialismo real no significó la victoria del capitalismo: persisten contradicciones del sistema capitalista (el hombre privado de dignidad, esclavo de necesidades ficticias).
\item La cultura tecnocrática ve a la sociedad como un agregado de individuos anónimos, cada uno pendiente solo de sí; el juicio ético queda subordinado a la eficiencia y a la innovación técnica.
\item Vicio de fondo: \textbf{materialismo utilitarista} --- solo vale lo ``eficaz'' y lo que rinde en términos económicos (``política del hacer'').
\item Frente a esto, Benedicto XVI propone una concepción humana integral del desarrollo: ``sin Dios el hombre no sabe adónde ir ni logra entender quién es''.
\end{itemize}

\medskip\noindent \textbf{2. Principios éticos, culturales y políticos de un nuevo humanismo}\par\nopagebreak
\noindent Se articulan en tres planos:\par\smallskip

\medskip\noindent \textbf{a) Plano ético --- ``libertad responsable''}\par\nopagebreak
\begin{itemize}
\item Parte del concepto de progreso como \textbf{vocación} (ya en Pablo VI): la vida es un ``don'' recibido, no autoproducido; toda persona está ``llamada a la vida'' y debe realizar ese proyecto libre y responsablemente.
\item Clave de la encíclica: la relación entre \textbf{verdad y caridad} (``caridad en la verdad''): el hombre está hecho para el amor y para la verdad, ambos son dones recibidos, no producidos por nosotros.
\item El ``desarrollo humano integral'' supone la \textbf{libertad responsable}: ninguna estructura externa puede garantizar el desarrollo por encima de la responsabilidad humana; solo en libertad responsable el desarrollo puede ser integralmente humano.
\end{itemize}

\medskip\noindent \textbf{b) Plano socio-cultural --- ``fraternidad''}\par\nopagebreak
\begin{itemize}
\item La globalización no es solo un fenómeno económico-financiero, sino también social y cultural: ``el mercado se ha transformado en el instrumento de una nueva cultura'' que impone su escala de valores.
\item La globalización \textbf{no es en sí buena ni mala}: ``será lo que la gente haga de ella'' --- puede generar nuevas esclavitudes o bien nuevas oportunidades de desarrollo, paz y derechos humanos, según se la oriente.
\item Necesidad de una \textbf{globalización solidaria}, con alma ética, que haga crecer la fraternidad. La razón sola puede establecer igualdad y convivencia cívica, pero no puede fundar la hermandad: esta nace de una vocación trascendente de Dios Padre.
\end{itemize}

\medskip\noindent \textbf{c) Plano político --- ``reciprocidad''}\par\nopagebreak
\begin{itemize}
\item La crisis del ``pensamiento único'' tecnocrático debilita la \textbf{democracia representativa} al corroer sus pilares: dignidad de la persona, solidaridad y subsidiariedad responsable.
\item Los valores democráticos no los crea el Estado: preexisten a la organización política y son punto de referencia normativo para la ley civil.
\item Riesgo: una democracia privada de su alma ética puede volverse ``democracia totalitaria'' (totalitarismo enmascarado).
\item Principio de \textbf{reciprocidad} en dos sentidos:
\begin{enumerate}
\item Las opciones políticas y las reformas deben abordarse con perspectiva \textbf{interdisciplinar} (valoración moral + investigación científica juntas); los derechos presuponen deberes. Dentro de este sentido, la estrecha conexión entre \textbf{ética personal y ética social}: el desarrollo es imposible sin ``hombres rectos'', con preparación profesional y coherencia moral.
\item Un progreso integralmente humano no puede prescindir de la \textbf{conciencia religiosa}: la democracia es un método, no es autosuficiente (Bobbio, Habermas, Croce); de ahí la ``laicidad positiva''.
\end{enumerate}
\item Reciprocidad también entre \textbf{razón y fe/religión}: la razón necesita ser purificada por la fe, y la fe purificada por la razón, para mostrar su rostro auténticamente humano. La Iglesia no ofrece soluciones técnicas ni se mezcla en política, pero tiene una misión de verdad al servicio de una sociedad a la medida del hombre.
\end{itemize}

\subsubsection{Ideas clave para retener (síntesis rápida --- Sorge Cap. I, V, VI)}
\begin{enumerate}
\item \textbf{Legitimidad} del discurso social: se apoya en la conexión orden moral--orden sobrenatural, y en la dimensión histórica de la Revelación.
\item Evolución terminológica: \textbf{doctrina social} (más estática/deductiva) $\rightarrow$ \textbf{enseñanza social} (más dinámica/inductiva) $\rightarrow$ hoy se prefiere \textbf{``discurso social''} (integra ambas, con elementos permanentes y contingentes).
\item Fases históricas del discurso social: ideología católica (1891-1931) $\rightarrow$ nueva cristiandad (1931-1958) $\rightarrow$ diálogo (1958-1978) $\rightarrow$ \textbf{profecía} (1978-2005, Juan Pablo II) $\rightarrow$ giro de Benedicto XVI (2009).
\item Juan Pablo II: tres ``evangelios'' --- \textbf{vida, trabajo y caridad} --- como respuesta profética a la crisis cualitativa de la ``cuestión social''.
\item Benedicto XVI (\textit{Caritas in veritate}): pasaje de la cuestión social clásica a la \textbf{cuestión antropológica}; crítica a la ideología tecnocrática; tres principios de un nuevo humanismo --- \textbf{libertad responsable (ético) / fraternidad (socio-cultural) / reciprocidad (político)}.
\item Fórmula clave de \textit{Caritas in veritate}: \textbf{``caridad en la verdad''} como motor del desarrollo humano integral.
\end{enumerate}



\subsection{1.3 Compendio DSI --- Introducción y Cap. II}

\subsubsection{INTRODUCCIÓN --- ``Un humanismo integral y solidario''}

\medskip\noindent \textbf{a) Al alba del tercer milenio (\S{}1--6)}\par\nopagebreak
\begin{itemize}
\item La Iglesia entra al tercer milenio guiada por Cristo, camino, verdad y vida. La salvación que Él trae es \textbf{integral}: alcanza también la economía, el trabajo, la técnica, la comunicación, la sociedad, la política y las relaciones internacionales.
\item La Iglesia ofrece su doctrina social como parte de su misión de anunciar el Evangelio: al anunciarlo, ``enseña al hombre... su dignidad propia y su vocación a la comunión de las personas'' y le descubre las exigencias de justicia y paz.
\item Esta doctrina tiene una unidad profunda que brota de la \textbf{Fe} (salvación integral), la \textbf{Esperanza} (justicia plena) y la \textbf{Caridad} (fraternidad universal en Cristo).
\item Descubrirse amado por Dios lleva al hombre a comprender su dignidad trascendente y a salir al encuentro del otro; solo el amor puede transformar radicalmente las relaciones humanas.
\item Panorama de pobrezas antiguas y nuevas (hambre, analfabetismo, falta de vivienda, drogadicción, marginación, desequilibrio ecológico, amenazas a la paz, violación de derechos humanos) que interpelan la conciencia y exigen compromiso.
\item La humanidad percibe que está ligada por un \textbf{destino único} (a menudo impuesto por la técnica o la economía) y necesita una mayor conciencia moral: de ahí la propuesta de un \textbf{humanismo integral y solidario}.
\end{itemize}

\medskip\noindent \textbf{b) El significado del documento (\S{}7--12)}\par\nopagebreak
\begin{itemize}
\item El Compendio ofrece principios de reflexión, criterios de juicio y directrices de acción; su difusión es una \textbf{prioridad pastoral}, porque enseñar y difundir la doctrina social forma parte de la misión evangelizadora de la Iglesia.
\item Elaborado por el Pontificio Consejo ``Justicia y Paz'' mediante amplia consulta (Dicasterios, Conferencias Episcopales, obispos y expertos).
\item Presenta de manera \textbf{orgánica y sistemática} (aunque sintética) los aspectos teológicos, filosóficos, morales, culturales y pastorales de la enseñanza social.
\item Importante: las citas del documento provienen de textos de \textbf{diversa autoridad} (documentos conciliares, encíclicas, discursos pontificios, documentos de Dicasterios) --- el lector debe ser consciente de que hay distintos grados de enseñanza.
\item El Compendio deja a las \textbf{Conferencias Episcopales} la aplicación concreta a las situaciones locales.
\item Destinatarios: en primer lugar los \textbf{Obispos}; también sacerdotes, religiosos, formadores, fieles laicos y comunidades cristianas; y además hermanos de otras Iglesias, otras religiones y todo hombre de buena voluntad comprometido con el bien común.
\end{itemize}

\medskip\noindent \textbf{c) Al servicio de la verdad plena del hombre (\S{}13--17)}\par\nopagebreak
\begin{itemize}
\item El documento coloca al \textbf{hombre entero} (cuerpo y alma, corazón y conciencia, inteligencia y voluntad) como eje de toda la exposición, siguiendo el estilo de \textit{Gaudium et spes}.
\item Retoma las preguntas radicales de la existencia humana (quién soy, sentido del dolor, la muerte, etc. --- recuerda el ``conócete a ti mismo'' de Delfos) para mostrar que la búsqueda de la verdad es constitutiva del ser humano.
\item Se identifican \textbf{tres grandes desafíos} actuales:
\begin{enumerate}
\item La \textbf{verdad del ser-hombre}: el límite entre naturaleza, técnica y moral (qué es el hombre, qué puede hacer, qué debe ser).
\item La gestión del \textbf{pluralismo} (de pensamiento, moral, cultural, religioso).
\item La \textbf{globalización}, que abre una nueva época en el destino de la humanidad.
\end{enumerate}
\end{itemize}

\medskip\noindent \textbf{d) Bajo el signo de la solidaridad, del respeto y del amor (\S{}18--19)}\par\nopagebreak
\begin{itemize}
\item La Iglesia camina junto a toda la humanidad, sin ser del mundo pero al servicio del mundo (cf. \textit{Gaudium et spes} 3).
\item El documento propone un \textbf{humanismo integral y solidario}, capaz de animar un nuevo orden social, económico y político, fundado en la dignidad y libertad de la persona humana, y realizado en la paz, la justicia y la solidaridad.
\end{itemize}



\subsubsection{CAPÍTULO SEGUNDO --- ``Misión de la Iglesia y doctrina social''}

\medskip\noindent \textbf{I. EVANGELIZACIÓN Y DOCTRINA SOCIAL}\par\nopagebreak

\medskip\noindent \textbf{a) La Iglesia, morada de Dios con los hombres (\S{}60--61)}\par\nopagebreak
\begin{itemize}
\item La Iglesia es solidaria de todo hombre y mujer, ``sacramento del amor de Dios'' y de esperanza; sirve a la salvación no en abstracto, sino en la historia concreta donde el hombre vive.
\item El hombre es un ser esencialmente relacional (vive en sociedad, familias, grupos intermedios); de ahí nace la sociedad con sus estructuras políticas, económicas, jurídicas y culturales, a la que la Iglesia se dirige con su doctrina social.
\end{itemize}

\medskip\noindent \textbf{b) Fecundar y fermentar la sociedad con el Evangelio (\S{}62--65)}\par\nopagebreak
\begin{itemize}
\item La Iglesia no solo se dirige al hombre \textit{en} la sociedad, sino que busca \textbf{fecundar la sociedad misma} con el Evangelio.
\item La sociedad, la política, la economía, el trabajo, el derecho y la cultura \textbf{no son ámbitos ajenos} a la salvación: atañen al hombre, que es ``el camino primero y fundamental de la Iglesia'' (Juan Pablo II).
\item La redención es de orden sobrenatural, pero lo sobrenatural \textbf{no excluye ni desprecia lo natural/humano}: lo asume y lo eleva (cf. Encarnación).
\item La Redención comienza con la \textbf{Encarnación}: el Hijo de Dios asume todo lo humano (menos el pecado). Por eso la Iglesia, con su doctrina social, es fiel a Cristo y se revela como ``sacramento universal de salvación''.
\end{itemize}

\medskip\noindent \textbf{c) Doctrina social, evangelización y promoción humana (\S{}66--68)}\par\nopagebreak
\begin{itemize}
\item La doctrina social es \textbf{parte integrante del ministerio de evangelización}. Existen vínculos profundos entre evangelización y promoción humana:
\begin{enumerate}
\item \textbf{Antropológicos}: el hombre a evangelizar está sujeto a problemas sociales y económicos concretos.
\item \textbf{Teológicos}: no se puede separar el plan de la creación del plan de la redención (que alcanza también las situaciones de injusticia).
\item \textbf{Evangélicos}: la caridad no puede predicarse sin promover, mediante justicia y paz, el crecimiento real del hombre.
\end{enumerate}
\item La doctrina social ``tiene de por sí el valor de un instrumento de evangelización''; es un camino propio para el ejercicio de la función profética de la Iglesia.
\item Aun así, la Iglesia \textbf{no entra en cuestiones técnicas} ni propone sistemas o modelos concretos de organización social: eso no corresponde a su misión, que es de orden religioso (recuerda \textit{Sorge Cap. I}: la Iglesia como ``experta en humanidad'').
\end{itemize}

\medskip\noindent \textbf{d) Derecho y deber de la Iglesia (\S{}69--71)}\par\nopagebreak
\begin{itemize}
\item Con su doctrina social la Iglesia ``se propone \textbf{ayudar al hombre en el camino de la salvación}'': es su \textbf{fin primordial y único}; no busca invadir competencias ajenas.
\item La Iglesia tiene el \textbf{derecho} de ser maestra de la verdad de fe y de la verdad moral, y esto configura también un \textbf{deber}: no puede renunciar a evangelizar lo social sin negarse a sí misma (``¡Ay de mí si no predicara el Evangelio!'' --- 1 Co 9,16).
\item No hay que reducir el hecho religioso a la esfera privada, ni tampoco reducir la salvación a algo puramente ultraterreno.
\item Es tarea de la Iglesia ``anunciar siempre y en todas partes los principios morales acerca del orden social'' y pronunciarse sobre cualquier realidad humana cuando lo exigen los derechos fundamentales de la persona o la salvación de las almas.
\end{itemize}

\medskip\noindent \textbf{II. LA NATURALEZA DE LA DOCTRINA SOCIAL}\par\nopagebreak

\medskip\noindent \textbf{a) Un conocimiento iluminado por la fe (\S{}72--75)}\par\nopagebreak
\begin{itemize}
\item La doctrina social es \textbf{de naturaleza teológica}, específicamente \textbf{teológico-moral}: ``una doctrina que debe orientar la conducta de las personas''.
\item La doctrina social \textbf{no fue pensada desde el inicio como sistema orgánico}, sino que se formó progresivamente a través de numerosas intervenciones magisteriales --- de ahí las oscilaciones históricas sobre su naturaleza y método (tema ya visto en Sorge Cap. I).
\item \textit{Sollicitudo rei socialis} (Juan Pablo II) clarifica: la doctrina social \textbf{no pertenece a la ideología, sino a la teología (moral)}. No es un sistema socioeconómico, sino la formulación cuidadosa del resultado de una reflexión sobre la vida del hombre en sociedad, a la luz de la fe y de la tradición eclesial.
\item Refleja \textbf{tres niveles de la enseñanza teológico-moral}:
\begin{enumerate}
\item \textbf{Fundante} --- de las motivaciones.
\item \textbf{Directivo} --- de las normas de la vida social.
\item \textbf{Deliberativo} --- de la conciencia, que media las normas generales en cada situación concreta.
\end{enumerate}
\item Fundamento esencial: la \textbf{Revelación bíblica y la Tradición}. Dos fuentes cognoscitivas: la \textbf{Revelación} y la \textbf{naturaleza humana} (razón); dos vías que se articulan: \textbf{fe y razón} (\textit{fides et ratio}).
\item Al ser un conocimiento iluminado por la fe pero racionalmente fundado, tiene \textbf{plausibilidad universal}: puede ser comprendida y compartida por todos.
\end{itemize}

\medskip\noindent \textbf{b) En diálogo cordial con todos los saberes (\S{}76--78)}\par\nopagebreak
\begin{itemize}
\item Tiene una fuerte \textbf{dimensión interdisciplinar}: dialoga con la filosofía (que ofrece los conceptos básicos: persona, sociedad, libertad, justicia, bien común, solidaridad, subsidiariedad, Estado, etc.) y con las ciencias humanas y sociales (que aportan conocimiento concreto, actual y descriptivo de la realidad social).
\item Este diálogo enriquece a la doctrina social con \textbf{competencia, concreción y actualidad}, y a la vez invita a las ciencias a abrirse a horizontes de sentido y valor.
\end{itemize}

\medskip\noindent \textbf{c) Expresión del ministerio de enseñanza de la Iglesia (\S{}79--80)}\par\nopagebreak
\begin{itemize}
\item La doctrina social \textbf{es de la Iglesia} porque toda la comunidad eclesial (sacerdotes, religiosos, laicos) participa en su elaboración, según carismas y tareas --- pero es el \textbf{Magisterio} (Papa y Obispos en comunión con él, con el ``\textit{munus docendi}'') quien la asume, interpreta y unifica.
\item Se articulan el Magisterio universal (Papa/Concilio) y el Magisterio episcopal (que aplica y actualiza localmente), en una \textbf{circularidad} que expresa la colegialidad.
\item Al ser parte de la enseñanza moral de la Iglesia, tiene la misma dignidad y autoridad: es \textbf{Magisterio auténtico} que exige aceptación y adhesión de los fieles (el grado de asentimiento depende de la naturaleza y autoridad de cada enseñanza concreta).
\end{itemize}

\medskip\noindent \textbf{d) Hacia una sociedad reconciliada en la justicia y en el amor (\S{}81--82)}\par\nopagebreak
\begin{itemize}
\item Objeto de la doctrina social: \textbf{el hombre}, llamado a la salvación y confiado por Cristo a la Iglesia.
\item La doctrina social realiza una doble tarea:
\begin{enumerate}
\item \textbf{Anuncio}: ofrece ``una visión global del hombre y de la humanidad'', con valores, criterios de juicio y también directrices de acción --- pero no busca estructurar u organizar la sociedad, sino formar las conciencias.
\item \textbf{Denuncia}: del pecado de injusticia y violencia, defendiendo especialmente los derechos de los pobres, los pequeños y los débiles.
\end{enumerate}
\item Finalidad \textbf{religiosa y moral}: religiosa porque la misión salvífica alcanza al hombre en toda su existencia personal y social; moral porque busca un ``humanismo pleno'', la liberación de todo lo que oprime al hombre y el desarrollo integral de todo el hombre y de todos los hombres.
\end{itemize}

\medskip\noindent \textbf{e) Un mensaje para los hijos de la Iglesia y para la humanidad (\S{}83--84)}\par\nopagebreak
\begin{itemize}
\item Primer destinatario: la \textbf{comunidad eclesial} entera. A los \textbf{laicos} corresponden de modo peculiar las responsabilidades seculares (políticas, económicas, administrativas) que derivan de la doctrina social, por la índole secular de su vocación.
\item Destinatario universal: \textbf{todos los hombres de buena voluntad}, no solo católicos ni cristianos.
\end{itemize}

\medskip\noindent \textbf{f) Bajo el signo de la continuidad y de la renovación (\S{}85--86)}\par\nopagebreak
\begin{itemize}
\item La doctrina social se caracteriza por dos rasgos simultáneos:
\begin{itemize}
\item \textbf{Continuidad}: fundada en valores universales (Revelación + naturaleza humana), no depende de culturas, ideologías u opiniones; se mantiene idéntica en su inspiración de fondo y sus principios fundamentales.
\item \textbf{Renovación}: constante atención a la historia y capacidad de abrirse a lo nuevo, sin diluirse en ello --- se describe como un \textbf{``taller'' siempre abierto}.
\end{itemize}
\end{itemize}

\medskip\noindent \textbf{III. LA DOCTRINA SOCIAL EN NUESTRO TIEMPO: APUNTES HISTÓRICOS}\par\nopagebreak

\medskip\noindent \textbf{a) El comienzo de un nuevo camino (\S{}87--88)}\par\nopagebreak
\begin{itemize}
\item El término ``doctrina social'' se remonta a \textbf{Pío XI}, pero designa el \textit{corpus} doctrinal que arranca formalmente con \textit{Rerum novarum} (León XIII, 1891) --- aunque la preocupación de la Iglesia por lo social no comenzó ahí, sino que tiene raíces en la Escritura, los Padres de la Iglesia y los grandes Doctores medievales.
\item La \textbf{revolución industrial} del s. XIX genera la primera gran ``cuestión social'': la \textbf{cuestión obrera}, por el conflicto entre capital y trabajo.
\end{itemize}

\noindent \textbf{b) De la \textit{Rerum novarum} hasta nuestros días (\S{}89--103) --- Recorrido histórico del Magisterio social}\par\smallskip

{\small\renewcommand{\arraystretch}{1.25}
\begin{longtable}{>{\raggedright\arraybackslash}p{0.439\linewidth}>{\raggedright\arraybackslash}p{0.032\linewidth}>{\raggedright\arraybackslash}p{0.439\linewidth}}
\toprule
\textbf{Papa / Documento} & \textbf{Año} & \textbf{Aporte clave} \\ \midrule \endhead
León XIII, \textit{Rerum novarum} & 1891 & Primera encíclica social; condición de los obreros; excluye el socialismo; defiende propiedad privada con función social, colaboración (no lucha de clases), derecho de asociación. ``Carta de ciudadanía'' de la Iglesia en la vida pública. \\ \addlinespace[2pt]
Pío XI, \textit{Quadragesimo anno} & 1931 & 40\textordmasculine{} aniversario; principio de \textbf{subsidiariedad}; salario familiar; rechaza el liberalismo ilimitado; reafirma función social de la propiedad. \\ \addlinespace[2pt]
Pío XI, \textit{Non abbiamo bisogno} & 1931 & Protesta contra el fascismo en Italia. \\ \addlinespace[2pt]
Pío XI, \textit{Mit brennender Sorge} & 1937 & Sobre la Iglesia católica en la Alemania nazi; llamado a la resistencia. \\ \addlinespace[2pt]
Pío XI, \textit{Divini Redemptoris} & 1937 & Crítica sistemática al comunismo ateo, ``intrínsecamente malo''. \\ \addlinespace[2pt]
Pío XII, Radiomensajes navideños & años 40 & Sin encíclica social propia, pero voz de la ``conciencia universal'' durante la guerra y posguerra; precursor del Vaticano II. \\ \addlinespace[2pt]
Juan XXIII, \textit{Mater et magistra} & 1961 & Palabras clave: comunidad y socialización; universaliza la cuestión social (Tercer Mundo). \\ \addlinespace[2pt]
Juan XXIII, \textit{Pacem in terris} & 1963 & Primera reflexión a fondo sobre \textbf{derechos humanos}; primer documento dirigido también ``a todos los hombres de buena voluntad''. \\ \addlinespace[2pt]
Concilio Vaticano II, \textit{Gaudium et spes} & 1965 & Constitución pastoral; Iglesia ``íntima y realmente solidaria del género humano''; estudia cultura, economía, familia, comunidad política, paz. \\ \addlinespace[2pt]
Concilio Vaticano II, \textit{Dignitatis humanae} & 1965 & Derecho a la \textbf{libertad religiosa}, fundado en la dignidad humana. \\ \addlinespace[2pt]
Pablo VI, \textit{Populorum progressio} & 1967 & ``El desarrollo es el nuevo nombre de la paz''; desarrollo integral y solidario. \\ \addlinespace[2pt]
Pablo VI, Pontificia Comisión \textit{Iustitia et Pax} / Jornada Mundial de la Paz & 1967--68 & Institucionaliza el compromiso social de la Iglesia. \\ \addlinespace[2pt]
Pablo VI, \textit{Octogesima adveniens} & 1971 & 80\textordmasculine{} aniversario de \textit{Rerum novarum}; insuficiencia de las ideologías; nuevos problemas (urbanización, juventud, mujer, ecología, medios). \\ \addlinespace[2pt]
Juan Pablo II, \textit{Laborem exercens} & 1981 & 90\textordmasculine{} aniversario; el trabajo como clave de toda la cuestión social; dimensión subjetiva del trabajo. \\ \addlinespace[2pt]
Juan Pablo II, \textit{Sollicitudo rei socialis} & 1987 & 20\textordmasculine{} aniversario de \textit{Populorum progressio}; distingue progreso y desarrollo; ``Opus solidaritatis pax''. \\ \addlinespace[2pt]
Juan Pablo II, \textit{Centesimus annus} & 1991 & Centenario de \textit{Rerum novarum}; principio de \textbf{solidaridad}; balance tras la caída del comunismo (1989). \\ \addlinespace[2pt]
\bottomrule\end{longtable}}

\medskip\noindent \textbf{c) A la luz y bajo el impulso del Evangelio (\S{}104)}\par\nopagebreak
\begin{itemize}
\item Este recorrido histórico muestra un \textit{corpus} doctrinal \textbf{renovado y en desarrollo constante}, no un sistema cerrado: la Iglesia, a la luz de la Palabra revelada y con la asistencia del Espíritu Santo, va leyendo los signos de los tiempos a medida que la historia avanza.
\end{itemize}

\subsubsection{Ideas clave para retener (síntesis rápida --- Compendio, Introducción y Cap. II)}
\begin{enumerate}
\item Propuesta central: un \textbf{humanismo integral y solidario}, fundado en la dignidad de la persona, realizado en paz, justicia y solidaridad.
\item El Compendio distingue \textbf{grados de autoridad} en sus citas (documentos conciliares y encíclicas $\neq$ discursos o documentos de Dicasterios).
\item La doctrina social \textbf{pertenece a la teología moral, no a la ideología} (clave: \textit{Sollicitudo rei socialis}), con tres niveles: fundante / directivo / deliberativo.
\item Dos fuentes (Revelación y naturaleza humana) y dos vías cognoscitivas (\textbf{fe y razón}) --- diálogo interdisciplinar con filosofía y ciencias humanas.
\item La doctrina social es del \textbf{Magisterio} (Papa + Obispos), aunque toda la comunidad eclesial participa en su elaboración.
\item Doble tarea: \textbf{anuncio} (visión global del hombre) y \textbf{denuncia} (del pecado de injusticia, defensa de los pobres).
\item Doble rasgo: \textbf{continuidad} (principios permanentes) + \textbf{renovación} (``taller siempre abierto'').
\item Línea histórica completa a memorizar: \textit{Rerum novarum (1891) $\rightarrow$ Quadragesimo anno (1931) $\rightarrow$ Mit brennender Sorge / Divini Redemptoris (1937) $\rightarrow$ Mater et Magistra (1961) $\rightarrow$ Pacem in terris (1963) $\rightarrow$ Gaudium et spes / Dignitatis humanae (1965) $\rightarrow$ Populorum progressio (1967) $\rightarrow$ Octogesima adveniens (1971) $\rightarrow$ Laborem exercens (1981) $\rightarrow$ Sollicitudo rei socialis (1987) $\rightarrow$ Centesimus annus (1991)}.
\end{enumerate}



\subsection{1.4 Francisco, \textit{Dilexit nos} (2024) --- Cap. I: ``La importancia del corazón''}

\noindent \textit{(Es el texto que figura en el programa 2026 para el Eje 1, temas ``Evangelio y vida'' y ``La caridad como fundamento''. La cátedra entregó además Dilexi te cap. IV, resumido en 1.1.)}\par\smallskip

\noindent \textbf{Idea central:} en un mundo ``líquido'', consumista y dominado por la tecnología, hace falta \textbf{volver al corazón}: el centro unificador de la persona, donde se decide lo importante y desde donde es posible el encuentro con los demás y con Dios. Solo desde el corazón puede cambiar el mundo; por eso la caridad (el amor) es el fundamento de toda la vida social.\par\smallskip

\subsubsection{A) Punto de partida (\S{}1--2)}
\begin{itemize}
\item ``Nos amó'' (Rm 8,37): Cristo nos amó primero, sin condiciones ni requisitos previos.
\item El símbolo del corazón sigue siendo válido: lo necesitamos cuando vivimos ``corriendo sin saber para qué'', como ``consumistas insaciables y esclavizados por los engranajes de un mercado'' al que no le interesa el sentido de la existencia.
\end{itemize}

\subsubsection{B) ¿Qué expresamos cuando decimos ``corazón''? (\S{}3--8)}
\begin{itemize}
\item En el griego clásico (\textit{kardia}), en Homero y Platón: centro corporal, anímico y espiritual; lugar donde se fraguan las decisiones; función ``sintetizadora'' de lo racional y lo tendencial.
\item El ser humano no es una suma de capacidades sino un mundo anímico-corpóreo con un \textbf{centro unificador} que da sentido y orientación.
\item En la Biblia: núcleo que está detrás de toda apariencia (discípulos de Emaús: ``¿no ardía nuestro corazón?'').
\item Es el \textbf{lugar de la sinceridad}, donde no se puede engañar ni disimular (Sansón y Dalila). ``Nada que valga la pena se construye sin el corazón''; la apariencia y la mentira solo ofrecen vacío (anécdota de las galletas huecas de la abuela, llamadas ``mentiras'').
\item Lleva a las preguntas decisivas: quién soy, qué busco, qué sentido quiero que tenga mi vida, quién soy frente a Dios.
\end{itemize}

\subsubsection{C) Volver al corazón (\S{}9--16)}
\begin{itemize}
\item Diagnóstico: ``sociedades de consumidores seriales que viven al día'', dominados por los ritmos y ruidos de la tecnología; el hombre corre el riesgo de \textbf{perder su centro}. ``Falta corazón''.
\item La desvalorización del corazón viene de lejos: racionalismo griego, idealismo, materialismo. La filosofía prefirió razón, voluntad o libertad, y el pensamiento desembocó en un \textbf{individualismo enfermizo}.
\item Si se devalúa el corazón se pierden las respuestas que la inteligencia sola no puede dar, el encuentro con los demás y la poesía.
\item Ejemplo literario: Stavroguin (Dostoyevski), según Guardini encarnación del mal porque \textbf{no tiene corazón}: no puede llegar a los demás. ``Es el corazón el que crea las posibilidades de encuentro''.
\item Todas las acciones deben ponerse bajo el ``dominio político'' del corazón; inteligencia y voluntad a su servicio.
\item ``Yo soy mi corazón'': es lo que me distingue y me pone en comunión con los demás. \textbf{Los algoritmos} muestran que pensamientos y decisiones son predecibles y manipulables; el corazón no.
\item ``Corazón'' es una palabra originaria que nombra a la persona como \textbf{totalidad} corpóreo-espiritual; ninguna ciencia la agota.
\end{itemize}

\subsubsection{D) El corazón que une los fragmentos (\S{}17--23)}
\begin{itemize}
\item Solo el corazón hace posibles vínculos auténticos y supera la \textbf{fragmentación del individualismo}. ``Anticorazón'' es una sociedad dominada por el narcisismo y la autorreferencia.
\item Conexión paradójica: se llega a ser uno mismo cuando se reconoce al otro.
\item María ``atesoraba todas estas cosas, ponderándolas en su corazón'': el corazón unifica la historia personal.
\item ``En el tiempo de la \textbf{inteligencia artificial} no podemos olvidar que para salvar lo humano hacen falta la poesía y el amor'': los pequeños detalles de la vida (el tenedor de las empanadas, la pelota de trapo) nunca estarán en los algoritmos.
\item El ser humano fue creado ante todo \textbf{para el amor}: amar y ser amado.
\item Las guerras, con complicidad o indiferencia de otros países, muestran que ``la sociedad mundial está perdiendo el corazón'' (las abuelas que lloran a sus nietos).
\item Pregunta decisiva: \textbf{¿tengo corazón?}
\end{itemize}

\subsubsection{E) El fuego (\S{}24--27)}
\begin{itemize}
\item Los Ejercicios de san Ignacio tienen por principio el \textit{affectus}: el orden de la vida nace del corazón, no de discursos racionales.
\item San Buenaventura: preguntar ``no a la luz, sino al fuego''. Newman: \textit{Cor ad cor loquitur} (el corazón habla al corazón).
\item Sin corazón, la práctica de la fe queda en ``moralismo autosuficiente''.
\end{itemize}

\subsubsection{F) El mundo puede cambiar desde el corazón (\S{}28--31)}
\begin{itemize}
\item ``Tomar en serio el corazón tiene \textbf{consecuencias sociales}'': los desequilibrios del mundo moderno se conectan con el desequilibrio que hunde sus raíces en el corazón humano (\textit{Gaudium et spes}).
\item El corazón unido al de Cristo es capaz del ``milagro social'': construir el Reino de amor y de justicia.
\item El corazón humano no es autosuficiente: es frágil y está herido; tiene \textbf{dignidad ontológica}, pero debe buscar una vida más digna y necesita el auxilio del amor divino.
\item Cristo es ``el corazón del mundo''. Oración final: que el mundo, que sobrevive entre guerras, desequilibrios socioeconómicos, consumismo y \textbf{uso antihumano de la tecnología}, recupere el corazón.
\end{itemize}

\subsubsection{Ideas clave para retener (Dilexit nos, cap. I)}
\begin{enumerate}
\item Corazón = centro unificador de toda la persona (no solo del alma), lugar de la sinceridad y de las decisiones.
\item Diagnóstico: sociedad líquida, consumista, tecnológica y narcisista; ``falta corazón''.
\item El corazón supera la fragmentación del individualismo y hace posible el encuentro.
\item Frente a los algoritmos y la inteligencia artificial: para salvar lo humano hacen falta poesía y amor.
\item Volver al corazón tiene consecuencias sociales: el mundo puede cambiar desde el corazón. La caridad es el fundamento.
\end{enumerate}



\clearpage\section{EJE N\textordmasculine{} 2 --- Del discurso a la propuesta social}

\subsection{2.1 Sorge --- Cap. VII, IX, X, XXIX}

\subsubsection{CAP. VII --- Del discurso a la propuesta social}

\medskip\noindent \textbf{1. Del ``discurso social'' a la ``propuesta social''}\par\nopagebreak
\begin{itemize}
\item El Evangelio no es solo un texto para leer, sino una orientación para la existencia. Por eso el discurso social no puede quedarse en afirmaciones abstractas: debe traducirse en principios de reflexión, criterios de juicio y orientaciones de acción.
\item La Iglesia interpreta los \textbf{``signos de los tiempos''} a la luz del Evangelio y de ahí formula una \textbf{``propuesta social''}: no es un modelo cerrado ni prefabricado, ni se limita a enumerar principios generales --- responde a situaciones históricas concretas y busca orientar éticamente la transformación social.
\end{itemize}

\medskip\noindent \textbf{2. ¿Por qué es necesaria una propuesta social?}\par\nopagebreak
\begin{itemize}
\item Las sociedades atraviesan una crisis de transición: riesgo de buscar solo soluciones técnicas inmediatas sin un \textbf{proyecto común}.
\item Problema central: la separación entre los problemas a resolver, los métodos usados y los valores que deberían orientar las soluciones.
\item Sin una cultura política sustentada en valores éticos compartidos, ninguna ingeniería institucional garantiza por sí sola una sociedad justa $\rightarrow$ se necesita consenso sobre valores, criterios de juicio, prioridades y objetivos comunes. Su ausencia lleva al \textbf{relativismo ético}.
\end{itemize}

\medskip\noindent \textbf{3. La propuesta social de la Iglesia}\par\nopagebreak
\begin{itemize}
\item Se resume en la idea de construir una \textbf{``civilización del amor''}, fundada en paz, solidaridad, justicia, libertad y dignidad de la persona.
\item La Iglesia no impone un proyecto político específico: ofrece un fundamento moral y espiritual para que creyentes y personas de buena voluntad construyan juntos, mediante el \textbf{diálogo} y no la imposición.
\end{itemize}

\medskip\noindent \textbf{4. Diálogo intercultural}\par\nopagebreak
\begin{itemize}
\item Necesario para construir una sociedad común. Problema histórico: confrontación entre cultura laica (razón como criterio fundamental de verdad) y cultura cristiana (razón abierta a una verdad trascendente).
\item El diálogo no implica que cada parte abandone su identidad, sino reconocer elementos de verdad en la posición del otro, distinguiendo el plano científico, filosófico, religioso y ético.
\end{itemize}

\medskip\noindent \textbf{5. Diálogo interreligioso}\par\nopagebreak
\begin{itemize}
\item No debe limitarse al encuentro creyentes/no creyentes; también incluye a las distintas religiones.
\item La doctrina social puede ser terreno común para promover dignidad humana, paz, lucha contra la pobreza y el hambre, cooperación y fraternidad.
\item Tiene fundamentación \textbf{teológica}, no solo estratégica: la humanidad tiene origen y destino comunes.
\end{itemize}

\noindent \textbf{Idea general del Cap. VII:} la propuesta social de la Iglesia es una \textbf{orientación ética y cultural}, no un modelo político cerrado. Núcleo: \textit{civilización del amor + dignidad humana + solidaridad + justicia + libertad + diálogo}.\par\smallskip



\subsubsection{CAP. IX --- Una ``gramática ética'' común}

\medskip\noindent \textbf{1. La necesidad de una ``gramática ética''}\par\nopagebreak
\begin{itemize}
\item No puede existir una verdadera propuesta social sin un conjunto mínimo de principios y valores éticos compartidos.
\item Comparación con la \textbf{gramática}: así como una lengua necesita reglas comunes para comunicarse, la sociedad necesita principios comunes para organizar la convivencia --- válidos para personas de diferentes creencias religiosas, culturales o políticas.
\item Cinco grandes principios: \textbf{personalista, solidaridad, subsidiariedad, bien común y una nueva concepción de la laicidad}.
\end{itemize}

\medskip\noindent \textbf{2. Repensar el concepto de laicidad}\par\nopagebreak
\begin{itemize}
\item Concepción clásica: separación Iglesia-Estado, autonomía de las instituciones civiles.
\item Evolución: laicidad ilustrada (razón, libertad e igualdad como únicos criterios) entendida como separación y contraposición; choque frontal con el \textit{Syllabus} de Pío IX (1864); Pío XI condena el ``laicismo'' pero reconoce la laicidad como distinción y autonomía recíproca (concordatos); Pío XII y Juan XXIII hablan de ``sana laicidad''; el Vaticano II la reconoce como \textbf{valor} con fundamento teológico: las realidades temporales tienen consistencia, leyes y fines propios (\textit{Gaudium et spes} 36).
\item Hoy la religión ya no se considera un fenómeno privado. El ``neolaicismo'' hace de la laicidad una ideología dogmática, una especie de ``religión de Estado''.
\item Nueva laicidad: respeto de la autonomía del Estado + respeto de la libertad religiosa + reconocimiento de la dimensión social de las religiones + diálogo + cooperación en la búsqueda del bien común.
\item \textbf{Laicidad $\neq$ laicismo}: la laicidad permite autonomía y colaboración mutua; el laicismo convierte la separación en ideología que excluye lo religioso del espacio público.
\end{itemize}

\medskip\noindent \textbf{3. Principio personalista}\par\nopagebreak
\begin{itemize}
\item La \textbf{persona humana es el centro} de la vida social, económica y política; vale por lo que es, no por lo que posee, produce o su utilidad --- tiene \textbf{dignidad trascendente}.
\item De aquí derivan: derechos fundamentales, libertad de conciencia, pensamiento, educación, asociación, derecho al trabajo.
\item La persona no puede ser instrumento de fines políticos, económicos o científicos. Aplicado a la vida humana: rechazo a su manipulación o experimentación --- la vida es un \textbf{continuum} cuya dignidad no depende de su grado de desarrollo o utilidad.
\end{itemize}

\medskip\noindent \textbf{4. Principio de solidaridad}\par\nopagebreak
\begin{itemize}
\item La persona es constitutivamente social; el reconocimiento de su dignidad individual lleva necesariamente a la solidaridad.
\item No es simple compasión, sino ``\textbf{determinación firme y perseverante de comprometerse con el bien común}''.
\item Solidaridad y justicia se complementan: la justicia da a cada uno lo suyo; la solidaridad impulsa a ocuparse activamente de los demás, especialmente los vulnerables --- puede ser instrumento de \textbf{liberación}.
\item Dimensión internacional: la interdependencia hace que la solidaridad no pueda limitarse al ámbito nacional; la paz es fruto de la solidaridad entre pueblos.
\item Evolución del principio: de la dimensión social de la persona y de la propiedad (León XIII, Pío XI) a la solidaridad como virtud moral (\textit{Sollicitudo rei socialis}). A la luz de la fe se vuelve caridad: gratuidad, perdón y reconciliación.
\item Tres criterios del Concilio: satisfacer primero las exigencias de la \textbf{justicia} (no dar como caridad lo que se debe por justicia); eliminar las \textbf{causas} y no solo los efectos de los males; ayudar de modo que quien recibe se vaya \textbf{liberando de la dependencia}.
\end{itemize}

\medskip\noindent \textbf{5. Principio de subsidiariedad}\par\nopagebreak
\begin{itemize}
\item Las instancias superiores no deben absorber funciones que pueden realizar bien las instancias inferiores; deben ayudar, coordinar, proteger y complementar, sin sustituir innecesariamente a personas, familias, asociaciones y comunidades.
\item Rechaza dos extremos: la \textbf{centralización excesiva} (burocracia, Estado asistencial) y el \textbf{abandono} de quienes necesitan apoyo.
\item Subsidiariedad implica \textbf{participación}: la ciudadanía no solo recibe decisiones desde arriba, participa activamente en la vida económica, social, cultural y política.
\end{itemize}

\medskip\noindent \textbf{6. Principio del bien común}\par\nopagebreak
\begin{itemize}
\item Finalidad que integra los principios anteriores; rechaza tanto la interpretación individualista (suma de intereses particulares) como la colectivista (absorbido en el interés de la colectividad).
\item Es el conjunto de condiciones sociales que permiten a personas y comunidades alcanzar su desarrollo integral: incluye dignidad, justicia, educación, salud, vivienda, trabajo, cultura, libertad, vida y ambiente.
\item Tiene una \textbf{dimensión global}: pobreza, desigualdad, degradación ambiental, drogas, enfermedades, paz --- superan la capacidad de un solo Estado y exigen cooperación internacional.
\item El bien común coincide más con la \textbf{calidad} de la vida humana que con la cantidad de bienes; incluye bienes como cultura, educación y dimensión espiritual.
\end{itemize}

\medskip\noindent \textbf{7. Una nueva laicidad en política}\par\nopagebreak
\begin{itemize}
\item Permite que creyentes y no creyentes colaboren políticamente sin renunciar a su identidad. Para los cristianos: respetar las reglas democráticas, no imponer la religión al Estado, traducir sus convicciones en argumentos comprensibles en el espacio público.
\item La participación política se entiende como \textbf{servicio}, no privilegio.
\end{itemize}

\medskip\noindent \textbf{Idea general del Cap. IX (tabla síntesis):}\par\nopagebreak

{\small\renewcommand{\arraystretch}{1.25}
\begin{longtable}{>{\raggedright\arraybackslash}p{0.142\linewidth}>{\raggedright\arraybackslash}p{0.788\linewidth}}
\toprule
\textbf{Principio} & \textbf{Idea central} \\ \midrule \endhead
Personalismo & La persona y su dignidad son el centro \\ \addlinespace[2pt]
Solidaridad & Compromiso firme con el bien común \\ \addlinespace[2pt]
Subsidiariedad & Las instancias superiores ayudan sin sustituir innecesariamente a las inferiores \\ \addlinespace[2pt]
Bien común & Condiciones sociales para el desarrollo integral de todos \\ \addlinespace[2pt]
Nueva laicidad & Autonomía, pluralismo, diálogo y cooperación \\ \addlinespace[2pt]
\bottomrule\end{longtable}}



\subsubsection{CAP. X --- Economía libre y solidaria}

\medskip\noindent \textbf{1. El capitalismo no ha vencido}\par\nopagebreak
\begin{itemize}
\item El siglo XX dejó el fracaso de dos modelos ideológicos: socialismo marxista y capitalismo liberal. El fracaso del socialismo no convierte al capitalismo en solución definitiva: también produjo desigualdad, exclusión, explotación, consumismo y concentración de poder. Se necesita un \textbf{nuevo modelo de desarrollo}.
\end{itemize}

\medskip\noindent \textbf{2-3. El capitalismo como ideología y su evolución histórica}\par\nopagebreak
\begin{itemize}
\item Distinción clave: capitalismo como \textbf{ideología} (concepción de sociedad, persona, trabajo, propiedad, consumo) vs. capitalismo como \textbf{sistema de producción}. La crítica de la Iglesia apunta principalmente a la ideología cuando coloca la economía por encima de la moral y la dignidad humana.
\item Etapas: \textbf{manchesteriano} (libertad económica extrema, mínima intervención estatal, dura explotación) $\rightarrow$ \textbf{monopolístico} (concentración de recursos y poder) $\rightarrow$ \textbf{consumista} (posguerra; el consumo como motor de la producción, generando necesidades artificiales, endeudamiento y pérdida de valores) $\rightarrow$ \textbf{societario} (el trabajador adquiere un papel más complejo en la empresa y el Estado social, aunque persisten problemas de propiedad y distribución).
\end{itemize}

\medskip\noindent \textbf{4. La industrialización}\par\nopagebreak
\begin{itemize}
\item Se distingue claramente entre \textbf{industrialización}, \textbf{ideología capitalista} y \textbf{sistema capitalista de producción}. La industrialización en sí puede tener aspectos positivos (productividad, progreso material); la crítica moral no va contra la tecnología, sino contra \textbf{cómo se organiza y para qué se usa} la producción.
\end{itemize}

\medskip\noindent \textbf{5. Capitalismo como sistema de producción}\par\nopagebreak
\begin{itemize}
\item La pregunta no es ``mercado sí / mercado no'', sino \textbf{qué reglas, valores y límites deben orientarlo}. El mercado debe estar subordinado a objetivos humanos y sociales: dignidad, justicia, trabajo, solidaridad, bien común, responsabilidad social. El beneficio puede ser legítimo pero no puede ser el criterio absoluto.
\item Elementos ``técnicos'' del sistema: propiedad privada de los medios de producción, división entre capital y trabajo, empresa, beneficio y mercado de libre concurrencia. La Iglesia no los discute en el plano técnico, sino que los juzga en el plano moral y pide orientarlos ética y jurídicamente.
\item Juan Pablo II (\textit{Centesimus annus}) invita a dejar de usar la palabra ``capitalismo'' y hablar de \textbf{``economía libre''}: ``una sociedad del trabajo libre, de la empresa y de la participación''.
\end{itemize}

\medskip\noindent \textbf{6. Una economía libre y solidaria}\par\nopagebreak
\begin{itemize}
\item Se puede superar la oposición entre libertad económica y solidaridad: el mercado debe funcionar dentro de un \textbf{marco ético, jurídico y político} (reglas democráticas, protección de trabajadores, límites a abusos, competencia justa, responsabilidad social). Libertad económica no significa ausencia de reglas.
\end{itemize}

\medskip\noindent \textbf{7. Globalización y economía postindustrial}\par\nopagebreak
\begin{itemize}
\item La caída del bloque comunista y el desarrollo tecnológico aceleraron la globalización de los mercados; la economía entra en etapa \textbf{postindustrial} (internacionalización, movilidad de capital, mercados financieros globales).
\item Fenómenos: ``riqueza sin naciones'', ``naciones sin riqueza'', hipercompetencia, deslocalización. Problema central: los Estados nacionales pierden capacidad de control sobre procesos globales, con riesgos para empleo, salarios y protección social.
\end{itemize}

\medskip\noindent \textbf{8. Sociedad del trabajo libre, la empresa y la participación}\par\nopagebreak
\begin{itemize}
\item La nueva economía coloca el \textbf{conocimiento y la capacidad humana} en el centro de la producción (no solo la fuerza física). La empresa debería entenderse como \textbf{comunidad de personas} orientada al bien común, no solo máquina de beneficios. La flexibilidad laboral no debe convertirse en precarización.
\end{itemize}

\medskip\noindent \textbf{9. Solidaridad económica internacional}\par\nopagebreak
\begin{itemize}
\item La interdependencia mundial exige una solidaridad que va más allá de la ayuda puntual: no basta transferir recursos si las estructuras siguen generando desigualdad. La cooperación internacional debe apuntar a desarrollo, justicia, participación y reducción de la pobreza.
\end{itemize}



\subsubsection{CAP. XXIX --- La crisis económica mundial: una oportunidad}

\medskip\noindent \textbf{1. La crisis como oportunidad}\par\nopagebreak
\begin{itemize}
\item La crisis económico-financiera iniciada con las hipotecas \textit{subprime} (2007) se interpreta como manifestación de una crisis más profunda: no solo financiera, sino económica, social, cultural, ética y de valores. Puede ser oportunidad para transformar el modelo económico.
\end{itemize}

\medskip\noindent \textbf{2. Naturaleza ``estructural'' de la crisis}\par\nopagebreak
\begin{itemize}
\item El problema no está solo en decisiones financieras puntuales: revela deficiencias más profundas del modelo económico y social (trabajo, familia, educación, instituciones, cultura, valores).
\item \textbf{Causas próximas:} comportamiento irresponsable de los operadores financieros (créditos \textit{subprime}, bancos convertidos en agencias especuladoras). \textbf{Causa última:} una economía de mercado que produce riqueza pero prescinde de las reglas éticas fundamentales (primado del bien común, destino universal de los bienes, prioridad del trabajo sobre el capital); en el fondo, la cultura económica neoliberal y su utilitarismo.
\item La guía de la globalización está en manos de pocos sujetos sin legitimación democrática (organismos financieros, grandes transnacionales).
\end{itemize}

\medskip\noindent \textbf{3. Implicaciones éticas}\par\nopagebreak
\begin{itemize}
\item Consecuencias humanas: desempleo, precariedad, pobreza, desigualdad, exclusión. Crítica a una lógica económica en la que \textbf{el beneficio y el capital terminan ocupando un lugar superior al trabajo y a la persona}. Las reglas del mercado por sí solas no bastan: se necesitan criterios éticos y jurídicos que limiten la lógica del beneficio cuando perjudica la dignidad humana y el bien común.
\end{itemize}

\medskip\noindent \textbf{4. El problema de la cultura consumista}\par\nopagebreak
\begin{itemize}
\item No se puede resolver una crisis generada por la lógica del consumo incentivando aún más consumo. La riqueza debe vincularse a trabajo, productividad, desarrollo humano, responsabilidad y sostenibilidad.
\end{itemize}

\medskip\noindent \textbf{5. La crisis como oportunidad de reforma}\par\nopagebreak
\begin{itemize}
\item Posibilidad de construir una \textbf{economía social de mercado}, reformando empresa, Estado social, mercados, finanzas y relaciones laborales.
\item La \textbf{empresa} debe verse como comunidad de personas, considerando no solo accionistas sino también trabajadores, clientes, proveedores, comunidad, generaciones futuras y medio ambiente.
\end{itemize}

\medskip\noindent \textbf{6. Tecnología y empresa}\par\nopagebreak
\begin{itemize}
\item Transición desde el trabajo basado en fuerza física hacia conocimiento, inteligencia, innovación y participación --- oportunidad para democratizar la empresa.
\end{itemize}

\medskip\noindent \textbf{7. Reforma del Estado social y sociedad civil}\par\nopagebreak
\begin{itemize}
\item El Estado no debe ser el único protagonista: se necesita interacción equilibrada entre \textbf{Estado, mercado y sociedad civil}. La sociedad civil debe tener papel activo (asociaciones, cooperación, economía solidaria, voluntariado) --- conecta directamente con el \textbf{principio de subsidiariedad} del Cap. IX.
\end{itemize}

\medskip\noindent \textbf{8. Conclusión del Cap. XXIX}\par\nopagebreak
\begin{itemize}
\item La salida de la crisis exige una \textbf{gobernanza equitativa y solidaria} de la economía. No se trata de eliminar el mercado sino de orientarlo con reglas, responsabilidad, solidaridad, participación y respeto a la dignidad humana.
\end{itemize}

\subsubsection{Ideas clave para retener (síntesis rápida --- Sorge Cap. VII, IX, X, XXIX)}
\begin{enumerate}
\item Secuencia lógica del bloque: crisis de la sociedad $\rightarrow$ necesidad de una \textbf{propuesta social} $\rightarrow$ búsqueda de una \textbf{``gramática ética'' común} $\rightarrow$ cinco principios (personalismo, solidaridad, subsidiariedad, bien común, nueva laicidad) $\rightarrow$ orientación ética de la política y de la economía $\rightarrow$ economía libre y solidaria $\rightarrow$ mercado orientado al bien común.
\item \textbf{Civilización del amor}: paz + justicia + libertad + solidaridad + dignidad humana --- núcleo de la propuesta social.
\item \textbf{Gramática ética}: personalismo (¿a quién sirve la sociedad? a la persona) + solidaridad (¿cómo relacionarse? comprometiéndose con el bien común) + subsidiariedad (¿cómo organizar las instituciones? ayudando sin sustituir) + bien común (¿objetivo? desarrollo integral de todos) + nueva laicidad (¿cómo convivir en el pluralismo? con autonomía, respeto y diálogo).
\item Distinción clave: \textbf{capitalismo como ideología} (sistema de valores que puede subordinar la persona a la lógica económica) vs. \textbf{capitalismo como sistema} (forma de organizar producción, orientable mediante normas éticas) --- y distinción adicional entre \textbf{industrialización} y ambos.
\item Fórmula: \textbf{economía libre + solidaria}, dentro de un marco ético, jurídico y político.
\item La crisis de 2007-2008 como crisis \textbf{estructural} (no solo financiera) y como \textbf{oportunidad} de reforma: economía social de mercado, empresa como comunidad de personas, equilibrio Estado-mercado-sociedad civil.
\end{enumerate}



\subsection{2.2 Compendio DSI --- Cap. IV y VII}

\subsubsection{CAPÍTULO CUARTO --- ``Los principios de la doctrina social de la Iglesia''}

\medskip\noindent \textbf{I. Significado y unidad (\S{}160--163)}\par\nopagebreak
\begin{itemize}
\item Por su \textbf{permanencia en el tiempo y universalidad de significado}, son el primer y fundamental parámetro para interpretar y valorar los fenómenos sociales. Tienen un significado profundamente \textbf{moral}: remiten a los fundamentos últimos de la vida social.
\item Los \textbf{principios permanentes} de la doctrina social son: la \textbf{dignidad de la persona humana} (ya visto en el capítulo precedente del Compendio, fundamento de todos los demás), el \textbf{bien común}, la \textbf{subsidiariedad} y la \textbf{solidaridad}.
\item Brotan del encuentro entre el mensaje evangélico (mandamiento del amor + justicia) y los problemas de la vida social.
\item Tienen carácter \textbf{general y fundamental}: van desde las relaciones interpersonales hasta las relaciones entre pueblos y Naciones. Deben apreciarse en su \textbf{unidad, conexión y articulación} --- no se pueden usar de forma aislada o desconectada.
\end{itemize}

\medskip\noindent \textbf{II. El principio del bien común (\S{}164--170)}\par\nopagebreak
\begin{itemize}
\item Deriva de la dignidad, unidad e igualdad de todas las personas; a él debe referirse todo aspecto de la vida social \textbf{para encontrar plenitud de sentido}. Es ``la dimensión social y comunitaria del bien moral'': la búsqueda del \textbf{sentido y la verdad} de las formas de vida social.
\item Definición clásica: \textbf{``el conjunto de condiciones de la vida social que hacen posible a las asociaciones y a cada uno de sus miembros el logro más pleno y más fácil de la propia perfección''}.
\item No es la simple suma de bienes particulares: es indivisible, y solo se alcanza, acrecienta y custodia \textbf{juntos}.
\item \textbf{Responsabilidad de todos}: nadie está exento de colaborar según sus capacidades. Incluye servicios esenciales (alimentación, vivienda, trabajo, educación, salud, información, libertad religiosa) y cooperación internacional.
\item \textbf{Tareas de la comunidad política}: el Estado debe garantizar cohesión y organización; en democracia, el gobierno debe buscar el bien de \textbf{toda} la comunidad, incluidas las minorías, no solo seguir a la mayoría.
\item El bien común \textbf{no es un fin autárquico}: tiene valor en relación a los fines últimos de la persona y su dimensión trascendente (Dios como fin último); una visión puramente materialista lo reduciría a mero bienestar socioeconómico.
\end{itemize}

\medskip\noindent \textbf{III. El destino universal de los bienes (\S{}171--184)}\par\nopagebreak
\begin{itemize}
\item \textbf{Origen y significado}: Dios destinó la tierra a todos los hombres y pueblos; es el \textbf{``primer principio de todo el ordenamiento ético-social''}, un derecho \textbf{natural y originario}, prioritario respecto a cualquier sistema socioeconómico --- todos los demás derechos (incluida la propiedad y el libre comercio) están \textbf{subordinados} a él.
\item No significa que todo esté a disposición de todos sin regulación: requiere intervenciones normativas y ordenamiento jurídico.
\item \textbf{Destino universal y propiedad privada}: la propiedad privada nace del trabajo humano y es ``ampliación de la libertad humana''; la Iglesia nunca la aceptó como derecho absoluto e intocable --- está \textbf{subordinada} al destino universal de los bienes, del que es un instrumento (medio, no fin). Tiene \textbf{función social}: el propietario no debe considerar sus bienes como exclusivamente suyos.
\item Se reconoce también la \textbf{propiedad comunitaria} (pueblos indígenas) como forma legítima, junto a la individual.
\item \textbf{Destino universal de los bienes y opción preferencial por los pobres}: reafirma la opción preferencial por los pobres como ``forma especial de primacía en el ejercicio de la caridad cristiana'' (cita textual de \textit{Sollicitudo rei socialis}, ya vista en \textit{Dilexi te} \S{}87). La limosna es también un acto de \textbf{justicia}, no solo de caridad: ``cuando damos a los pobres las cosas indispensables no les hacemos liberalidades personales, sino que les devolvemos lo que es suyo''.
\end{itemize}

\medskip\noindent \textbf{IV. El principio de subsidiariedad (\S{}185--188)}\par\nopagebreak
\begin{itemize}
\item Presente desde \textit{Rerum novarum}; protege a la persona, familia y cuerpos intermedios de los abusos de las instancias superiores.
\item Cita clave de \textit{Quadragesimo anno}: no se puede quitar a las comunidades menores lo que ellas pueden hacer y dárselo a una sociedad mayor --- toda acción social debe \textbf{ayudar} a los miembros del cuerpo social, no destruirlos o absorberlos.
\item Contrasta con la centralización, burocratización y el \textbf{``Estado asistencial''} (que provoca pérdida de energías humanas y aumento excesivo del aparato público).
\item El Estado puede ejercer una \textbf{función de suplencia} en situaciones excepcionales (grave desequilibrio o injusticia), pero no debe prolongarse más de lo estrictamente necesario.
\end{itemize}

\medskip\noindent \textbf{V. La participación (\S{}189--191)}\par\nopagebreak
\begin{itemize}
\item Consecuencia característica de la subsidiariedad: el ciudadano contribuye directa o indirectamente a la vida cultural, económica, política y social. Es un \textbf{deber}, no solo un derecho.
\item Es uno de los pilares de los ordenamientos democráticos: \textbf{``toda democracia debe ser participativa''}.
\item Riesgos: formas de participación insuficientes (abstención, desinterés), regímenes totalitarios que niegan este derecho, o burocracia excesiva que lo anula de hecho.
\end{itemize}

\medskip\noindent \textbf{VI. El principio de solidaridad (\S{}192--196)}\par\nopagebreak
\begin{itemize}
\item Surge de la creciente \textbf{interdependencia} entre personas y pueblos (telecomunicaciones, comercio, información), que convive con fuertes desigualdades entre países desarrollados y en desarrollo.
\item Doble dimensión: \textbf{principio social} (las ``estructuras de pecado'' deben transformarse en estructuras de solidaridad mediante leyes y reglas justas) y \textbf{virtud moral} --- cita clave de \textit{Sollicitudo rei socialis}: no es ``un sentimiento superficial'', sino \textbf{``la determinación firme y perseverante de empeñarse por el bien común... para que todos seamos verdaderamente responsables de todos''}.
\item Vínculos: solidaridad--bien común, solidaridad--destino universal de los bienes, solidaridad--paz.
\item Su cumbre es la vida de Jesucristo, ``solidario con la humanidad hasta la muerte de cruz'': a la luz de la fe, la solidaridad se reviste de gratuidad, perdón y reconciliación.
\end{itemize}

\medskip\noindent \textbf{VII. Los valores fundamentales de la vida social (\S{}197--203)}\par\nopagebreak
\begin{itemize}
\item Sobre la libertad (\S{}200): su plenitud consiste en ``la \textbf{capacidad de disponer de sí mismo con vistas al auténtico bien}'', en el horizonte del bien común; incluye la capacidad de rechazar lo moralmente negativo.
\item Relación de \textbf{reciprocidad} entre principios y valores. Los valores esenciales son: \textbf{la verdad, la libertad, la justicia y el amor}.
\item \textbf{La verdad}: vivir en la verdad ordena la convivencia; exige compromiso educativo, especialmente en comunicación pública y economía.
\item \textbf{La libertad}: signo de la dignidad humana; no debe entenderse de forma individualista o como autonomía arbitraria --- se ejercita en relaciones reguladas por la verdad y la justicia.
\item \textbf{La justicia}: ``dar a Dios y al prójimo lo que les es debido''. Formas clásicas: \textbf{conmutativa, distributiva y legal}, más la \textbf{justicia social} (desarrollo de la justicia general con dimensión mundial y estructural). ``Por sí sola, la justicia no basta... puede llegar a negarse a sí misma si no se abre a la fuerza más profunda que es el amor'' $\rightarrow$ de ahí la fórmula \textbf{``Opus solidaritatis pax''} (la paz, fruto de la solidaridad).
\end{itemize}

\medskip\noindent \textbf{VIII. La vía de la caridad (\S{}204--208)}\par\nopagebreak
\begin{itemize}
\item La caridad es criterio supremo y universal de toda la ética social --- ``la vía más excelente'' (1 Co 12,31).
\item Los valores de verdad, justicia y libertad nacen de la fuente de la caridad. La caridad \textbf{presupone y trasciende} la justicia: la justicia ``corrige'' con el amor (\textit{summum ius, summa iniuria} --- la justicia extrema puede ser injusticia).
\item La caridad se convierte en \textbf{caridad social y política}: no se limita a relaciones interpersonales, sino que interviene sobre la comunidad social y política en su conjunto, buscando el bien posible para todos.
\end{itemize}



\subsubsection{CAPÍTULO SÉPTIMO --- ``La vida económica''}

\medskip\noindent \textbf{I. Aspectos bíblicos (\S{}323--327)}\par\nopagebreak
\begin{itemize}
\item Antiguo Testamento: doble postura ante la riqueza --- apreciación de los bienes necesarios (la abundancia como bendición), pero condena de su \textbf{mal uso} (estafas, usura, explotación --- tradición profética). La pobreza aceptada con espíritu religioso predispone a la confianza en Dios.
\item Jesús asume esta tradición y la lleva a plenitud: instaura el Reino de Dios, que hace posible una convivencia en justicia, fraternidad, solidaridad y ``compartir''.
\item La actividad económica debe entenderse como respuesta agradecida a la vocación que Dios da al hombre (cultivar y custodiar el jardín --- Gn 2,15): la buena administración de los bienes es obra de justicia (parábola de los talentos).
\end{itemize}

\medskip\noindent \textbf{b) La riqueza existe para ser compartida (\S{}328--329)}\par\nopagebreak
\begin{itemize}
\item Los bienes, aun poseídos legítimamente, conservan \textbf{destino universal}; toda acumulación indebida es inmoral.
\item Los Padres de la Iglesia (Clemente de Alejandría, san Juan Crisóstomo, san Basilio, san Gregorio Magno) insisten en que el rico es solo \textbf{administrador}, no dueño absoluto: las riquezas se deben hacer circular; dar a quien tiene necesidad es ``pagar una deuda'', no un favor.
\end{itemize}

\medskip\noindent \textbf{II. Moral y economía (\S{}330--335)}\par\nopagebreak
\begin{itemize}
\item La economía y la moral tienen principios propios, pero \textbf{no están separadas}: la actividad económica está sujeta a un fin dado por Dios Creador (cita de \textit{Quadragesimo anno}, Pío XI).
\item \textbf{``El hombre es el autor, el centro y el fin de toda la vida económico-social''} (\textit{Gaudium et spes}). La eficiencia económica y el desarrollo solidario son fines \textbf{compatibles}, no alternativos.
\item Es necesario combatir, con justicia y caridad, las \textbf{``estructuras de pecado''} que generan pobreza y subdesarrollo.
\item Valoración del \textbf{``capitalismo''}: si se entiende como sistema que reconoce el papel positivo de la empresa, el mercado y la propiedad privada dentro de un sólido marco ético-jurídico, la respuesta es positiva (mejor llamarlo ``economía de mercado'' o ``economía libre''); si se entiende como sistema donde la libertad económica no está encuadrada en ningún marco jurídico ni ético, la respuesta es \textbf{absolutamente negativa} --- cita clave de \textit{Centesimus annus}.
\end{itemize}

\medskip\noindent \textbf{III. Iniciativa privada y empresa (\S{}336--345)}\par\nopagebreak
\begin{itemize}
\item La libertad económica es un \textbf{derecho inalienable}; negarla o limitarla ``en nombre de una pretendida igualdad'' destruye la subjetividad creativa del ciudadano.
\item \textbf{a) La empresa y sus fines}: la empresa debe servir al bien común mediante la producción de bienes y servicios; tiene también una \textbf{función social} (encuentro, colaboración, valoración de capacidades). No es solo ``sociedad de capitales'', sino también ``\textbf{sociedad de personas}'' (capital + trabajo). El beneficio es indicador válido del buen funcionamiento, pero no garantiza por sí solo que la empresa sirva bien a la sociedad --- puede haber ``balances correctos'' con trabajadores humillados en su dignidad. La \textbf{usura} está moralmente condenada como forma extrema de injusticia.
\item \textbf{b) El papel del empresario y del dirigente}: implica virtudes (diligencia, prudencia, fiabilidad, lealtad); no puede regirse solo por criterios financieros --- debe respetar la dignidad de los trabajadores y favorecer a la familia (especialmente a las madres trabajadoras).
\end{itemize}

\medskip\noindent \textbf{IV. Instituciones económicas al servicio del hombre (\S{}346--360)}\par\nopagebreak
\begin{itemize}
\item \textbf{a) El papel del libre mercado}: instrumento eficaz para asignar recursos y responder a necesidades, pero \textbf{no se autolegitima}: debe subordinarse a fines morales. Riesgo de ``idolatría del mercado'' --- hay bienes que no son ni pueden ser simples mercancías.
\item \textbf{b) La acción del Estado}: debe combinar subsidiariedad (crear condiciones favorables a la iniciativa privada) y solidaridad (proteger a la parte más débil). \textit{``La solidaridad sin subsidiariedad puede degenerar en asistencialismo; la subsidiariedad sin solidaridad puede alimentar el localismo egoísta''} --- fórmula clave para el parcial. El Estado puede intervenir directamente en situaciones excepcionales (función de suplencia) y regula mediante finanza pública justa (impuestos como deber de solidaridad).
\item \textbf{c) La función de los cuerpos intermedios}: hay bienes colectivos que no dependen solo del mercado ni son competencia exclusiva del Estado --- ahí actúan las organizaciones sin fines de lucro y la sociedad civil, en ``democracia económica''.
\item \textbf{d) Ahorro y consumo}: la elección de invertir ``en un lugar y no en otro... es siempre una opción moral y cultural''. Crítica al \textbf{consumismo}, que orienta hacia el ``tener'' en vez del ``ser''.
\end{itemize}

\medskip\noindent \textbf{V. Las ``res novae'' en economía (\S{}361--373)}\par\nopagebreak
\begin{itemize}
\item \textbf{Solidaridad entre generaciones (\S{}367):} es moralmente ilícito y económicamente contraproducente descargar los costos actuales sobre las generaciones futuras.
\item \textbf{e) Una gran obra educativa y cultural (\S{}375--376):} la economía es solo una dimensión de la actividad humana; si se absolutiza, la causa está en un sistema sociocultural que ignora la dimensión ética y religiosa. Hace falta educar a los consumidores, a los productores y a los profesionales de los medios.
\item \textbf{a) La globalización: oportunidades y riesgos}: proceso de creciente integración económico-financiera mundial. Genera esperanzas (reducción de costos, aceleración del desarrollo) pero también riesgos: aumento de desigualdades entre países, incluso dentro de los países ricos. Objetivo: \textbf{``globalizar la solidaridad''}, ``una globalización sin dejar a nadie al margen''. La globalización \textbf{no debe ser un nuevo tipo de colonialismo} --- debe respetar la diversidad cultural y religiosa de los pueblos.
\item \textbf{b) El sistema financiero internacional}: los mercados financieros son útiles para financiar la economía real, pero corren el riesgo de volverse \textbf{autorreferenciales}, desconectados de la economía real, generando inestabilidad --- especialmente peligroso para los países más frágiles.
\item \textbf{c) La función de la comunidad internacional}: ante la pérdida de centralidad de los Estados nacionales frente a los mercados globales, se necesita una \textbf{autoridad/comunidad internacional} capaz de dirigir estos procesos con criterios no solo económicos sino también morales.
\item \textbf{d) Un desarrollo integral y solidario}: tarea fundamental --- ``promover a todos los hombres y a todo el hombre'' (\textit{Populorum progressio}), con distribución equitativa de recursos y conciencia de la interdependencia que une a los pueblos en un único destino.
\end{itemize}

\subsubsection{Ideas clave para retener (síntesis rápida --- Compendio Cap. IV y VII)}
\begin{enumerate}
\item \textbf{Cuatro principios permanentes} de la DSI: dignidad de la persona (base) + \textbf{bien común + subsidiariedad + solidaridad}, más los \textbf{cuatro valores}: verdad, libertad, justicia y amor/caridad (caridad como ``vía más excelente'', que corrige y trasciende a la justicia).
\item El \textbf{destino universal de los bienes} es el ``primer principio del ordenamiento ético-social'': la propiedad privada existe pero está subordinada a él (es medio, no fin) y tiene función social.
\item \textbf{Subsidiariedad + solidaridad deben ir siempre juntas}: subsidiariedad sin solidaridad = localismo egoísta; solidaridad sin subsidiariedad = asistencialismo --- fórmula muy citable para el examen.
\item La actividad económica tiene fundamento bíblico: los bienes son dones de Dios que deben administrarse y compartirse (destino universal, no acumulación indebida).
\item Distinción clave en la valoración del ``capitalismo'': positiva si hay marco ético-jurídico (mejor llamarlo ``economía libre/de mercado''); negativa si la libertad económica no está encuadrada en ningún marco --- mismo esquema que en Sorge Cap. X (ideología vs. sistema).
\item El mercado es un instrumento \textbf{útil pero no autolegitimado}: necesita el Estado (subsidiariedad + solidaridad) y los cuerpos intermedios (sociedad civil) actuando de forma complementaria.
\item La globalización trae oportunidades y riesgos; el criterio moral es \textbf{``globalizar la solidaridad''}, sin que se convierta en un nuevo colonialismo cultural.
\item Puente directo con Sorge Cap. VII, IX, X y XXIX: los mismos principios (personalismo/dignidad, solidaridad, subsidiariedad, bien común) y la misma crítica a la autonomización del mercado y de las finanzas respecto de la ética.
\end{enumerate}



\clearpage\section{EJE N\textordmasculine{} 3 --- La presencia social de la Iglesia}

\subsection{3.1 Francisco --- Fratelli Tutti, Cap. V y VI}

\subsubsection{CAPÍTULO V --- La mejor política}

\medskip\noindent \textbf{1. Idea central}\par\nopagebreak
\begin{itemize}
\item Para construir una comunidad mundial basada en la fraternidad y la amistad social hace falta una \textbf{``mejor política''}: una actividad orientada al \textbf{bien común}, no una lucha por el poder, los intereses económicos o el beneficio electoral. Debe estar al servicio de las personas, especialmente de las más vulnerables, y buscar una sociedad abierta donde todos tengan un lugar.
\end{itemize}

\medskip\noindent \textbf{2. Crítica al populismo y al liberalismo}\par\nopagebreak
\begin{itemize}
\item \textbf{Popular no es populista:} la noción de \textbf{pueblo} es legítima y necesaria (sin ella no hay ``democracia'', gobierno del pueblo). Pueblo es una categoría \textbf{``mítica''}, no lógica ni mística: identidad común hecha de lazos sociales y culturales, proceso lento hacia un proyecto común; es una categoría \textbf{abierta}, que incorpora al diferente.
\item El liderazgo popular degenera en \textbf{``insano populismo''} cuando instrumentaliza la cultura del pueblo para un proyecto personal y la perpetuación en el poder, avasalla las instituciones o cae en el \textbf{inmediatismo} (responder a exigencias para ganar votos). Los planes asistenciales son solo ``respuestas pasajeras''.
\item \textbf{``El gran tema es el trabajo'':} lo verdaderamente popular es asegurar a todos la posibilidad de desarrollar sus capacidades; ayudar con dinero es solución provisoria.
\item \textbf{Límites de las visiones liberales:} ven la sociedad como mera suma de intereses e ignoran la fragilidad humana (concupiscencia). ``El mercado solo no resuelve todo'', aunque se lo presente como ``dogma de fe neoliberal''; el ``derrame'' no resuelve la inequidad.
\item El desprecio hacia los débiles puede aparecer bajo distintas formas políticas:
\begin{itemize}
\item \textbf{Populismo}: manipula al pueblo y usa sus necesidades como instrumento político en lugar de trabajar por su dignidad y bienestar real.
\item \textbf{Liberalismo}: puede priorizar los intereses económicos de los sectores poderosos.
\end{itemize}
\item En ambos casos aparece la misma dificultad: pensar una sociedad abierta, incluir a los más débiles, respetar las culturas y poner la política al servicio del bien común. Francisco no opone simplemente dos modelos, sino que critica \textbf{toda} forma política que termine excluyendo grupos.
\end{itemize}

\medskip\noindent \textbf{3. La política como servicio al bien común}\par\nopagebreak
\begin{itemize}
\item La política debe recuperar una visión amplia y de largo plazo, sin limitarse a resultados electorales, encuestas, marketing político o soluciones inmediatas.
\item Pregunta clave: \textbf{qué sociedad se está construyendo y qué consecuencias tendrán las decisiones para las generaciones futuras}. El político es un \textbf{constructor}, capaz de transformar las condiciones que producen exclusión.
\end{itemize}

\medskip\noindent \textbf{4. La política frente a la economía}\par\nopagebreak
\begin{itemize}
\item \textbf{El poder internacional:} ante el debilitamiento de los Estados nacionales frente a lo económico-financiero, hacen falta instituciones internacionales más fuertes, una reforma de la ONU y de la arquitectura financiera, el respeto de los acuerdos (\textit{pacta sunt servanda}: ``la fuerza del derecho'' y no ``el derecho de la fuerza'') y preferencia por los acuerdos multilaterales.
\item \textbf{La política no debe quedar subordinada a la economía}: esta puede aportar recursos y eficiencia, pero no puede determinar por sí sola qué es bueno para la sociedad.
\item La política necesita: visión integral, diálogo con distintas disciplinas, capacidad de pensar más allá de los intereses económicos inmediatos, y orientar la economía hacia el bien común. Sin política, la economía queda dominada por una lógica puramente eficientista.
\end{itemize}

\medskip\noindent \textbf{5. Una política orientada hacia los más vulnerables}\par\nopagebreak
\begin{itemize}
\item Debe ocuparse de quienes sufren exclusión: trata de personas, explotación sexual, trabajo esclavo, tráfico de drogas y armas, terrorismo, crimen organizado, hambre, falta de vivienda y de salud.
\item Las instituciones deben ser \textbf{efectivas}, no solo declarativas. \textbf{Principio fundamental: la dignidad de la persona debe estar por encima de la eficiencia, los intereses económicos y los cálculos políticos.}
\end{itemize}

\medskip\noindent \textbf{6. Derechos humanos y política internacional}\par\nopagebreak
\begin{itemize}
\item Todavía no existe una verdadera universalización de los derechos humanos básicos. Prioridades \textbf{mínimas impostergables}: combatir el hambre, garantizar alimentación, vivienda y salud, combatir la trata de personas, proteger a los vulnerables.
\end{itemize}

\medskip\noindent \textbf{7. La caridad política}\par\nopagebreak
\begin{itemize}
\item La política es ``una altísima vocación, una de las formas más preciosas de la caridad, porque busca el bien común''. \textbf{Pueblo y persona son términos correlativos.}
\item \textbf{Amor ``elícito'' y amor ``imperado'':} el primero son los actos dirigidos directamente a personas y pueblos; el segundo, los que impulsan a crear instituciones más sanas, regulaciones más justas y estructuras más solidarias. Ejemplo: ayudar a un anciano a cruzar un río es caridad; el político que construye un puente también ejerce la caridad.
\item La caridad política es \textbf{amor preferencial por los últimos} y necesita la luz de la verdad (razón y fe) para no quedar en sentimentalismo.
\item \textbf{Movimientos populares:} son ``sembradores de cambio'', \textbf{``poetas sociales''} que gestan formas de economía popular y producción comunitaria; sin ellos ``la democracia se atrofia''. Políticas sociales \textbf{con} los pobres y \textbf{de} los pobres, no solo \textbf{hacia} los pobres.
\item La política también puede ser una forma de \textbf{caridad}: buscar el bien de todos, aun con renuncias personales o partidarias. Implica escuchar, aceptar diferencias, buscar puntos de encuentro, trabajar por el bien común superando intereses particulares.
\item Imagen del \textbf{poliedro}: una sociedad de múltiples diferencias que, sin desaparecer, conviven y contribuyen a una unidad más amplia.
\end{itemize}

\medskip\noindent \textbf{8. Diferencia frente a uniformidad}\par\nopagebreak
\begin{itemize}
\item La fraternidad no elimina las diferencias. La uniformidad también puede destruir la riqueza cultural y social.
\item \textbf{Diferencia $\neq$ división}; \textbf{Unidad $\neq$ uniformidad}. La unidad verdadera integra las diferencias, no las elimina.
\end{itemize}

\medskip\noindent \textbf{9. El político y la dignidad humana}\par\nopagebreak
\begin{itemize}
\item El político debe preocuparse menos por su imagen y más por la transformación real de la sociedad. No ``¿cuánto apoyo tengo?'' sino ``¿qué hice avanzar? ¿qué problemas solucioné? ¿cuánta paz social generé?''. La política se mide por su \textbf{fecundidad social}.
\end{itemize}



\subsubsection{CAPÍTULO VI --- Diálogo y amistad social}

\medskip\noindent \textbf{1. Idea central}\par\nopagebreak
\begin{itemize}
\item El \textbf{diálogo} es herramienta fundamental para construir amistad social: acercarse, expresarse, escuchar, mirar al otro, conocerse, comprender, buscar puntos de contacto. No es solo intercambiar opiniones, sino construir una relación que genera comprensión y cooperación.
\end{itemize}

\medskip\noindent \textbf{2. El diálogo como alternativa}\par\nopagebreak
\begin{itemize}
\item Se presenta como tercera vía entre dos extremos: \textbf{indiferencia egoísta $\leftrightarrow$ protesta violenta}. Permite que distintas generaciones, culturas y sectores participen en un proyecto común.
\end{itemize}

\medskip\noindent \textbf{3. Diálogo no es intercambio superficial de opiniones}\par\nopagebreak
\begin{itemize}
\item Se diferencia del intercambio digital en redes sociales, que puede generar mucha información sin auténtico encuentro. Requiere apertura, escucha, paciencia, búsqueda de la verdad, reconocimiento del otro y disposición a modificar la propia posición.
\end{itemize}

\medskip\noindent \textbf{4. Diálogo y verdad}\par\nopagebreak
\begin{itemize}
\item \textbf{``El relativismo no es la solución'':} bajo una supuesta tolerancia, deja que los poderosos interpreten los valores según su conveniencia. El diálogo es el camino para reconocer valores que están \textbf{más allá de todo consenso} y son \textbf{nunca negociables}; el fundamento es la dignidad inalienable de todo ser humano, que la razón puede reconocer.
\item No todas las opiniones tienen automáticamente el mismo valor, ni la verdad depende del consenso. Existen \textbf{valores fundamentales que no deberían quedar sometidos a consensos circunstanciales}.
\item Fórmula: \textbf{diálogo $\rightarrow$ búsqueda de la verdad}, pero \textbf{consenso $\neq$ verdad absoluta}.
\end{itemize}

\medskip\noindent \textbf{5. Cultura del encuentro}\par\nopagebreak
\begin{itemize}
\item La sociedad como \textbf{poliedro}: distintas caras y perspectivas forman una unidad sin perder sus particularidades.
\item Implica: reconocer diferencias, buscar puntos de contacto, tender puentes, incluir a las periferias, aprender de los demás, construir proyectos compartidos. Nadie es inútil o prescindible; las periferias pueden aportar perspectivas que los centros de poder no perciben.
\end{itemize}

\medskip\noindent \textbf{6. La paz social}\par\nopagebreak
\begin{itemize}
\item No consiste en eliminar conflictos o silenciar diferencias. Francisco: la verdadera paz es \textbf{trabajosa y artesanal}.
\item Una paz superficial (impuesta, silenciada, controlada, excluyente) es frágil. La verdadera paz requiere encuentro, integración, reconocimiento de las diferencias, diálogo y participación --- más lenta, pero más sólida.
\end{itemize}

\medskip\noindent \textbf{7. Reconocer al otro}\par\nopagebreak
\begin{itemize}
\item De ahí la necesidad de un \textbf{pacto social} que sea también \textbf{``pacto cultural''}: respetar las diversas cosmovisiones (por ejemplo, los pueblos originarios) y aceptar ceder algo por el bien común (``realismo dialogante'').
\item La amistad social exige reconocer el derecho del otro a existir, expresarse, ser diferente y participar. Negar el valor de determinados grupos es una forma de violencia menos visible, pero que igual destruye los vínculos sociales.
\end{itemize}

\medskip\noindent \textbf{8. Amabilidad como práctica social}\par\nopagebreak
\begin{itemize}
\item No es solo buenos modales: prestar atención, escuchar, ofrecer tiempo, evitar la agresividad, facilitar la convivencia. Como práctica social, transforma la forma de debatir y facilita el consenso.
\end{itemize}

\subsubsection{Conceptos clave y esquema (Fratelli Tutti V-VI)}

{\small\renewcommand{\arraystretch}{1.25}
\begin{longtable}{>{\raggedright\arraybackslash}p{0.204\linewidth}>{\raggedright\arraybackslash}p{0.726\linewidth}}
\toprule
\textbf{Concepto} & \textbf{Significado} \\ \midrule \endhead
Mejor política & Política orientada al bien común y a la dignidad humana \\ \addlinespace[2pt]
Caridad política & Servicio al otro mediante decisiones orientadas al bien común \\ \addlinespace[2pt]
Diálogo & Encuentro basado en escucha, comprensión y búsqueda de puntos comunes \\ \addlinespace[2pt]
Cultura del encuentro & Forma de vida que integra diferencias y construye puentes \\ \addlinespace[2pt]
Poliedro & Imagen de una sociedad plural donde las diferencias forman una unidad \\ \addlinespace[2pt]
Paz social & Resultado de un proceso de encuentro e integración, no de silenciar conflictos \\ \addlinespace[2pt]
\bottomrule\end{longtable}}

\medskip\noindent \textbf{Esquema de continuidad Cap. V $\rightarrow$ Cap. VI:}\par\nopagebreak
\begin{esquema}
\noindent\hspace*{0.0em}FRATERNIDAD $\rightarrow$ AMISTAD SOCIAL $\rightarrow$ MEJOR POLÍTICA $\rightarrow$ BIEN COMÚN $\rightarrow$ DIGNIDAD HUMANA\par
\noindent\hspace*{0.0em}$\rightarrow$ DIÁLOGO $\rightarrow$ CULTURA DEL ENCUENTRO $\rightarrow$ INTEGRACIÓN DE LAS DIFERENCIAS $\rightarrow$ PAZ SOCIAL\par
\end{esquema}
\begin{itemize}
\item \textbf{Cap. V} responde: ¿cómo debe actuar la política para construir una sociedad fraterna? $\rightarrow$ orientándose al bien común, protegiendo la dignidad humana, incluyendo a los vulnerables.
\item \textbf{Cap. VI} responde: ¿cómo pueden convivir personas y grupos diferentes? $\rightarrow$ mediante el diálogo, la escucha y la cultura del encuentro.
\end{itemize}



\subsection{3.2 Sorge --- Cap. XIII y XXVIII}

\subsubsection{CAP. XIII --- De la propuesta a la ``presencia social''}

\medskip\noindent \textbf{1. De la propuesta a la acción}\par\nopagebreak
\begin{itemize}
\item El discurso social de la Iglesia tiene una \textbf{``propuesta social''} propia: líneas de fondo (principios, criterios de juicio, orientaciones operativas) para elaborar un proyecto cultural cristianamente inspirado, compartible con todos los hombres de buena voluntad.
\item Esta propuesta ya tiene nombre: \textbf{``La civilización del amor''} --- no es un slogan ni un reclamo moralista, sino un proyecto de sociedad, un desafío histórico (cita la Instrucción \textit{Libertatis conscientia}).
\item El discurso social está \textbf{``esencialmente orientado a la acción''}: se traduce en ``audaces programas de acción'' para la liberación de millones de personas. La Iglesia siempre ha reivindicado este derecho-deber de presencia y acción social (cita \textit{Mater et magistra}: la Iglesia tiene derecho y deber no solo de tutelar los principios éticos, sino de intervenir autorizadamente en el orden temporal).
\item Aunque la Iglesia no tenga un ``modelo técnico'' de sociedad, sí puede formular una propuesta social y auspiciar que los cristianos la traduzcan en realidad.
\end{itemize}

\medskip\noindent \textbf{El paso de la ``propuesta'' a la ``acción'' se da a través de dos formas de presencia:}\par\nopagebreak
\begin{enumerate}
\item La \textbf{presencia social y política} propia y específica de los \textbf{fieles laicos}.
\item Una presencia social ``diversa'' de la comunidad cristiana \textbf{en cuanto tal} (propia de quienes ejercen un ministerio oficial --- obispos, sacerdotes).
\end{enumerate}

\medskip\noindent \textbf{1. El papel insustituible de los fieles laicos}\par\nopagebreak
\begin{itemize}
\item El laico recibe la llamada al compromiso social y político \textbf{directamente de Cristo en el bautismo}, no ``por delegación'' del obispo o del párroco.
\item Es necesidad urgente contar con laicos maduros (hombres y mujeres) comprometidos con la evangelización y el servicio cristiano al mundo (cita Juan Pablo II, \textit{Christifideles laici}: ``a nadie le es lícito permanecer ocioso'').
\item Sin un laicado maduro y comprometido, la Iglesia no puede hacerse socialmente presente hoy de modo adecuado.
\item El mundo no es una prisión de la que huir, sino el \textbf{``espacio teológico''} donde todos los cristianos ---y en primer lugar los laicos--- están llamados a vivir y trabajar para anunciar la salvación.
\item Imágenes evangélicas de la sal, la luz y la levadura (Juan Pablo II) tienen aplicación específica a los laicos: expresan su plena inserción y participación en el mundo, y la originalidad de su participación destinada a la difusión del Evangelio.
\item El laico está llamado a ser \textbf{``hombre de la síntesis''}: síntesis entre fe e historia, fe y cultura, fe y política --- la misma síntesis que la Iglesia entera debe realizar en su ``presencia social''.
\item Recuadro \textit{Compendio} \S{}544-545: la espiritualidad laica no es ni ``espiritualismo intimista'' ni ``activismo social'', sino una síntesis vital que confiere unidad, significado y esperanza a la existencia.
\end{itemize}

\medskip\noindent \textbf{2. El movimentismo católico}\par\nopagebreak
\begin{itemize}
\item Así como el ``discurso social'' siempre acabó convirtiéndose en ``propuesta social'', la propuesta social nunca se quedó en mera propuesta, sino que se transformó en \textbf{presencia} y \textbf{acción social}: el fenómeno del \textbf{movimentismo católico}.
\item Sorge distingue \textbf{cuatro fases} del movimentismo católico, correspondientes a las cuatro etapas del discurso social (ya vistas en el Cap. I de Sorge):
\end{itemize}

{\small\renewcommand{\arraystretch}{1.25}
\begin{longtable}{>{\raggedright\arraybackslash}p{0.250\linewidth}>{\raggedright\arraybackslash}p{0.231\linewidth}>{\raggedright\arraybackslash}p{0.430\linewidth}}
\toprule
\textbf{Fase} & \textbf{Período aprox.} & \textbf{Características} \\ \midrule \endhead
\textbf{Movimiento católico} & desde fines s. XVIII & Forma embrionaria; los laicos actúan bajo la guía rigurosa del papa y del clero, oponiéndose a la laicización del Estado. Dentro de esta fase surgen los \textbf{``intransigentes''} (nostálgicos del viejo orden, contra los ``errores'' del liberalismo) y los \textbf{``católicos liberales''} (buscan conciliar fe y progreso histórico). \\ \addlinespace[2pt]
\textbf{Catolicismo social} & desde \textit{Rerum novarum} (1891) & León XIII compromete a la Iglesia en la solución de la ``cuestión obrera''; nacen los movimientos democristianos; se supera el clima de ``ciudadela sitiada''; se consolida (1\textordfeminine{} mitad s. XX) el movimiento democristiano. También surge la \textbf{Acción Católica} (organización laical más moderna, sustituyendo a la antigua ``Obra de los Congresos''). \\ \addlinespace[2pt]
\textbf{Mundo católico} & tras la 2\textordfeminine{} Guerra Mundial (con Pío XII) & Ante el ``peligro comunista'', el movimentismo asume las características de un verdadero ``mundo católico'': se movilizan los laicos en múltiples asociaciones, todas apoyando al partido democristiano --- nace el \textbf{``colateralismo''} entre Iglesia, movimientos católicos y partido católico. Esta situación dura hasta el Vaticano II y solo queda superada tras la caída del Muro de Berlín (1989). \\ \addlinespace[2pt]
\textbf{Área cultural de inspiración cristiana} & actual (fase todavía por definir) & Tras el ocaso del ``mundo católico'', se busca un nuevo modo de presencia social adecuado a la transición hacia el tercer milenio: reafirmación de valores éticos y culturales desde un \textbf{área ``intermedia''}, autónoma tanto del área eclesial como del área política, gestionada por fieles laicos, abierta a creyentes y no creyentes. Todavía es una propuesta, ``un camino por inventar''. \\ \addlinespace[2pt]
\bottomrule\end{longtable}}

\medskip\noindent \textbf{3. Dos conceptos de política y de compromiso político}\par\nopagebreak
\begin{itemize}
\item El Concilio subraya la importancia de distinguir dos conceptos de ``política'' (comentario recogido en el documento final de Puebla):
\begin{itemize}
\item \textbf{Política con ``P'' mayúscula} (acepción amplia, ``cultura política''): corresponde precisar los valores fundamentales de toda comunidad, conciliando igualdad con libertad, autoridad pública con participación, soberanía nacional con solidaridad internacional. En este sentido amplio, entran también las actividades sociales, asistenciales, de voluntariado --- no elaboradas por partidos ni por el Estado, sino nacidas espontáneamente del compromiso social. En este sentido, la política interesa también a la Iglesia y a sus pastores.
\item \textbf{política con ``p'' minúscula} (acepción corriente, ``praxis'' política): la de los partidos, sindicatos, gobierno, instituciones --- traducción ``técnica'' de valores y necesidades en opciones operativas. \textbf{Esta política partidista es terreno propio de los laicos}, no de los pastores. Estos últimos deben abandonar toda política-partidista para conservar la libertad de evangelizar lo político.
\end{itemize}
\item Ambas son ``dos aspectos distintos, pero inseparables entre sí, de la misma realidad'': cuando los partidos e instituciones pierden en su praxis la conexión con su bagaje cultural, están destinados a morir, ``como si perdieran el alma''.
\end{itemize}

\noindent \textbf{Idea general del Cap. XIII:} el ``discurso social'' desemboca necesariamente en ``propuesta social'' (civilización del amor) y esta, a su vez, en ``presencia'' y ``acción social'' concretas --- protagonizadas específicamente por los laicos, a través de las sucesivas formas históricas del movimentismo católico, distinguiendo siempre Política (cultura) de política (praxis partidista).\par\smallskip



\subsubsection{CAP. XXVIII --- Política y ``valores no negociables''}

\noindent \textit{(Nota: este capítulo tiene un fuerte anclaje en el contexto italiano contemporáneo --- crisis política italiana, cardenal Martini, etc. --- que Sorge usa como caso de estudio para una reflexión de alcance general sobre católicos y política.)}\par\smallskip

\medskip\noindent \textbf{1. Cómo se plantea hoy el problema de los valores en política}\par\nopagebreak
\begin{itemize}
\item El problema se agudiza porque la cuestión social se ha convertido en \textbf{``cuestión antropológica''} (ver Sorge Cap. VI / Benedicto XVI), y porque los católicos militan hoy en todos los partidos, conviviendo con tradiciones culturales diversas que no comparten la misma comprensión de los valores que el Magisterio considera \textbf{``no negociables''}.
\item Pregunta de fondo: en un contexto cultural y políticamente fragmentado, ¿pueden los católicos comportarse a la vez como auténticos cristianos y sinceros demócratas? ¿pueden crear unidad respetando la diversidad sin traicionar sus valores?
\item Diagnóstico de crisis política (ejemplo italiano): predominio de intereses personales o corporativos sobre el bien común, desconfianza hacia la política, alto abstencionismo, tensión entre poderes del Estado, riesgo de ``autoritarismo populista'' (quien tiene el poder ``manda'' en vez de ``gobernar'').
\item A esto se suma la crisis eclesial interna: necesidad de purificación, de una fe madura que se traduzca en obras (``sin obras, la fe está muerta'' --- St 2,17), y de repensar el rol de los laicos.
\end{itemize}

\medskip\noindent \textbf{2. La cuestión de los llamados ``valores no negociables''}\par\nopagebreak
\begin{itemize}
\item Los católicos no están tan divididos sobre el \textbf{primado de los valores} en sí ni sobre la necesidad de comprometerse social y políticamente, sino sobre \textbf{el modo} de defender los valores ``no negociables'' en la situación actual del país.
\item Dos posturas:
\begin{enumerate}
\item Reafirmar primero los valores y principios ``absolutos'' y ``no negociables'' de manera clara, y luego dialogar sin complejos con las dinámicas culturales del tiempo.
\item Sin restar importancia al testimonio de los valores evangélicos, no empezar por reafirmarlos de forma abstracta, sino por \textbf{compartir concretamente} los problemas sociales, culturales y morales de la gente, para avanzar juntos gradualmente hacia la verdad completa, confiando en el Espíritu Santo.
\end{enumerate}
\item No es solo una cuestión de método: Sorge cita a Benedicto XVI (\textit{Deus caritas est}): el testimonio de la caridad (también en el plano cultural, social y político) \textbf{no debe ser un medio para el proselitismo} --- el amor es gratuito, no busca imponer la fe a los demás.
\item Clave (cardenal Martini): el \textbf{bien común no es una definición filosófica abstracta}, sino algo que debe perseguirse de manera concreta y proporcional a las situaciones históricas reales; su consecución pasa por el convencimiento, la paciencia, la afirmación progresiva y gradual de los valores --- a veces con ``duras renuncias'' en nombre de una concordia civil superior, con la mirada puesta en un bien superior.
\item Distinción importante: los \textbf{principios y valores son siempre en sí mismos absolutos y ``no negociables''}, pero su \textbf{traducción histórica} está sujeta a las condiciones de tiempo y lugar, al consenso democrático, al cambio de las costumbres sociales y a la mediación política. No se trata de confiar al consenso de la mayoría la verificación de la verdad de un valor, sino de asumir con autonomía una responsabilidad respecto al cambio de las costumbres --- tarea propia de la ética política.
\item Propuesta de síntesis: un \textbf{``neopersonalismo solidario y laico''}, que profundice y actualice los valores (libertad, igualdad, solidaridad, paz, dignidad de la persona) ya presentes en las Constituciones democráticas modernas --- combinando el personalismo (tradición católico-democrática), la solidaridad (tradición socialista) y la laicidad (tradición liberal-democrática). Este programa concreto implica: repensar el ejercicio y tutela de los derechos humanos, la reforma de las instituciones, las reglas del mercado, el Estado social, la defensa del trabajo, el rechazo del racismo, la política migratoria, la tutela del ambiente, la dimensión mundial del desarrollo y la colaboración internacional.
\end{itemize}

\medskip\noindent \textbf{3. Nuevas perspectivas para los católicos en política}\par\nopagebreak
\begin{itemize}
\item La tarea de traducir la fe en opciones eficaces corresponde sobre todo a los \textbf{fieles laicos}, en diálogo con los hombres de buena voluntad, respetando la laicidad y las reglas democráticas.
\item Los cristianos deben ser \textbf{ciudadanos activos y responsables}, no espectadores pasivos: deben ``ensuciarse las manos'' junto a todos los ciudadanos de buena voluntad para dar un alma ética a la política.
\item Deben mantener la \textbf{libertad de militar} en distintos partidos democráticos (siempre con coherencia respecto a los valores evangélicos), evitando depender de líderes carismáticos únicos.
\item El ``primer partido'' de los cristianos debe ser el \textbf{territorio}: el barrio, la ciudad, la provincia, la región donde viven --- trasladando la confrontación desde la contraposición electoral estéril hacia el debate sobre el proyecto de sociedad y los valores comunes.
\item Se requiere diálogo entre los propios católicos, comprensión recíproca entre las distintas opciones políticas legítimas, y sobre todo una renovada \textbf{formación} de la conciencia cívica y moral --- especialmente de los más jóvenes: laicos maduros, espiritualmente formados y profesionalmente preparados.
\end{itemize}

\subsubsection{Ideas clave para retener (síntesis rápida --- Sorge Cap. XIII y XXVIII)}
\begin{enumerate}
\item Secuencia del ``discurso social'': discurso $\rightarrow$ \textbf{propuesta social} (``civilización del amor'') $\rightarrow$ \textbf{presencia y acción social} concreta (protagonizada por los laicos).
\item Doble forma de presencia: la de los \textbf{laicos} (política, con ``p'' minúscula incluida) y la \textbf{específica de la comunidad cristiana como tal} (pastores, sin política partidista).
\item \textbf{Cuatro fases del movimentismo católico}, en paralelo a las cuatro fases del discurso social: movimiento católico $\rightarrow$ catolicismo social $\rightarrow$ mundo católico $\rightarrow$ área cultural de inspiración cristiana (actual, aún por definir).
\item Distinción Política (P mayúscula, cultura de valores) / política (p minúscula, praxis partidista) --- ambas inseparables entre sí.
\item \textbf{Valores ``no negociables''}: son absolutos en sí mismos, pero su traducción histórica es gradual, dialogada y mediada políticamente --- no queda sujeta al consenso de la mayoría en cuanto a su verdad, pero sí en cuanto a su aplicación concreta.
\item Propuesta de síntesis cultural: \textbf{``neopersonalismo solidario y laico''} (personalismo + solidaridad + laicidad).
\item Los laicos deben ser protagonistas: ciudadanía activa, formación de la conciencia cívica, compromiso territorial concreto.
\end{enumerate}



\subsection{3.3 Compendio DSI --- Cap. VIII}

\subsubsection{CAPÍTULO OCTAVO --- ``La comunidad política''}

\medskip\noindent \textbf{I. Aspectos bíblicos (\S{}377--383)}\par\nopagebreak
\begin{itemize}
\item El rey de Israel es \textbf{elegido y consagrado (ungido) por Yahvéh} y debe \textbf{hacer visible su señorío} y su designio de salvación: defender a los débiles y asegurar la justicia.
\item \textbf{El señorío de Dios}: Israel reconoce originalmente solo el señorío de Yahvéh; cuando pide un rey, Samuel advierte los riesgos del poder despótico. El rey ideal (culminando en David y, plenamente, en Jesucristo) debe ser defensor de los débiles y garante de la justicia.
\item \textbf{Jesús y la autoridad política}: rechaza el poder opresivo y despótico (``los jefes de las naciones las tiranizan...''), pero nunca rechaza directamente a las autoridades de su tiempo. Ante el tributo al César: ``dar a Dios lo que es de Dios'' --- condena implícitamente la divinización del poder temporal, pero no niega su legitimidad. Jesús rechaza el mesianismo político de dominación: Él vino ``a servir y dar su vida'', no a dominar.
\item \textbf{Las primeras comunidades cristianas}: San Pablo (Rm 13) pide sumisión por razones de conciencia a la autoridad, en cuanto está al servicio de Dios para el bien; San Pedro exhorta a una obediencia libre y responsable. Cuando el poder se autodiviniza y exige sumisión absoluta, se convierte en la ``Bestia'' del Apocalipsis --- frente a esto, la resistencia de los mártires.
\item Conclusión bíblica: el poder político procede de Dios, es parte del orden creado, y su significado auténtico es el \textbf{servicio}, no el dominio.
\end{itemize}

\medskip\noindent \textbf{II. El fundamento y el fin de la comunidad política (\S{}384--391)}\par\nopagebreak
\begin{itemize}
\item \textbf{La persona humana es el fundamento y el fin de la convivencia política}: por su apertura a la Trascendencia y a los demás, la vida social es una dimensión esencial del hombre, no una ``sobrecarga accidental''.
\item La comunidad política existe para lograr el \textbf{bien común}, un fin que de otro modo sería inalcanzable.
\item El \textbf{pueblo} no es una masa amorfa manipulable, sino un conjunto de personas que comparten vida y valores, con conciencia y responsabilidad propias. Las \textbf{minorías} tienen derecho a existir, mantener su cultura, lengua y religión; y también deberes, especialmente cooperar al bien común del Estado donde están insertas --- el terrorismo como vía de reivindicación es siempre injustificable.
\item \textbf{Derechos humanos}: la comunidad política tiende al bien común cuando defiende y promueve los derechos fundamentales; estos son norma objetiva anterior al Estado. Doble tarea complementaria: evitar privilegios de algunos grupos, pero también evitar que ``defender los derechos de todos'' impida el desarrollo pleno de cada uno.
\item \textbf{Amistad civil}: la convivencia adquiere su pleno significado cuando se funda en la \textbf{amistad civil y la fraternidad}, no solo en el derecho (que regula intereses). La amistad es desinterés, donación, disponibilidad al otro --- actuación más auténtica del principio de fraternidad (inseparable de libertad e igualdad). La justicia es ``medida mínima del amor''; el ideal es la \textbf{civilización del amor}.
\end{itemize}

\medskip\noindent \textbf{III. La autoridad política (\S{}393--405)}\par\nopagebreak
\begin{itemize}
\item Los valores morales esenciales no se fundan en mayorías de opinión: son elementos de la \textbf{ley natural}, punto de referencia normativo de la \textbf{ley civil}. La ley humana es ley en cuanto conforme a la recta razón (\S{}397--398).
\item \textbf{Fundamento}: necesaria por la naturaleza social del hombre; deriva en última instancia de Dios como autor del orden social. Debe garantizar el bien común sin suplantar la libertad de personas y grupos. El \textbf{pueblo es el sujeto de la autoridad} (soberanía popular), que transfiere su ejercicio a representantes pero conserva la facultad de control y sustitución.
\item \textbf{La autoridad como fuerza moral}: debe guiarse por la ley moral, no por criterios puramente sociológicos o de fuerza. Cita clave: si se niega la existencia de una ley moral objetiva y de Dios, esos preceptos se desintegran. Las leyes injustas dejan de ser leyes: ``se convierten en un acto de violencia'' (Santo Tomás).
\item \textbf{Objeción de conciencia}: derecho humano elemental; el ciudadano no está obligado a obedecer leyes contrarias al orden moral, a los derechos fundamentales o al Evangelio --- quien objeta debe estar protegido de sanciones penales, legales, disciplinarias o económicas.
\item \textbf{Derecho de resistencia}: legítimo ante violación grave y reiterada de derechos fundamentales, pero el recurso a las armas solo se justifica con condiciones estrictas: violaciones ciertas/graves/prolongadas, agotamiento de otros recursos, sin provocar males peores, esperanza fundada de éxito, sin mejor alternativa previsible --- siempre preferible la \textbf{resistencia pasiva}.
\item \textbf{Infligir penas}: el Estado tiene derecho y deber de imponer penas proporcionadas; doble finalidad: reinserción del culpable y justicia reconciliadora. Prohibición absoluta de la \textbf{tortura}, en cualquier circunstancia. Se valora positivamente la creciente aversión mundial a la \textbf{pena de muerte}: aunque la enseñanza tradicional no la excluye en teoría como único medio de legítima defensa social, hoy tales casos son ``muy raros, por no decir prácticamente inexistentes'' --- se prefieren métodos incruentos, más conformes a la dignidad humana.
\end{itemize}

\medskip\noindent \textbf{IV. El sistema de la democracia (\S{}406--416)}\par\nopagebreak
\begin{itemize}
\item \textbf{Principio moral para los medios de comunicación (\S{}416):} la \textbf{persona y la comunidad humana son el fin y la medida} del uso de los medios; el bien de las personas no se realiza al margen del bien común.
\item La Iglesia \textbf{aprecia la democracia} en la medida en que asegura participación, control y sustitución pacífica de gobernantes --- pero exige un Estado de derecho y una recta concepción de la persona (cita clave de \textit{Centesimus annus}).
\item \textbf{Los valores y la democracia}: una democracia auténtica no es solo respeto formal de reglas, sino aceptación convencida de valores (dignidad, derechos, bien común). El mayor riesgo: el \textbf{relativismo ético} --- sin verdad última, las convicciones se instrumentalizan para el poder: ``una democracia sin valores se convierte con facilidad en un totalitarismo visible o encubierto''. La democracia es un \textbf{instrumento}, no un fin: su moralidad depende de los fines que persigue y los medios que usa.
\item \textbf{Instituciones}: se valora la \textbf{división de poderes} (Estado de derecho: ``es soberana la ley, no la voluntad arbitraria''); elecciones libres y rendición de cuentas como mecanismos de control.
\item \textbf{Componente moral de la representación}: quienes gobiernan deben ejercer virtudes de servicio (paciencia, modestia, moderación, caridad, generosidad). La \textbf{corrupción política} es de las deformaciones más graves: traiciona la moral y la justicia social, genera desconfianza y clientelismo.
\item \textbf{Instrumentos de participación}: los partidos deben interpretar las aspiraciones de la sociedad civil orientándolas al bien común, ser democráticos internamente; el referéndum es otra herramienta de participación directa.
\item \textbf{Información y democracia}: la información es ``servicio del bien común''; derecho a una información fundada en verdad, libertad, justicia y solidaridad. Riesgo: concentración de medios y vínculos estrechos entre poder político, financiero y mediático.
\end{itemize}

\medskip\noindent \textbf{V. La comunidad política al servicio de la sociedad civil (\S{}417--420)}\par\nopagebreak
\begin{itemize}
\item Distinción entre \textbf{comunidad política} y \textbf{sociedad civil}: la sociedad civil es el conjunto de relaciones y recursos culturales/asociativos relativamente autónomos del ámbito político y económico.
\item \textbf{Primado de la sociedad civil}: la comunidad política está al servicio de la sociedad civil (y en definitiva, de las personas), no al revés --- es la sociedad civil la que justifica la existencia de la comunidad política.
\item \textbf{Principio de subsidiariedad aplicado}: la vida democrática debe crecer desde el tejido social; se valora el \textbf{``tercer sector''} (voluntariado, cooperación privada-social, distinto del Estado y del mercado) como espacio privilegiado para la dimensión social de la persona y para recomponer una ética pública basada en solidaridad y diálogo fraterno.
\end{itemize}

\medskip\noindent \textbf{VI. El Estado y las comunidades religiosas (\S{}421--427)}\par\nopagebreak
\begin{itemize}
\item \textbf{A) La libertad religiosa como derecho humano fundamental}: proclamada por \textit{Dignitatis humanae} (Vaticano II); ``la verdad no se impone de otra manera que por la fuerza de la misma verdad''. No es un derecho ilimitado: tiene justos límites determinados por la prudencia política según el bien común y el orden moral objetivo. Un Estado puede reconocer especialmente a una comunidad religiosa por sus vínculos históricos, sin que esto discrimine a otros grupos religiosos.
\item \textbf{B) Iglesia católica y comunidad política --- Autonomía e independencia}: Iglesia y comunidad política son de naturaleza diferente (fines distintos), \textbf{independientes y autónomas} cada una en su terreno (Vaticano II). La Iglesia no tiene competencia sobre la estructura política ni valora programas políticos salvo por sus implicaciones religiosas y morales.
\item \textbf{Colaboración}: la autonomía recíproca no excluye la colaboración, ya que ambas sirven a la vocación personal y social de los mismos hombres.
\item La Iglesia reivindica: libertad de expresión, enseñanza y evangelización; libertad de culto público; libertad de organización interna; libertad de elección y nombramiento de sus ministros; libertad de construir edificios religiosos; libertad de asociarse con fines religiosos, educativos, culturales, de salud y caritativos.
\item Existen formas jurídicas estables (concordatos, acuerdos) para prevenir y resolver conflictos entre Iglesia y Estado.
\end{itemize}

\subsubsection{Ideas clave para retener (síntesis rápida --- Compendio Cap. VIII)}
\begin{enumerate}
\item La autoridad política tiene fundamento en la naturaleza social del hombre y, en última instancia, en Dios; su ejercicio auténtico es \textbf{servicio}, no dominio (fundamento bíblico: Jesús ``vino a servir'').
\item La \textbf{persona humana es fundamento y fin} de la comunidad política; el pueblo no es masa manipulable.
\item La convivencia política se completa con la \textbf{amistad civil y la fraternidad} (no solo el derecho) --- ideal: la civilización del amor.
\item La autoridad es fuerza \textbf{moral}: las leyes injustas dejan de ser leyes. De ahí derivan la \textbf{objeción de conciencia} (siempre legítima) y el \textbf{derecho de resistencia} (con condiciones estrictas, y siempre preferible la vía pasiva).
\item Rechazo absoluto de la \textbf{tortura}; posición muy restrictiva sobre la \textbf{pena de muerte} (hoy prácticamente nunca necesaria).
\item La democracia se aprecia como sistema, pero \textbf{no es un fin en sí misma}: sin valores compartidos (contra el relativismo ético) degenera en totalitarismo encubierto.
\item \textbf{Primado de la sociedad civil} sobre la comunidad política --- y el \textbf{principio de subsidiariedad} aplicado también aquí (valoración del ``tercer sector'').
\item Relación Iglesia-Estado: \textbf{autonomía e independencia mutua + colaboración} --- ni confusión de roles ni separación excluyente. Libertad religiosa como derecho humano fundamental, con límites solo por bien común y orden moral objetivo.
\item Puente con Sorge Cap. XIII: la misma distinción entre ``Política'' (P mayúscula, cultura/valores --- donde interviene la Iglesia) y ``política'' (p minúscula, praxis partidista --- terreno de los laicos) es coherente con la autonomía Iglesia/Estado del Compendio.
\item Puente con Fratelli Tutti V-VI: la ``mejor política'' de Francisco (bien común, dignidad, diálogo) actualiza y da lenguaje contemporáneo a estos mismos principios clásicos del Compendio (autoridad como servicio, democracia con valores, sociedad civil, amistad social).
\end{enumerate}



\clearpage\section{EJE N\textordmasculine{} 4 --- Debates sociales: el pensamiento social de la Iglesia en diálogo con las ciencias}

\subsection{4.1 Compendio DSI --- Cap. III: ``La persona humana y sus derechos''}

\subsubsection{I. Doctrina social y principio personalista (\S{}105--107)}
\begin{itemize}
\item La Iglesia ve en cada hombre la \textbf{imagen viva de Dios}, con una dignidad ``incomparable e inalienable''. Cristo, con su encarnación, ``se ha unido en cierto modo con todo hombre''.
\item Toda la vida social es expresión de su protagonista: \textbf{la persona humana}, que es \textbf{sujeto, fundamento y fin} de la vida social.
\item \textbf{Principio personalista:} toda la doctrina social se desarrolla a partir del principio que afirma la \textbf{inviolable dignidad de la persona humana}. El hombre es ``el corazón y el alma'' de la enseñanza social.
\end{itemize}

\subsubsection{II. La persona humana ``imago Dei'' (\S{}108--123)}
\medskip\noindent \textbf{a) Criatura a imagen de Dios (\S{}108--114)}\par\nopagebreak
\begin{itemize}
\item El hombre es criatura querida por Dios por sí misma, \textbf{capaz de Dios} (\textit{capax Dei}): capaz de conocerse, poseerse, darse libremente y entrar en comunión.
\item Dimensión relacional y social inscrita en la naturaleza: ``no es bueno que el hombre esté solo''.
\item \textbf{Hombre y mujer} tienen la misma dignidad y son de igual valor; su reciprocidad (``unidad de los dos'') es imagen de Dios.
\item La vida humana es \textbf{sagrada e inviolable}; el hombre es custodio de la vida de los otros.
\item Relación con las demás criaturas: \textbf{señorío responsable}, no dominio arbitrario.
\item Relación consigo mismo: interioridad, el ``corazón''.
\end{itemize}

\medskip\noindent \textbf{b) El drama del pecado (\S{}115--119)}\par\nopagebreak
\begin{itemize}
\item El pecado original rompe la comunión con Dios, consigo mismo, con los demás y con la creación: raíz de toda alienación personal y social.
\item \textbf{Todo pecado es personal y a la vez social.} Algunos son sociales por su objeto: agresión directa al prójimo, contra la justicia, contra los derechos, contra el bien común.
\item El pecado social no diluye la responsabilidad personal.
\item \textbf{Estructuras de pecado:} nacen de pecados personales, se consolidan y condicionan la conducta de los hombres. Dos actitudes en su raíz (\textit{Sollicitudo rei socialis}): \textbf{el afán de ganancia exclusiva} y \textbf{la sed de poder}, ``a cualquier precio''.
\end{itemize}

\medskip\noindent \textbf{c) Universalidad del pecado y universalidad de la salvación (\S{}120--123)}\par\nopagebreak
\begin{itemize}
\item \textbf{Realismo cristiano:} ve los abismos del pecado, pero a la luz de la esperanza, mayor que todo mal, dada por la redención de Cristo. Ni optimismo ingenuo ni pesimismo.
\end{itemize}

\subsubsection{III. La persona humana y sus múltiples dimensiones (\S{}124--151)}
\noindent La doctrina social rechaza las \textbf{visiones reductivas} del hombre: ni individuo absoluto, ni simple célula de un organismo; ni puro espíritu ni pura materia.\par\smallskip

\medskip\noindent \textbf{A. La unidad de la persona (\S{}127--129)}\par\nopagebreak
\begin{itemize}
\item El hombre es \textbf{unidad de alma y cuerpo}. Por su corporeidad unifica en sí los elementos del mundo material; por su espiritualidad supera la totalidad de las cosas.
\item Ni \textbf{espiritualismo} que desprecia la realidad del cuerpo, ni \textbf{materialismo} que considera el espíritu una mera manifestación de la materia.
\end{itemize}

\medskip\noindent \textbf{B. Apertura a la trascendencia y unicidad (\S{}130--134)}\par\nopagebreak
\begin{itemize}
\item \textbf{Abierta a la trascendencia:} el hombre está abierto al infinito (Dios) y a todos los seres creados; sale de sí hacia el otro.
\item \textbf{Única e irrepetible:} existe como un ``yo'' capaz de autocomprenderse, autoposeerse y autodeterminarse. Pero \textbf{la persona no se reduce a su conciencia o a su inteligencia}: es la persona la que está en la base de esos actos, aunque falten.
\item \textbf{Respeto de la dignidad humana:} una sociedad justa solo se realiza en el respeto de la dignidad trascendente de la persona. ``El orden social y su progreso deben subordinarse al bien de la persona... \textbf{y no al contrario}'' (\textit{Gaudium et spes}).
\item La persona \textbf{no puede ser instrumentalizada} para fines ajenos a su propio desarrollo (económicos, sociales o políticos).
\item Los cambios sociales auténticos dependen de la \textbf{conversión interior} y de las capacidades morales de las personas.
\end{itemize}

\medskip\noindent \textbf{C. La libertad de la persona (\S{}135--143)}\par\nopagebreak
\begin{itemize}
\item \textbf{Valor y límites:} la libertad es signo eminente de la imagen divina; Dios quiso ``dejar al hombre en manos de su propia decisión''. No es arbitrio ni autonomía ilimitada: ``debe detenerse ante el árbol de la ciencia del bien y del mal''. Su ejercicio depende de condiciones económicas, sociales, jurídicas, políticas y culturales.
\item \textbf{Vínculo con la verdad y la ley natural:} la libertad se ejerce siguiendo el juicio de la conciencia, que no crea la verdad sino que la reconoce. La \textbf{ley natural} es universal e inmutable, participación de la ley eterna en la criatura racional; une a los hombres por encima de las culturas y es fundamento de derechos y deberes y de la \textbf{ley civil}.
\item La libertad está herida por el pecado y necesita ser liberada por Cristo.
\end{itemize}

\medskip\noindent \textbf{D. La igual dignidad de todas las personas (\S{}144--148)}\par\nopagebreak
\begin{itemize}
\item ``Dios no hace acepción de personas'': la dignidad de todo hombre ante Dios es el fundamento de la \textbf{igualdad y fraternidad}, también entre pueblos y Estados.
\item \textbf{Masculino y femenino:} diferencia en igual dignidad; complementariedad recíproca. La mujer es complemento del varón como el varón lo es de la mujer.
\item \textbf{Personas con discapacidad} (el texto dice ``minusválidas''): sujetos plenamente humanos, titulares de derechos y deberes; hay que promover su participación en la vida familiar y social.
\end{itemize}

\medskip\noindent \textbf{E. La sociabilidad humana (\S{}149--151)}\par\nopagebreak
\begin{itemize}
\item La persona es \textbf{constitutivamente un ser social}: por naturaleza, no por contrato.
\item La sociabilidad no lleva automáticamente a la comunión: está amenazada por la soberbia y el egoísmo.
\item Pluralismo social: el bien común depende de un sano pluralismo. La \textbf{familia, la comunidad civil y la comunidad religiosa} responden más inmediatamente a la naturaleza humana; otras sociedades nacen de la libre voluntad (asociaciones intermedias, ``socialización'').
\end{itemize}

\subsubsection{IV. Los derechos humanos (\S{}152--159)}
\begin{itemize}
\item \textbf{Valor:} la Declaración Universal de 1948 es ``una piedra miliar en el camino del progreso moral de la humanidad'' (Juan Pablo II).
\item \textbf{Raíz y fuente:} la raíz es la dignidad de todo ser humano, que se descubre ante todo \textbf{con la razón} y se confirma a la luz de la fe. La fuente última \textbf{no es la voluntad de los hombres, ni el Estado, ni los poderes públicos}, sino el hombre mismo y Dios su Creador.
\item \textbf{Características:} \textbf{universales} (todos, sin excepción), \textbf{inviolables} (inherentes a la persona) e \textbf{inalienables} (nadie puede privar de ellos). Deben tutelarse \textbf{en su conjunto} (indivisibles).
\item \textbf{Especificación} (\textit{Centesimus annus} 47): derecho a la vida desde la concepción; a vivir en una familia unida; a desarrollar la inteligencia y la libertad en la búsqueda de la verdad; al trabajo; a fundar una familia; y la \textbf{libertad religiosa}, ``fuente y síntesis'' de los demás. El primero es el \textbf{derecho a la vida}, desde la concepción hasta la muerte natural.
\item \textbf{Derechos y deberes} son inseparables: quienes reivindican sus derechos y olvidan sus deberes ``se asemejan a los que derriban con una mano lo que con la otra construyen'' (Juan XXIII, \textit{Pacem in terris}).
\item \textbf{Derechos de los pueblos y de las Naciones:} ``lo que es verdad para el hombre lo es también para los pueblos'': autodeterminación e independencia, existencia, lengua y cultura propias, modelar su vida según sus tradiciones (sin violar derechos humanos ni oprimir minorías).
\item \textbf{Distancia entre la letra y el espíritu:} contradicción entre la proclamación solemne y la realidad (guerras, genocidios, nuevas esclavitudes como la trata, los niños soldado, la explotación laboral, la droga y la prostitución). La Iglesia se compromete a \textbf{defender y promover} los derechos en doble dirección: \textbf{anuncio} de su fundamento cristiano y \textbf{denuncia} de sus violaciones; ``el anuncio es siempre más importante que la denuncia''. Advertencia: una afirmación excesiva de igualdad puede dar lugar a un individualismo en el que cada uno reivindica derechos sin hacerse responsable del bien común.
\end{itemize}

\subsubsection{Ideas clave para retener (Compendio, cap. III)}
\begin{enumerate}
\item \textbf{Principio personalista:} la persona es sujeto, fundamento y fin de la vida social; su dignidad es inviolable.
\item Fundamento: el hombre es \textbf{imagen de Dios}; unidad de alma y cuerpo (ni espiritualismo ni materialismo).
\item La persona no se reduce a su conciencia, inteligencia o utilidad, y no puede ser instrumentalizada.
\item El orden social se subordina a la persona, \textbf{no al contrario}.
\item Libertad ligada a la \textbf{verdad y a la ley natural}; no es autonomía sin límites.
\item Pecado personal y social; \textbf{estructuras de pecado} (afán de ganancia y sed de poder).
\item Derechos humanos: universales, inviolables, inalienables; su fuente no es el Estado sino el hombre y Dios; inseparables de los deberes.
\end{enumerate}



\subsection{4.2 Dicasterio para la Doctrina de la Fe --- \textit{Dignitas infinita} (2024)}

\noindent \textbf{Datos:} Declaración ``sobre la dignidad humana'' del Dicasterio para la Doctrina de la Fe (prefecto: card. Víctor M. Fernández), aprobada por Francisco el 25 de marzo de 2024, fechada el 2 de abril y publicada el 8 de abril, tras cinco años de elaboración. Las tres primeras partes recuerdan principios; la cuarta, ampliada por pedido expreso del Papa, denuncia violaciones actuales. Llama a \textit{Fratelli tutti} una ``Carta Magna'' de las tareas actuales para salvaguardar la dignidad. Conmemora los \textbf{75 años de la Declaración Universal de los Derechos Humanos} (1948). El título viene de una frase de Juan Pablo II (1980). Es además el \textbf{eje pedagógico transversal} de la cátedra.\par\smallskip

\subsubsection{Introducción (\S{}1--9)}
\begin{itemize}
\item \textbf{Tesis (\S{}1):} toda persona posee una \textbf{dignidad infinita}, fundada inalienablemente en su propio ser, ``más allá de toda circunstancia y en cualquier estado o situación en que se encuentre''. Es plenamente reconocible por la \textbf{sola razón} y la Revelación la confirma.
\item \textbf{Cuádruple distinción del concepto de dignidad (\S{}7--8):}
\end{itemize}

{\small\renewcommand{\arraystretch}{1.25}
\begin{longtable}{>{\raggedright\arraybackslash}p{0.074\linewidth}>{\raggedright\arraybackslash}p{0.418\linewidth}>{\raggedright\arraybackslash}p{0.418\linewidth}}
\toprule
\textbf{Tipo} & \textbf{En qué consiste} & \textbf{¿Se puede perder?} \\ \midrule \endhead
\textbf{Ontológica} & Corresponde a la persona por el solo hecho de existir y de ser querida, creada y amada por Dios & \textbf{Nunca}: es indeleble, vale en cualquier circunstancia \\ \addlinespace[2pt]
\textbf{Moral} & Se refiere al ejercicio de la libertad: obrar conforme o contra la conciencia & Sí: quien obra el mal se comporta de modo ``indigno'' (pero conserva la ontológica) \\ \addlinespace[2pt]
\textbf{Social} & Condiciones en las que vive una persona (por ejemplo, pobreza extrema) & Sí: hay condiciones de vida ``indignas'' \\ \addlinespace[2pt]
\textbf{Existencial} & Vivencia de la propia vida (enfermedad grave, violencia familiar, adicciones, soledad) & Sí: se puede percibir la propia vida como ``indigna'' \\ \addlinespace[2pt]
\bottomrule\end{longtable}}

\begin{itemize}
\item La más importante es la \textbf{ontológica}; las otras tres pueden faltar sin que ella desaparezca.
\item Definición clásica de persona (Boecio): ``sustancia individual de naturaleza racional''. La dignidad es \textbf{intrínseca}, no concedida por otros ni dependiente de dotes o cualidades.
\end{itemize}

\subsubsection{1. Una conciencia progresiva de la centralidad de la dignidad humana (\S{}10--16)}
\begin{itemize}
\item \textbf{Antig\''uedad clásica:} primera intuición, pero desde una perspectiva social: dignidad según el \textbf{rango} de cada uno; todavía lejos de fundar el respeto a toda persona.
\item \textbf{Perspectiva bíblica:} el hombre creado \textbf{a imagen y semejanza de Dios}; los profetas denuncian la injusticia con los pobres; Jesús se identifica con los últimos (``a mí me lo hicisteis'', Mt 25) y rompe barreras culturales y cultuales.
\item \textbf{Pensamiento cristiano:} Padres de la Iglesia; Tomás de Aquino (``persona significa lo que en toda naturaleza es perfectísimo''); humanismo cristiano del Renacimiento; ecos en Descartes y Kant; \textbf{personalismo} del siglo XX (intersubjetividad).
\item La dignidad es \textbf{conferida por Dios, no pretendida ni merecida}; no la conceden otros seres humanos, por eso no puede ser retirada.
\item \textbf{Hoy:} la Declaración de 1948; el Vaticano II (\textit{Dignitatis humanae}, \textit{Gaudium et spes}).
\end{itemize}

\subsubsection{2. La Iglesia anuncia, promueve y se hace garante de la dignidad humana (\S{}17--22)}
\noindent Triple convicción de fe:\par\smallskip
\begin{enumerate}
\item \textbf{Creación:} imagen \textbf{indeleble} de Dios en todo ser humano, que abarca a la persona entera (alma y cuerpo, varón y mujer).
\item \textbf{Encarnación:} Cristo \textbf{eleva} la dignidad del hombre; al unirse a todo hombre confirma que cada ser humano tiene una dignidad inestimable, por el solo hecho de pertenecer a la comunidad humana. Novedad de Jesús: el ser humano es \textbf{más digno de respeto y amor cuanto más débil, miserable y sufriente}.
\item \textbf{Resurrección:} vocación a la \textbf{plenitud} de la dignidad, la comunión con Dios.
\end{enumerate}
\begin{itemize}
\item \textbf{Compromiso con la propia libertad:} todos están llamados a manifestar esa dignidad en su vida; el pecado puede herirla y oscurecerla, pero \textbf{nunca borrarla}.
\end{itemize}

\subsubsection{3. La dignidad, fundamento de los derechos y de los deberes humanos (\S{}23--32)}
\begin{itemize}
\item \textbf{Respeto incondicional:} la dignidad no depende de la capacidad de razonar o de la autonomía; si así fuera, no tendrían dignidad el niño no nacido, el anciano dependiente o la persona con discapacidad mental. No se debe sustituir ``dignidad humana'' por ``dignidad personal'' entendida como capacidad de razonar.
\item \textbf{Crítica a la multiplicación arbitraria de derechos:} la dignidad no puede usarse para justificar cualquier preferencia individual o deseo subjetivo.
\item \textbf{Fundamento objetivo de la libertad:} la dignidad no se funda en el individualismo ni en la autorreferencia; existe una naturaleza humana común.
\item \textbf{Estructura relacional de la persona:} contra una libertad autorreferencial e individualista; la dignidad incluye la capacidad de asumir \textbf{obligaciones hacia los otros}. También respeto por las demás criaturas: solo es sostenible un \textbf{``antropocentrismo situado''}; el cuidado del ambiente pertenece a la dignidad del hombre.
\item El \textbf{relativismo moral} no es clave de convivencia pacífica sino origen de división y de negación de la dignidad (Benedicto XVI).
\item \textbf{Liberación de condicionamientos:} la libertad real exige condiciones económicas, sociales, jurídicas, políticas y culturales; Cristo libera la libertad.
\end{itemize}

\subsubsection{4. Algunas violaciones graves de la dignidad humana (\S{}33--62)}
\noindent Antes de la lista recuerda (\S{}34) que \textit{Gaudium et spes} 27 ya condenaba homicidio, genocidio, aborto, eutanasia, suicidio, mutilaciones, tortura, condiciones infrahumanas de vida, encarcelamientos arbitrarios, esclavitud, prostitución, trata y condiciones laborales degradantes; y agrega que la \textbf{pena de muerte} ``viola la dignidad inalienable de toda persona humana más allá de cualquier circunstancia''.\par\smallskip

{\small\renewcommand{\arraystretch}{1.25}
\begin{longtable}{>{\raggedright\arraybackslash}p{0.375\linewidth}>{\raggedright\arraybackslash}p{0.555\linewidth}}
\toprule
\textbf{Violación} & \textbf{Idea clave} \\ \midrule \endhead
\textbf{El drama de la pobreza} & Una de las mayores injusticias; desigual distribución de la riqueza. Es ilusorio distinguir sin más entre países ricos y pobres. ``No existe peor pobreza que aquella que priva del trabajo y de la dignidad del trabajo''. Todos somos responsables \\ \addlinespace[2pt]
\textbf{La guerra} & Tragedia que niega la dignidad humana; siempre es una ``derrota de la humanidad''. Hoy es muy difícil sostener la ``guerra justa''. La guerra en nombre de la religión es una guerra contra la religión \\ \addlinespace[2pt]
\textbf{El trabajo de los emigrantes} & Primeras víctimas de las múltiples formas de pobreza; se les niega dignidad en origen y en destino \\ \addlinespace[2pt]
\textbf{La trata de personas} & ``Actividad innoble, una verg\''uenza para nuestras sociedades''; es un \textbf{crimen contra la humanidad}: desfigura a la víctima y deshumaniza a quien la comete. Incluye comercio de órganos, explotación sexual, trabajo esclavo \\ \addlinespace[2pt]
\textbf{Los abusos sexuales} & Dejan profundas cicatrices; también dentro de la Iglesia \\ \addlinespace[2pt]
\textbf{Las violencias contra las mujeres} & Desigualdades, feminicidio, aborto forzado, poligamia \\ \addlinespace[2pt]
\textbf{El aborto} & La dignidad de la vida vale desde la concepción; denuncia el lenguaje ambiguo (``interrupción del embarazo'') \\ \addlinespace[2pt]
\textbf{La maternidad subrogada} & El niño se convierte en ``mero objeto''; viola la dignidad del niño y de la mujer. ``Un hijo es siempre un don y nunca el objeto de un contrato''; no existe un ``derecho al hijo''. Pide su prohibición universal \\ \addlinespace[2pt]
\textbf{La eutanasia y el suicidio asistido} & Usa un concepto erróneo de dignidad contra la vida misma (``leyes de muerte digna''). El sufrimiento no hace perder la dignidad. Sí a los cuidados paliativos, no al encarnizamiento terapéutico. ``La vida es un derecho, no la muerte'' \\ \addlinespace[2pt]
\textbf{El descarte de las personas con discapacidad} & ``Cultura del descarte''; se pide inclusión y participación activa \\ \addlinespace[2pt]
\textbf{La teoría de género} & Primero: toda persona, cualquiera sea su tendencia sexual, debe ser respetada; es contrario a la dignidad encarcelar, torturar o matar por orientación sexual. Luego critica la teoría de género como ``colonización ideológica'': pretende disponer de sí sin reconocer la vida como don y negar la diferencia sexual. Sexo biológico y género ``se pueden distinguir pero no separar'' \\ \addlinespace[2pt]
\textbf{El cambio de sexo} & El cuerpo participa de la dignidad de la imagen de Dios; por regla general \textbf{corre el riesgo de atentar} contra la dignidad recibida desde la concepción. No incluye la asistencia médica por anomalías genitales \\ \addlinespace[2pt]
\textbf{La violencia digital} & Nuevas formas de violencia en redes: noticias falsas, calumnia, ciberacoso, pornografía, explotación, adicción, aislamiento. Las tecnologías deben servir a la dignidad \\ \addlinespace[2pt]
\bottomrule\end{longtable}}

\subsubsection{Conclusión (\S{}63--66)}
\begin{itemize}
\item El respeto de la dignidad de toda persona, \textbf{en cualquier circunstancia}, debe estar en el centro del compromiso por el bien común y de todo ordenamiento jurídico.
\item Corresponde también a cada persona y comunidad la realización concreta de la dignidad, y a los Estados protegerla y garantizar las condiciones para que florezca.
\end{itemize}

\subsubsection{Ideas clave para retener (Dignitas infinita)}
\begin{enumerate}
\item Dignidad \textbf{infinita, inalienable, intrínseca}: no depende de circunstancias, capacidades ni conducta.
\item Cuatro sentidos: \textbf{ontológica} (nunca se pierde), \textbf{moral, social, existencial}.
\item Tres fundamentos de fe: creación a imagen de Dios, \textbf{encarnación} (que confirma y eleva la dignidad de todo ser humano), vocación a la comunión con Dios.
\item La dignidad es fundamento de derechos \textbf{y deberes}; no justifica cualquier deseo individual.
\item Trece violaciones graves (pobreza, guerra, emigrantes, trata, abusos, violencia contra las mujeres, aborto, maternidad subrogada, eutanasia, descarte de discapacitados, teoría de género, cambio de sexo, violencia digital), más el rechazo de la \textbf{pena de muerte}.
\end{enumerate}



\subsection{4.3 Francisco --- \textit{Laudato Si'} (2015), Cap. III: ``Raíz humana de la crisis ecológica''}

\noindent \textbf{Idea central:} la crisis ecológica tiene una raíz humana: un modo de entender la vida y la acción que se ha desviado. Dos nombres para esa raíz: el \textbf{paradigma tecnocrático} dominante y el \textbf{antropocentrismo desviado}.\par\smallskip

\subsubsection{I. La tecnología: creatividad y poder (\S{}102--105)}
\begin{itemize}
\item Valoración positiva: la tecnociencia bien orientada es ``producto maravilloso de la creatividad humana'', remedió males innumerables y puede producir cosas valiosas y bellas.
\item Pero otorga un \textbf{tremendo poder} a quienes tienen el conocimiento y sobre todo el poder económico para usarlo. Nada garantiza que se use bien.
\item ``El inmenso crecimiento tecnológico \textbf{no estuvo acompañado de un desarrollo del ser humano en responsabilidad, valores, conciencia}''. Falta una ética sólida, una cultura y una espiritualidad que pongan límites.
\end{itemize}

\subsubsection{II. Globalización del paradigma tecnocrático (\S{}106--114)}
\begin{itemize}
\item El problema de fondo es que la humanidad asumió la tecnología junto con un \textbf{paradigma homogéneo y unidimensional}: un sujeto que, mediante el método científico, \textbf{posee y domina} un objeto exterior, tratado como algo informe y totalmente disponible para su manipulación.
\item De ahí la idea de un \textbf{crecimiento infinito o ilimitado}, que supone la ``mentira'' de la disponibilidad infinita de los bienes del planeta.
\item Los objetos técnicos \textbf{no son neutros}: crean un entramado que \textbf{condiciona los estilos de vida} y orienta las posibilidades sociales \textbf{en la línea de los intereses de determinados grupos de poder}.
\item El paradigma tiende a dominar también \textbf{la economía y la política}: la economía asume todo desarrollo tecnológico en función del rédito; las finanzas ahogan a la economía real. El mercado por sí mismo no garantiza el desarrollo humano integral ni la inclusión social.
\item La \textbf{fragmentación de los saberes} (especialización) impide ver el conjunto; se pierde el sentido de la totalidad.
\item \textbf{Las soluciones meramente técnicas} que aparecen ante cada problema ambiental aíslan cosas que están conectadas y \textbf{esconden los verdaderos y más profundos problemas} del sistema mundial.
\item La libertad humana \textbf{es capaz de limitar la técnica}, orientarla y ponerla al servicio de un progreso más sano, humano, social e integral.
\item Propuesta: una mirada distinta, un pensamiento, una política, un programa educativo, un estilo de vida y una espiritualidad que resistan al paradigma. Una \textbf{``valiente revolución cultural''}: no volver a las cavernas, pero sí aminorar la marcha y recuperar valores.
\end{itemize}

\subsubsection{III. Crisis y consecuencias del antropocentrismo moderno (\S{}115--136)}
\begin{itemize}
\item El \textbf{antropocentrismo moderno} colocó la razón técnica por encima de la realidad: la naturaleza ya no es norma ni refugio, sino mero espacio y materia para una obra.
\item La forma correcta de entender al hombre como ``señor'' del universo es la de \textbf{administrador responsable} (no el ``sueño prometeico'' de dominio).
\item Un \textbf{antropocentrismo desviado} da lugar a un estilo de vida desviado. Pero la respuesta no es un ``biocentrismo'': \textbf{``no hay ecología sin una adecuada antropología''}. No se puede exigir al hombre un compromiso con el mundo si no se reconocen sus capacidades peculiares de conocimiento, voluntad, libertad y responsabilidad.
\item Todo está conectado: no es compatible defender la naturaleza y justificar el aborto, ni cuidar otros seres sin cuidar a los humanos más frágiles (embrión, pobres, personas con discapacidad).
\end{itemize}

\medskip\noindent \textbf{El relativismo práctico (\S{}122--123)}\par\nopagebreak
\begin{itemize}
\item Cuando el ser humano se coloca en el centro, da prioridad absoluta a sus conveniencias circunstanciales y \textbf{todo lo demás se vuelve relativo}.
\item Es la misma lógica que lleva a \textbf{aprovecharse de otro y tratarlo como mero objeto}: trabajos forzados, esclavitud por deudas, explotación sexual de niños, abandono de ancianos, trata, narcotráfico, comercio de órganos y de pieles de animales en extinción. Es la \textbf{lógica del ``usa y tira''}, del descarte.
\item Sin verdades objetivas ni principios sólidos, las leyes se ven como imposiciones arbitrarias.
\end{itemize}

\medskip\noindent \textbf{Necesidad de preservar el trabajo (\S{}124--129)}\par\nopagebreak
\begin{itemize}
\item En cualquier planteo de ecología integral es indispensable incorporar \textbf{el valor del trabajo}.
\item El trabajo es una \textbf{necesidad, parte del sentido de la vida} en esta tierra, camino de maduración, de desarrollo humano y de realización personal (ejemplo: el monacato benedictino, ``ora et labora'').
\item \textbf{No debe buscarse que el progreso tecnológico reemplace cada vez más el trabajo humano}: con eso la humanidad se dañaría a sí misma.
\item Ayudar a los pobres con dinero debe ser siempre una \textbf{solución provisoria}; el gran objetivo es permitirles una vida digna a través del trabajo.
\item Para que haya trabajo hay que promover una economía que favorezca la \textbf{diversidad productiva y la creatividad empresarial} (pequeños productores).
\end{itemize}

\medskip\noindent \textbf{Innovación biológica a partir de la investigación (\S{}130--136)}\par\nopagebreak
\begin{itemize}
\item La experimentación con animales es legítima solo dentro de límites razonables y para cuidar o salvar vidas humanas; es contrario a la dignidad humana hacerlos sufrir inútilmente.
\item \textbf{Organismos genéticamente modificados (OGM):} \textbf{es difícil emitir un juicio general}; pueden ser muy diversos y requieren distintas consideraciones. Los riesgos no siempre se deben a la técnica sino a su aplicación inadecuada o excesiva.
\begin{itemize}
\item Beneficios: en algunas regiones produjeron crecimiento económico.
\item Dificultades: \textbf{concentración de tierras} en pocas manos, desaparición de pequeños productores, precarización, pérdida de biodiversidad, \textbf{oligopolios} de semillas (más aún con semillas estériles).
\item Hace falta una \textbf{discusión científica y social responsable y amplia}, con \textbf{información completa} y con todos los afectados (agricultores, consumidores, autoridades, científicos), más investigación \textbf{autónoma e interdisciplinaria}.
\end{itemize}
\item Incoherencia: algunos ecologistas defienden la integridad del ambiente pero no aplican los mismos principios a la \textbf{vida humana} (experimentación con embriones vivos). \textbf{La técnica separada de la ética difícilmente podrá autolimitar su poder.}
\end{itemize}

\subsubsection{Ideas clave para retener (Laudato Si', cap. III)}
\begin{enumerate}
\item Raíz de la crisis: \textbf{paradigma tecnocrático} (dominio sobre un objeto disponible, crecimiento ilimitado) y \textbf{antropocentrismo desviado}.
\item La técnica \textbf{no es neutra}: condiciona estilos de vida según los intereses de grupos de poder; domina economía y política.
\item Las soluciones solo técnicas \textbf{ocultan los verdaderos problemas}.
\item \textbf{Relativismo práctico} = lógica del ``usa y tira'': justifica explotación y descarte.
\item El \textbf{trabajo} es parte del sentido de la vida; la técnica no debe reemplazarlo sin más.
\item \textbf{OGM:} ni rechazo total ni aceptación incondicional; discernimiento caso por caso, debate amplio, información completa.
\item ``No hay ecología sin una adecuada antropología.''
\end{enumerate}



\subsection{4.4 Sorge --- Cap. XXI y XXIII}

\subsubsection{CAP. XXI --- Biotecnologías: ¿se puede manipular la vida?}

\medskip\noindent \textbf{Planteo}\par\nopagebreak
\begin{itemize}
\item Las biotecnologías son técnicas que emplean organismos vivientes e intervienen en sus procesos biológicos para obtener productos (medicina, agricultura, zootecnia, alimentación, química, ambiente).
\item El sector más revolucionario es la ingeniería genética aplicada al hombre, pero el capítulo se limita a \textbf{animales y plantas}: con el ADN recombinante se obtienen \textbf{organismos genéticamente modificados (OGM)} (``supertopo'', ``oveja-cabra'', oveja Dolly, plantas transgénicas resistentes).
\item Opinión dividida: unos las rechazan en bloque como manipulación inaceptable; otros subrayan solo las ventajas (más producción, beneficio para países pobres, menos pesticidas).
\item \textbf{Nadie puede hoy valorar con exactitud los riesgos}; por eso resulta inaplicable el \textbf{principio de precaución} (Convención sobre biodiversidad, 1992: un producto solo es legítimo si hay certeza moral de que no conlleva riesgos graves e irreversibles). Aun así es \textbf{obligado} hacer todo lo posible por preverlos y prevenirlos.
\item Preguntas ineludibles: efectos sobre la salud, sobre los equilibrios ecológicos, y si las ventajas se repartirán equitativamente o serán monopolio de los países ricos y las multinacionales.
\item La verdadera pregunta no es técnica sino sobre el significado de la vida y de la naturaleza.
\end{itemize}

\medskip\noindent \textbf{1. La dimensión ética de las biotecnologías}\par\nopagebreak
\begin{itemize}
\item El verdadero problema es \textbf{cultural}: se refiere a la concepción misma de la vida. La era tecnológica acelera una \textbf{transmutación de valores}: la eficiencia como criterio fundamental; el bienestar identificado con el consumismo; concepción inmanente de la persona; individualismo, subjetivismo y relativismo ético vistos como ``madurez civil''.
\item El hombre \textbf{no puede reducirse al determinismo de sus mecanismos biológicos}: es pensamiento, espíritu, libertad, sentido de lo divino. ``Cada vez que el ser pierde la primacía sobre el obrar, el hombre pierde la libertad y se destruye''.
\item Sin embargo, \textbf{sería insensato condenar las biotecnologías en cuanto tales}: abren perspectivas de humanización (prevenir malformaciones, curar enfermedades). Las desviaciones no dependen de la tecnología en sí, sino del \textbf{uso} que el hombre hace de ella.
\item El progreso científico y técnico \textbf{nunca es un fin en sí mismo}: está esencialmente orientado al hombre. Por eso las biotecnologías tienen una dimensión \textbf{intrínsecamente ética}, que es la que les da pleno significado.
\item No se resta importancia a la \textbf{autonomía} de la ciencia: es búsqueda de la verdad, y el conocimiento es siempre un bien, un factor de liberación. \textbf{No hay que temer el progreso, sino favorecerlo}; preocupa su uso distorsionado.
\item Hablar solo de ``límites éticos'' es \textbf{reductor}: no se trata de impedir su uso, sino de hacer realidad su pleno significado.
\item \textbf{Del mero hecho de que algo sea técnicamente posible no se sigue que sea lícito hacerlo.} El juicio moral pertenece a un orden distinto del científico, pero no separado: solo la responsabilidad moral hace humano el progreso.
\item Las biotecnologías \textbf{no pueden manipular la vida a su antojo, sino únicamente ayudarla a desarrollarse y crecer}.
\end{itemize}

\medskip\noindent \textbf{2. Frente a los desafíos de las biotecnologías}\par\nopagebreak
\begin{itemize}
\item \textbf{Dos posturas equivocadas:} (a) quien, desde una visión sagrada e intangible de la naturaleza, niega \textit{a priori} toda modificación; (b) quien, en nombre de la primacía de la ciencia, ni se plantea el problema moral.
\item \textbf{Postura coherente:} preguntarse responsablemente por los peligros y asumir el deber moral de garantizar la seguridad de la vida humana y respetar toda forma de vida, evitando ``imprudencias, destrucciones y crueldades inútiles''.
\item \textbf{Criterio moral:} el recurso a las biotecnologías se justifica solo si está garantizada la \textbf{tutela sustancial del ecosistema} y con ventajas proporcionadas (reducir contaminación y consumo de recursos no renovables, luchar contra enfermedades incurables).
\item Derecho de los ciudadanos a una \textbf{información transparente} (etiquetado claro de los alimentos transgénicos) y derecho de los \textbf{países pobres} a acceder a las biotecnologías.
\item \textbf{Iniciativas necesarias:}
\begin{enumerate}
\item Afrontar el problema desde una \textbf{óptica global}, superando el positivismo de la primacía absoluta de la ciencia.
\item \textbf{Formación}: investigadores con competencia científica y además ética, jurídica y política; sensibilizar a la gente.
\item Iniciativas concretas, sobre todo de los cristianos, para garantizar la \textbf{primacía del ser sobre el hacer} y del bien común sobre los intereses particulares: una \textbf{cultura de la vida}.
\item Un \textbf{vasto consenso popular} en la sociedad democrática, también para impedir la especulación de las multinacionales.
\end{enumerate}
\item Contra las patentes: \textbf{``La vida no es un invento industrial. La vida no se fabrica. La vida no se patenta.''} Reconocer a unos pocos el derecho a disponer de las biotecnologías sería legitimar nuevas formas de monopolio y colonialismo.
\item Conclusión: solo integrando la investigación en un \textbf{humanismo integral, abierto a Dios}, las biotecnologías alcanzan su finalidad: servir a una calidad de vida cada vez más digna del hombre.
\end{itemize}

\subsubsection{CAP. XXIII --- La dimensión ética de la investigación científica}

\noindent \textbf{Planteo:} los descubrimientos (clonación de animales superiores, madre subrogada, cromosoma artificial) reavivan la \textit{vexata quaestio} de las fronteras de la ciencia y su relación con la ética. Se acusa a la Iglesia de rechazar la ciencia. Hay que distinguir dos cuestiones: la autonomía de la investigación y su dimensión ética.\par\smallskip

\medskip\noindent \textbf{1. Autonomía de la investigación científica y técnica}\par\nopagebreak
\begin{itemize}
\item Nadie la pone en duda, \textbf{tampoco la Iglesia}. El Vaticano I (1870): fe y razón no pueden discordar. El Vaticano II (\textit{Gaudium et spes} 36): la ciencia y la técnica tienen \textbf{leyes, finalidades y métodos propios}, inscritos en la naturaleza de las cosas, que no dependen del conocimiento religioso.
\item Razón última: ciencia y técnica son un hecho de \textbf{conocimiento humano}; la búsqueda de la verdad es siempre un bien en sí y es humanizadora.
\item No puede haber oposición entre lo verdadero conocido por la ciencia y lo verdadero conocido por la fe: ambos tienen origen en el mismo Dios.
\item La Iglesia \textbf{no solo no teme sino que favorece positivamente} el progreso científico. Dice \textbf{``sí'' a la ciencia y ``no'' al cientifismo}.
\item \textbf{Cientifismo:} deformación del concepto de ciencia con una \textbf{doble pretensión absolutista}: (1) ser la \textbf{única forma auténtica de conocimiento} de toda realidad, incluido el hombre; (2) \textbf{no estar sometida a ningún vínculo ético}, ni en los fines ni en los medios, ni en la investigación ni en la aplicación.
\end{itemize}

\medskip\noindent \textbf{2. Dimensión ética de la investigación}\par\nopagebreak
\begin{itemize}
\item La investigación, como toda búsqueda de la verdad, está \textbf{intrínsecamente orientada al hombre} y a su crecimiento integral (material y espiritual, individual y social, inmanente y trascendente).
\item Por eso su autonomía \textbf{no es absoluta sino relativa}: se armoniza con las otras formas de conocimiento.
\item El cientifismo considera la ciencia como fin en sí misma y trata al hombre como \textbf{epifenómeno} de mecanismos físico-biológico-sociales. Cuando se desconoce la dignidad trascendente de la persona, el hombre puede ser \textbf{manipulado, instrumentalizado, degradado a cobaya} (manipulación genética indiscriminada, producción in vitro de embriones, comercio de órganos, procreación artificial desaprensiva).
\item Hablar de ``fronteras'' o límites éticos \textbf{no es anticientífico}: es reafirmar la naturaleza propia de la investigación. La ciencia abre posibilidades, pero \textbf{no le corresponde decir qué uso debe hacerse de ellas}; no puede sustituir a la fe, los valores morales, el arte o la política (Juan Pablo II).
\end{itemize}

\medskip\noindent \textbf{3. Ciencia y ética, un encuentro necesario}\par\nopagebreak
\begin{itemize}
\item El futuro no está en una \textbf{dicotomía antinatural} ni en una \textbf{indebida interferencia mutua}, sino en la \textbf{síntesis armónica}, respetando la recíproca autonomía.
\item Si la ciencia se sustrae a todo condicionamiento ético puede volverse contra el hombre, sometiéndolo a \textbf{nuevas formas de esclavitud y totalitarismo}.
\item Se repite el principio: \textbf{que algo pueda hacerse científica o técnicamente no significa que sea lícito hacerlo}.
\item Conclusión: la autonomía de la ciencia no tiene nada que perder y sí todo que ganar con el reconocimiento de su dimensión ética; y la ética no tiene nada que temer del progreso. La Iglesia impulsa a los cristianos a mostrar la fecundidad del encuentro entre \textbf{inteligencia y conciencia, ciencia y fe}.
\end{itemize}

\subsubsection{Ideas clave para retener (Sorge XXI y XXIII)}
\begin{enumerate}
\item Ni condena en bloque ni aceptación acrítica: las biotecnologías no son malas en sí; el problema es su \textbf{uso} y la cultura que las orienta.
\item La dimensión ética es \textbf{intrínseca} a la ciencia y la técnica, porque están orientadas al hombre.
\item \textbf{Técnicamente posible no equivale a moralmente lícito.}
\item Las biotecnologías no pueden manipular la vida a su antojo, solo ayudarla a desarrollarse. ``La vida no se patenta.''
\item Autonomía de la ciencia: legítima pero \textbf{relativa}, no absoluta.
\item \textbf{Sí a la ciencia, no al cientifismo} (doble pretensión: único conocimiento válido y sin vínculo ético).
\item Síntesis armónica entre ciencia y ética: ni dicotomía ni interferencia.
\end{enumerate}



\subsection{4.5 Faggioni --- ``Transhumanismo. Volar más allá de la naturaleza humana''}

\noindent \textit{(Texto de la carpeta de bibliografía, para el tema 2 del eje: posthumanismo, transhumanismo y humanismo integral.)}\par\smallskip

\subsubsection{1. Definición y raíces culturales}
\begin{itemize}
\item Término acuñado por el biólogo \textbf{Julian Huxley} (1927). \textbf{Nick Bostrom:} movimiento intelectual y cultural que afirma la posibilidad y el deseo de \textbf{mejorar fundamentalmente la condición humana} mediante la razón aplicada, sobre todo con tecnologías para \textbf{eliminar el envejecimiento} y \textbf{potenciar} las capacidades intelectuales, físicas y psicológicas.
\item Ideas recurrentes: \textbf{mejoramiento} (\textit{improvement}), \textbf{potenciación} (\textit{enhancement}) y \textbf{superación} (\textit{overcoming}) de los límites humanos.
\item No niega la naturaleza humana: la \textbf{reduce a su expresión material} y proyecta transformarla hasta superarla.
\item Raíces: mitos de inmortalidad (Gilgamesh, Prometeo, Ícaro); el humanismo y la revolución científica (\textbf{Francis Bacon}: saber traducido en hacer); empirismo de Hume, materialismo de La Mettrie, dualismo cartesiano, evolucionismo de Darwin. Supone una \textbf{visión reduccionista del hombre}.
\item \textbf{No} depende del superhombre de Nietzsche (crecimiento interior de individuos excepcionales): el hombre nuevo transhumanista es una criatura potenciada por la biotecnología.
\item Antecedentes: Haldane (genética), Turing y Minsky (inteligencia artificial), criogenia (Ettinger, 1962), FM-2030, extropianismo (Max More), \textbf{World Transhumanist Association} (1998, Bostrom y Pearce).
\end{itemize}

\subsubsection{2. El rostro del transhumanismo}
\begin{itemize}
\item Nace de la \textbf{insatisfacción por la condición humana natural} (envejecimiento, enfermedad, límites). Es, en última instancia, \textbf{el rechazo del límite, de la finitud y de la vulnerabilidad}.
\item Medios: neurociencia, inteligencia artificial, nanotecnología, biotecnología genética; hibernación, transferencia del cerebro, ectogénesis, separación total entre sexualidad y procreación, fluidez de género.
\item Objetivo: crear \textbf{seres posthumanos}; los estadios intermedios son los ``transhumanos'' o ``humanos +''.
\item Paso de la técnica \textbf{convencional} (ayuda al hombre a sobrevivir dominando el mundo natural) a la técnica \textbf{innovadora} (reproduce, replasma y recrea la naturaleza, incluido el organismo humano). Del \textit{homo faber} al \textbf{\textit{homo creator}} (Anders). De una \textbf{antropología de lo incompleto} a una \textbf{antropología de la perfección}.
\item La Declaración transhumanista defiende el ``derecho moral'' a usar la tecnología para ampliar capacidades y llama \textbf{tecnofobia} a toda oposición.
\end{itemize}

\subsubsection{3. ¿Transhumano o inhumano? (crítica)}
\begin{itemize}
\item \textbf{Dante} usa ``transhumanar'' (Paraíso I): el hombre es elevado \textbf{por gracia} más allá de lo humano, en la línea del \textbf{cumplimiento} de su naturaleza y en la lógica del \textbf{don}. El transhumanismo lo entiende como \textbf{superación} de la naturaleza, en la lógica de la \textbf{manipulación y el dominio}.
\item Error central: \textbf{reducir el anhelo de autotrascendencia a la manipulación del organismo biopsíquico}. Se vuelve un proyecto \textbf{eugenésico} más, que pierde de vista el valor inalienable de cada ser humano.
\item \textit{Dignitas personae} (2008) n. 27: en esas propuestas hay ``insatisfacción o hasta rechazo del valor del ser humano como criatura y persona finita''; el hombre pretende \textbf{sustituir al Creador} (argumento de la \textit{hybris}).
\item La técnica es \textbf{``asimbólica''}: es puro medio, no trae en sí valores. Los criterios deben venir de una interpretación de la \textbf{naturaleza humana} como portadora de un proyecto (\textit{telos}). Sin ella queda solo el \textbf{deseo} y la \textbf{eficiencia}.
\item \textbf{Contradicción dramática:} el deseo se respeta por ser humano, pero se vuelve contra la naturaleza de la que surge: ``el hombre se afirmaría en cuanto libre negándose en cuanto humano''.
\item Críticos de distinta procedencia: \textbf{Fukuyama, Habermas} (la eugenesia positiva crea asimetría entre padres e hijos y debilita la igualdad democrática), \textbf{Kass, Pellegrino}, el informe \textit{Beyond Therapy}. Punto común: la naturaleza humana es la última defensa de la dignidad y de la democracia.
\item La \textbf{antropología personalista} evita el biologicismo: el hombre es una realidad multidimensional pero unitaria; por la unidad de alma y cuerpo, toda alteración del cuerpo repercute en lo espiritual.
\item Conclusión: el posthumano corre el riesgo de ser inhumano. \textbf{``El transhumanismo no es un humanismo.''}
\end{itemize}

\subsubsection{4. ¿El hombre es perfectible?}
\begin{itemize}
\item Son \textbf{lícitas} las intervenciones terapéuticas, incluida la terapia génica \textbf{somática} (\textit{Dignitas personae} 26); la germinal por ahora no es aceptable.
\item Las manipulaciones de \textbf{mejora} favorecen una mentalidad eugenésica, estigmatizan a quienes no poseen ciertas dotes y atentan contra la \textbf{igualdad} de todos los seres humanos; ¿quién decide qué es ``mejor''? Riesgo de crear \textbf{grupos humanos diferentes} y nuevas marginaciones (Juan Pablo II). Corresponde aplicar el \textbf{principio de precaución}.
\item Faggioni propone una \textbf{tercera vía} entre transhumanistas y bioconservadores: entre ``una ética paralizada por el temor y una ética ebria por la ciencia'', una \textbf{ética de la responsabilidad del hombre por el hombre}. Podrían admitirse mejoras que respeten la identidad, la igualdad y la integridad de la naturaleza humana, en continuidad con sus potencialidades: se trata de \textbf{mejorar las condiciones biológicas de la existencia personal, no de crear un hombre ``mejor''}.
\item Fundamento: el \textbf{señorío ministerial} (no absoluto) del hombre sobre su vida (\textit{Evangelium vitae} 52).
\end{itemize}

\subsubsection{Ideas clave para retener (transhumanismo)}
\begin{enumerate}
\item Transhumanismo = mejorar y superar la condición humana por la tecnología; meta: el \textbf{posthumano}.
\item Supone una antropología \textbf{reduccionista y materialista} y el \textbf{rechazo de la finitud}.
\item Crítica cristiana: la dignidad no depende de capacidades; el hombre es criatura finita, unidad de cuerpo y alma; la técnica no da por sí sola los fines.
\item Alternativa: \textbf{humanismo integral} (Compendio, \textit{Caritas in veritate}): desarrollo de todo el hombre y de todos los hombres, abierto a la trascendencia.
\item Terapia sí; mejora eugenésica no; ética de la responsabilidad.
\end{enumerate}

\subsubsection{Síntesis del Eje 4}
\begin{esquema}
\noindent\hspace*{0.0em}PERSONA (imagen de Dios, unidad cuerpo-alma, dignidad infinita)\par
\noindent\hspace*{1.7em}$\rightarrow$ fundamento de DERECHOS y DEBERES\par
\noindent\hspace*{1.7em}$\rightarrow$ criterio para juzgar la CIENCIA y la TÉCNICA\par
\noindent\hspace*{4.4em}- autonomía legítima pero relativa (sí ciencia, no cientifismo)\par
\noindent\hspace*{4.4em}- posible no es igual a lícito\par
\noindent\hspace*{4.4em}- paradigma tecnocrático / transhumanismo = reducción del hombre\par
\noindent\hspace*{1.7em}$\rightarrow$ propuesta: HUMANISMO INTEGRAL + ética de la responsabilidad\par
\end{esquema}



\clearpage\section{EJE N\textordmasculine{} 5 --- Aportes del papa Francisco al pensamiento social de la Iglesia}

\subsection{5.1 Camacho --- ``Diez años del papa Francisco. Aporte al pensamiento social de la Iglesia''}

\noindent \textbf{Propósito del autor:} no recorrer temas (propiedad, democracia, ecología), sino identificar \textbf{claves o ejes transversales} que muestran la coherencia del pensamiento de Francisco, inseparable de sus \textbf{acciones y gestos}.\par\smallskip

\subsubsection{A) Tomando perspectiva: más allá de los contenidos}
\begin{itemize}
\item La doctrina social no es atemporal e inmutable: es un \textbf{proceso reflexivo} en el que aparecen temas nuevos (el desarrollo a mediados del siglo XX, la ecología en el XXI).
\item \textbf{Siglo XIX:} la Iglesia vive en \textbf{confrontación con la Modernidad}, con postura \textbf{defensiva}; se distancia a la vez del \textbf{socialismo} (reafirma el orden vigente y la propiedad privada) y del \textbf{liberalismo} (critica la libertad absoluta y el Estado liberal).
\item \textbf{Vaticano II = punto de inflexión:} la Modernidad deja de ser incompatible con la fe; ``el hombre moderno es, por fin, aceptado como \textbf{interlocutor}'', con talante crítico pero sin descalificación.
\item Nueva eclesiología: de la Iglesia como \textbf{``sociedad perfecta''} jerárquica (definición jurídica) a la Iglesia como \textbf{misterio} y \textbf{pueblo de Dios}; ``sacramento'' de la unión con Dios y de la unidad del género humano (\textit{Lumen gentium}). Consecuencia: \textbf{revalorización del laicado}; la misión compete a todos.
\item Para la doctrina social: ya no relaciones ``entre cabezas'' (Estado e Iglesia), sino presencia de los laicos en las estructuras de la sociedad. Es un conjunto de principios \textbf{para la acción}.
\end{itemize}

\subsubsection{B) Dos documentos de referencia}
\begin{itemize}
\item Cinco grandes documentos: \textit{Evangelii gaudium} (2013), \textit{Laudato si'} (2015), \textit{Amoris laetitia} (2016), \textit{Gaudete et exsultate} (2018), \textit{Fratelli tutti} (2020).
\item \textbf{\textit{Evangelii gaudium}:} documento \textbf{programático} del pontificado; el que Francisco dejaría como legado (``un plagio de \textit{Evangelii nuntiandi} de Pablo VI y del Documento de Aparecida''). No es encíclica social sino sobre la evangelización. Tres aspectos: (1) \textbf{``Iglesia en salida''} misionera (``prefiero una Iglesia accidentada, herida y manchada por salir a la calle''); (2) la misión es de \textbf{todo el pueblo de Dios}; (3) la evangelización incluye la \textbf{dimensión social}: inclusión de los pobres, y diálogo y paz social.
\item \textbf{\textit{Fratelli tutti}:} sí es encíclica social; sintetiza sus preocupaciones y propone el horizonte de la \textbf{fraternidad universal}, involucrando a las religiones.
\end{itemize}

\subsubsection{C) Rasgos de su pensamiento social}
\begin{enumerate}
\item \textbf{Más inductivo:} parte del análisis de la realidad (contraste con Benedicto XVI, más \textbf{deductivo}: de la teología a la antropología y de ahí a la sociedad).
\item \textbf{Menos ``clásico'', menos doctrinal y más personal:} interpela al lector, estilo pastoral (el capítulo del buen samaritano como meditación); entrevistas, improvisaciones.
\item \textbf{Los gestos} importan tanto o más que las palabras y a menudo las preceden: la \textbf{sinodalidad}, los encuentros con \textbf{movimientos populares}, la \textbf{Economía de Francisco}.
\end{enumerate}

\subsubsection{D) Las seis claves}

\medskip\noindent \textbf{Primera clave: el tiempo es superior al espacio}\par\nopagebreak
\begin{itemize}
\item Es uno de los \textbf{cuatro principios} de \textit{Evangelii gaudium} para construir un pueblo: (1) \textbf{el tiempo es superior al espacio}; (2) \textbf{la unidad prevalece sobre el conflicto}; (3) \textbf{la realidad es más importante que la idea}; (4) \textbf{el todo es superior a la parte}. Inspiración última: \textbf{Romano Guardini} (las ``oposiciones'' polares).
\item Tensión entre \textbf{plenitud} (el tiempo: horizonte, utopía) y \textbf{límite} (el espacio: el momento, la coyuntura).
\item Significa \textbf{``iniciar procesos más que poseer espacios''}: trabajar a largo plazo sin obsesionarse por resultados inmediatos; paciencia ante las dificultades. Dar prioridad al espacio lleva a querer tener todo resuelto ya y a ocupar espacios de poder.
\item Los otros tres principios son desarrollos del primero. De la atención a la complejidad nace la imagen del \textbf{poliedro} frente a la \textbf{esfera} (lo homogéneo y acabado).
\end{itemize}

\medskip\noindent \textbf{Segunda clave: personalismo frente a individualismo; raíces ignacianas}\par\nopagebreak
\begin{itemize}
\item La crítica al \textbf{individualismo} es constante: es causa de muchos males; el \textbf{nacionalismo} es una de sus expresiones; es la \textbf{estrategia de los poderosos}, porque disgrega a los pueblos.
\item Alternativa: la \textbf{persona} como ser constitutivamente \textbf{social}, en relación (antropología cristiana, frente al individualismo liberal y al colectivismo marxista).
\item Matiz \textbf{ignaciano}: la persona como \textbf{sujeto único e irrepetible}, libre y autónomo, con un camino propio ante Dios (Ejercicios espirituales; conversión de Ignacio en Loyola y Manresa). Clave: el \textbf{discernimiento}.
\item Ejemplo: \textit{Amoris laetitia}: la moral no se reduce al cumplimiento de la norma; discernir cada caso, gradualidad, conciencia.
\item \textbf{Antropocentrismo desviado} (\textit{Laudato si'}): Francisco no critica el antropocentrismo sino su desviación: reducir todo lo exterior al sujeto, incluidas las otras personas, a \textbf{instrumento}. Su origen es la universalización del \textbf{paradigma tecnocrático}.
\item De ahí la crítica al sistema económico: \textbf{``economía de la exclusión y la inequidad''}, \textbf{``economía que mata''}, \textbf{``idolatría del dinero''} (\textit{Evangelii gaudium}). El conflicto ambiental y el social tienen un origen común.
\end{itemize}

\medskip\noindent \textbf{Tercera clave: personalismo frente al populismo y al liberalismo}\par\nopagebreak
\begin{itemize}
\item \textit{Fratelli tutti} 155: el \textbf{desprecio de los débiles} puede esconderse en formas \textbf{populistas} (los usan demagógicamente) o \textbf{liberales} (al servicio de los intereses económicos de los poderosos). Ambas son expresiones del \textbf{individualismo}.
\item Dos aportes: (a) reflexión sobre \textbf{la política} (crítica a la polarización y al marketing; ``mejor política''; bien común, subsidiariedad, solidaridad), en continuidad con la tradición; (b) la \textbf{centralidad del pueblo}, muy propia de él.
\item \textbf{Pueblo:} categoría \textbf{``mítica, no mística''}: identidad común hecha de lazos sociales y culturales, proceso lento hacia un proyecto común. Ligada a la \textbf{teología del pueblo} argentina. Es también concepto teológico: el \textbf{``pueblo fiel''}, pueblo de Dios. ``Nadie se salva solo.''
\item El \textbf{populismo} instrumentaliza al pueblo al servicio de un líder; se distingue de lo \textbf{popular}.
\end{itemize}

\medskip\noindent \textbf{Cuarta clave: la realidad vista desde los últimos}\par\nopagebreak
\begin{itemize}
\item Opción preferencial por los pobres, con tres notas:
\begin{enumerate}
\item Los pobres son los \textbf{excluidos, los descartados}: ya no ``explotados'' sino \textbf{``sobrantes''}, están \textbf{fuera}. \textbf{Cultura del descarte}.
\item Los pobres \textbf{``tienen mucho que enseñarnos''}: ``el mundo se ve más claramente \textbf{desde las periferias}''. No es paternalismo.
\item Son \textbf{agentes de su propio desarrollo} y promotores del cambio social.
\end{enumerate}
\item \textbf{Movimientos populares:} encuentros mundiales (Roma 2014, Santa Cruz de la Sierra 2015, Roma 2016, virtual 2021). Francisco los llama \textbf{``poetas sociales''}: crean esperanza donde hay descarte. La sociedad se construye \textbf{desde abajo}.
\end{itemize}

\medskip\noindent \textbf{Quinta clave: la fraternidad como sueño en un mundo globalizado}\par\nopagebreak
\begin{itemize}
\item La fraternidad completa la tríada \textbf{libertad, igualdad, fraternidad}: sin fraternidad la libertad se reduce a soledad y autonomía, y la igualdad queda en abstracción.
\item Es un \textbf{sueño ``con los pies en la tierra''}: fraternidad universal que se construye \textbf{desde lo cercano}, desde abajo, no por decreto.
\item Diagnóstico: ``indiferencia cómoda, fría y globalizada''; olvido de que ``estamos todos en la misma barca''.
\item Caminos: \textbf{diálogo} (``acercarse, expresarse, escucharse, mirarse, conocerse...''), reconciliación, acogida de \textbf{migrantes} (Lampedusa, 2013), respeto de la diversidad cultural.
\item Tarea común de las \textbf{religiones}: Documento sobre la fraternidad humana (\textbf{Abu Dabi}, 2019, con el Gran Imán Al-Tayyeb).
\end{itemize}

\medskip\noindent \textbf{Sexta clave: los cristianos en un mundo laico, plural y religiosamente diverso}\par\nopagebreak
\begin{itemize}
\item \textbf{Doble discurso:} exhortaciones (\textit{Evangelii gaudium}, \textit{Amoris laetitia}) dirigidas a la comunidad cristiana; encíclicas (\textit{Laudato si'}, \textit{Fratelli tutti}) dirigidas a toda la humanidad.
\item Los documentos terminan con un capítulo sobre \textbf{espiritualidad}: el pensamiento debe movilizar a los creyentes para la acción.
\item Santidad para todos (\textit{Gaudete et exsultate}).
\item \textbf{Sinodalidad:} no solo un evento sino un \textbf{estilo de vida} de la Iglesia: ``caminar juntos''; no es democracia ni mero consenso, sino \textbf{escucha del Espíritu}.
\end{itemize}

\subsubsection{Ideas clave para retener (Camacho)}
\begin{enumerate}
\item Vaticano II: de la confrontación al \textbf{diálogo} con la Modernidad; Iglesia pueblo de Dios; laicos protagonistas.
\item Francisco: método \textbf{inductivo}, estilo pastoral, primacía de los \textbf{gestos}.
\item Seis claves: tiempo sobre espacio / personalismo vs. individualismo / contra populismo y liberalismo / realidad desde los últimos / fraternidad / cristianos en un mundo plural.
\item Cuatro principios de \textit{Evangelii gaudium} e imagen del \textbf{poliedro}.
\item Vocabulario propio: Iglesia en salida, cultura del descarte, periferias, poetas sociales, economía que mata, paradigma tecnocrático, antropocentrismo desviado, pueblo, sinodalidad.
\end{enumerate}



\subsection{5.2 Francisco --- \textit{Laudato Si'} (2015), Cap. IV: ``Una ecología integral''}

\noindent \textbf{Idea central:} porque \textbf{``todo está conectado''}, los problemas actuales requieren una mirada que tenga en cuenta todos los factores de la crisis mundial: una \textbf{ecología integral}, que incorpore claramente las \textbf{dimensiones humanas y sociales}.\par\smallskip

\subsubsection{I. Ecología ambiental, económica y social (\S{}138--142)}
\begin{itemize}
\item La ecología estudia las relaciones entre los organismos y su ambiente; exige pensar las condiciones de vida de una sociedad y cuestionar modelos de desarrollo, producción y consumo.
\item ``Ambiente'' es la relación entre la naturaleza y la sociedad que la habita: somos parte de ella. Los ecosistemas tienen un \textbf{valor intrínseco}, independiente de su uso.
\item \textbf{``No hay dos crisis separadas, una ambiental y otra social, sino una sola y compleja crisis socio-ambiental.''} La solución requiere una aproximación integral para \textbf{combatir la pobreza, devolver la dignidad a los excluidos y cuidar la naturaleza}.
\item \textbf{Ecología económica:} el crecimiento debe considerar la realidad de modo más amplio; la protección del ambiente es parte del desarrollo.
\item \textbf{Ecología social:} es necesariamente \textbf{institucional} (familia, comunidad local, nación, vida internacional); la debilidad de las instituciones daña el ambiente.
\end{itemize}

\subsubsection{II. Ecología cultural (\S{}143--146)}
\begin{itemize}
\item Junto con el patrimonio natural hay un \textbf{patrimonio histórico, artístico y cultural} amenazado.
\item La visión consumista y la economía globalizada tienden a \textbf{homogeneizar las culturas}. Las \textbf{soluciones meramente técnicas} o las normativas uniformes corren el riesgo de atender síntomas que no responden a los problemas más profundos; hace falta el protagonismo de los actores locales. La desaparición de una cultura puede ser tan grave como la de una especie.
\item Especial atención a las \textbf{comunidades aborígenes}: para ellas la tierra no es un bien económico sino don de Dios y de los antepasados, espacio sagrado. Deben ser los \textbf{principales interlocutores}.
\end{itemize}

\subsubsection{III. Ecología de la vida cotidiana (\S{}147--155)}
\begin{itemize}
\item La calidad de vida depende de los ambientes donde se vive: vivienda, barrio, trabajo, \textbf{transporte}, espacios públicos.
\item Creatividad de los pobres para humanizar ambientes precarios; la \textbf{falta de vivienda} es grave; integrar los barrios precarios en vez de erradicarlos.
\item \textbf{Ecología humana:} existe una ``ley moral inscrita en la naturaleza humana''. \textbf{El cuerpo} nos pone en relación directa con el ambiente y los demás: \textbf{aceptarlo como don de Dios}, cuidarlo y respetar sus significados es necesario para acoger el mundo entero como don; valorar el propio cuerpo en su \textbf{femineidad o masculinidad}. No es sana una actitud que pretenda ``cancelar la diferencia sexual''.
\end{itemize}

\subsubsection{IV. El principio del bien común (\S{}156--158)}
\begin{itemize}
\item La ecología integral es \textbf{inseparable de la noción de bien común} (definición de \textit{Gaudium et spes}: conjunto de condiciones de la vida social que hacen posible a las asociaciones y a cada miembro el logro más pleno de la propia perfección).
\item Supone: respeto a la \textbf{persona} y sus derechos; bienestar social y desarrollo de los \textbf{grupos intermedios} (subsidiariedad), en especial la \textbf{familia}; \textbf{paz social} (estabilidad y seguridad), que requiere \textbf{justicia distributiva}.
\item En las condiciones actuales de inequidad, el bien común se convierte en un llamado a la \textbf{solidaridad} y en una \textbf{opción preferencial por los más pobres}.
\end{itemize}

\subsubsection{V. Justicia entre las generaciones (\S{}159--162)}
\begin{itemize}
\item \textbf{``La noción de bien común incorpora también a las generaciones futuras.''} No se puede hablar de desarrollo sostenible sin \textbf{solidaridad intergeneracional}.
\item La tierra que recibimos \textbf{pertenece también a los que vendrán}: es un préstamo. No es una actitud opcional sino \textbf{una cuestión básica de justicia}.
\item ``¿Qué tipo de mundo queremos dejar a quienes nos sucedan?'' La pregunta toca el \textbf{sentido de la existencia} y nuestra propia dignidad.
\item \textbf{``El ritmo de consumo, de desperdicio y de alteración del medio ambiente ha superado las posibilidades del planeta, de tal manera que el estilo de vida actual, por ser insostenible, sólo puede terminar en catástrofes.''}
\item La dificultad de tomar en serio el desafío se relaciona con un \textbf{deterioro ético y cultural}: individualismo, búsqueda egoísta de satisfacción inmediata. Además de la solidaridad intergeneracional urge la \textbf{solidaridad intrageneracional} (con los pobres de hoy).
\end{itemize}

\subsubsection{Ideas clave para retener (Laudato Si', cap. IV)}
\begin{enumerate}
\item ``Todo está conectado'': \textbf{una sola crisis socio-ambiental}.
\item Ecología integral = ambiental + económica + social + cultural + de la vida cotidiana (y humana).
\item Toda verdadera ecología incluye las \textbf{dimensiones humanas y sociales}: clamor de la tierra y clamor de los pobres.
\item El cuerpo como \textbf{don} que hay que aceptar y respetar.
\item Bien común + \textbf{solidaridad intergeneracional} (justicia, no opción).
\item El estilo de vida actual es \textbf{insostenible} y solo puede terminar en catástrofes.
\end{enumerate}



\subsection{5.3 Francisco --- \textit{Fratelli tutti} (2020), Cap. I y VII}

\subsubsection{CAPÍTULO I --- Las sombras de un mundo cerrado (\S{}9--55)}

\noindent \textbf{Idea central:} un diagnóstico de las tendencias del mundo actual que \textbf{desfavorecen la fraternidad universal}.\par\smallskip

\begin{itemize}
\item \textbf{Sueños que se rompen en pedazos (\S{}10--12):} retroceden los sueños de integración (Europa unida, América Latina); resurgen \textbf{nacionalismos cerrados}, resentidos y agresivos. \textbf{``Abrirse al mundo''} es una expresión que \textbf{fue cooptada por la economía y las finanzas}: hoy se refiere exclusivamente a la apertura a intereses extranjeros o a la libertad de los poderes económicos para invertir sin trabas. La sociedad globalizada ``nos hace más cercanos, pero no más hermanos''.
\item \textbf{El fin de la conciencia histórica (\S{}13--14):} ``deconstruccionismo'' que desprecia la historia; \textbf{colonización cultural}; se vacían de sentido palabras como democracia, libertad, justicia, unidad.
\item \textbf{Sin un proyecto para todos (\S{}15--17):} la política se reduce a \textbf{exasperar, exacerbar y polarizar}; marketing e inmediatismo; se niega al otro el derecho a existir y opinar.
\item \textbf{El descarte mundial (\S{}18--21):} partes de la humanidad parecen \textbf{sacrificables}: los no nacidos, los \textbf{ancianos} (abandonados en la pandemia), los pobres, las personas con discapacidad. La riqueza aumenta, pero con inequidad; \textbf{racismo} que se oculta y reaparece.
\item \textbf{Derechos humanos no suficientemente universales (\S{}22--24):} en la práctica no son iguales para todos; desigualdad de las \textbf{mujeres}; formas modernas de \textbf{esclavitud} y trata.
\item \textbf{Conflicto y miedo (\S{}25--28):} \textbf{``tercera guerra mundial en etapas''} (en \S{}259: ``guerra mundial a pedazos''); cultura de \textbf{muros}; soledad y miedo que alimentan las mafias.
\item \textbf{Globalización y progreso sin un rumbo común (\S{}29--31):} avances científicos junto a deterioro ético; ``indiferencia cómoda, fría y globalizada''.
\item \textbf{Las pandemias y otros flagelos de la historia (\S{}32--36):} el Covid-19 mostró que somos una comunidad mundial en \textbf{la misma barca}: \textbf{``nadie se salva solo''}. Riesgo del ``sálvese quien pueda''.
\item \textbf{Sin dignidad humana en las fronteras (\S{}37--41):} tanto regímenes populistas como planteos económicos liberales quieren evitar la llegada de \textbf{migrantes}, tratados como menos valiosos, menos humanos; mentalidad xenófoba incompatible con la fe. También hay un \textbf{derecho a no emigrar} (tener condiciones para quedarse en la propia tierra).
\item \textbf{La ilusión de la comunicación (\S{}42--50):} el mundo digital expone y vigila la intimidad; \textbf{agresividad sin pudor} en las redes; \textbf{información sin sabiduría}: acumulación de datos no es sabiduría; se pierde el silencio, la escucha y el encuentro real.
\item \textbf{Sometimientos y autodesprecios (\S{}51--53):} países que copian modelos ajenos con baja autoestima; ``no hay peor alienación que experimentar que no se tienen raíces''.
\item \textbf{Esperanza (\S{}54--55):} a pesar de las sombras, Dios sigue sembrando bien (los trabajadores esenciales en la pandemia). Invitación a la \textbf{esperanza}, que es audaz.
\end{itemize}

\subsubsection{CAPÍTULO VII --- Caminos de reencuentro (\S{}225--270)}

\noindent \textbf{Idea central:} en muchos lugares hacen falta caminos de paz que cicatricen las heridas: \textbf{artesanos de paz} dispuestos a generar procesos de sanación y reencuentro con ingenio y audacia.\par\smallskip

\medskip\noindent \textbf{Recomenzar desde la verdad (\S{}226--227)}\par\nopagebreak
\begin{itemize}
\item Reencuentro no es volver a antes del conflicto. Los que han estado enfrentados deben conversar \textbf{desde la verdad}, clara y desnuda, cultivando una \textbf{``memoria penitencial''}. ``El pueblo tiene el derecho de saber qué pasó.''
\item \textbf{Verdad, justicia y misericordia} son inseparables. Verdad es contar a las familias lo ocurrido con sus desaparecidos, reconocer el dolor de las víctimas.
\end{itemize}

\medskip\noindent \textbf{La arquitectura y la artesanía de la paz (\S{}228--232)}\par\nopagebreak
\begin{itemize}
\item El camino a la paz no implica homogeneizar la sociedad.
\item \textbf{``Arquitectura'' de la paz:} intervienen las instituciones de la sociedad. \textbf{``Artesanía'' de la paz:} nos involucra a todos; los procesos efectivos son transformaciones artesanales obradas por los pueblos.
\item Nadie puede ser encasillado para siempre por lo que hizo; no hay punto final en la construcción de la paz.
\end{itemize}

\medskip\noindent \textbf{Sobre todo con los últimos (\S{}233--235)}\par\nopagebreak
\begin{itemize}
\item La paz social no es solo ausencia de violencia. \textbf{Sin igualdad de oportunidades}, la agresión y la guerra encuentran caldo de cultivo. Opción por los más pobres.
\end{itemize}

\medskip\noindent \textbf{El valor y el sentido del perdón (\S{}236--245)}\par\nopagebreak
\begin{itemize}
\item \textbf{El conflicto es inevitable}; no hay que ocultarlo ni proponer una falsa ``reconciliación'' que beneficie al poderoso.
\item \textbf{Luchas legítimas y perdón:} perdonar \textbf{no es renunciar a los propios derechos} frente a un poderoso corrupto o un criminal. Amar al opresor no es consentir que siga oprimiendo: es buscar que deje de hacerlo. Se exige justicia \textbf{sin venganza ni odio}.
\item \textbf{La verdadera superación:} el conflicto se resuelve en un plano superior que conserva las virtualidades de las polaridades: \textbf{``la unidad es superior al conflicto''}.
\end{itemize}

\medskip\noindent \textbf{La memoria (\S{}246--249)}\par\nopagebreak
\begin{itemize}
\item No se puede exigir un ``perdón social'' a quien sufrió: la reconciliación es un hecho personal y no se puede decretar una ``reconciliación general''. Nunca proponer el olvido.
\item Hay que \textbf{mantener viva la memoria}: la \textbf{Shoah}, \textbf{Hiroshima y Nagasaki}, las persecuciones, la trata de esclavos, las matanzas étnicas. No ``dar vuelta la página''. Recordar también a quienes eligieron el bien.
\end{itemize}

\medskip\noindent \textbf{Perdón sin olvidos (\S{}250--254)}\par\nopagebreak
\begin{itemize}
\item \textbf{Perdonar no es olvidar.} El perdón libre y sincero renuncia a ser poseído por la misma fuerza destructiva (romper el círculo de la venganza), pero \textbf{no implica impunidad}: la justicia se busca por amor a la justicia misma, por respeto a las víctimas y para prevenir nuevos crímenes.
\item La violencia ejercida desde las estructuras y el poder del \textbf{Estado} no está al mismo nivel que la de grupos particulares; pero toda víctima inocente merece el mismo respeto.
\end{itemize}

\medskip\noindent \textbf{La guerra y la pena de muerte (\S{}255--270)}\par\nopagebreak
\begin{itemize}
\item Dos situaciones extremas presentadas como soluciones, que son \textbf{respuestas falsas}.
\item \textbf{La injusticia de la guerra:} ``toda guerra deja al mundo peor''; es \textbf{``un fracaso de la política y de la humanidad''}. Con las armas actuales los riesgos superan siempre la utilidad: \textbf{hoy es muy difícil sostener los criterios de la ``guerra justa''}. \textbf{``¡Nunca más la guerra!''}
\item Las \textbf{armas nucleares}: su eliminación total es imperativo moral y humanitario. Propuesta: con el dinero de las armas, un \textbf{Fondo mundial} contra el hambre.
\item \textbf{La pena de muerte:} es \textbf{inadmisible}; la Iglesia se compromete a su \textbf{abolición en todo el mundo} (en línea con Juan Pablo II). El asesino no pierde su dignidad personal. También rechaza la \textbf{cadena perpetua} (``pena de muerte oculta'') y las ejecuciones extrajudiciales.
\end{itemize}

\subsubsection{Ideas clave para retener (Fratelli tutti I y VII)}
\begin{enumerate}
\item Cap. I: sombras = nacionalismos cerrados, ``abrirse al mundo'' reducido a lo económico, polarización, \textbf{descarte}, derechos no universales, guerra a pedazos, muros, trato a migrantes, ilusión de la comunicación digital. Cierre en \textbf{esperanza}.
\item Lección de la pandemia: \textbf{``nadie se salva solo''}.
\item Cap. VII: paz = \textbf{verdad + justicia + misericordia}; \textbf{arquitectura} (instituciones) y \textbf{artesanía} (todos) de la paz.
\item \textbf{Perdón sin olvido y sin impunidad}; memoria (Shoah, Hiroshima).
\item \textbf{Guerra:} fracaso de la política y de la humanidad; casi imposible hablar de ``guerra justa''. \textbf{Pena de muerte:} inadmisible.
\end{enumerate}

\subsubsection{Síntesis del Eje 5}
\begin{esquema}
\noindent\hspace*{0.0em}MÉTODO DE FRANCISCO: inductivo (ver la realidad) + gestos + discernimiento\par
\noindent\hspace*{1.7em}DIAGNÓSTICO: individualismo $\rightarrow$ paradigma tecnocrático $\rightarrow$ descarte\par
\noindent\hspace*{8.8em}(Fratelli tutti I, Laudato Si' III)\par
\noindent\hspace*{1.7em}CRITERIO:    persona y pueblo; mirar desde los últimos y las periferias\par
\noindent\hspace*{1.7em}PROPUESTA:   ecología integral (Laudato Si' IV)\par
\noindent\hspace*{8.8em}fraternidad y amistad social (Fratelli tutti)\par
\noindent\hspace*{8.8em}mejor política, diálogo, caminos de reencuentro (V, VI, VII)\par
\noindent\hspace*{1.7em}ESTILO:      iniciar procesos (tiempo \textgreater{} espacio), poliedro, sinodalidad\par
\end{esquema}



\clearpage\section{Cuadro final --- Documentos y fórmulas para reconocer en el examen}

{\small\renewcommand{\arraystretch}{1.25}
\begin{longtable}{>{\raggedright\arraybackslash}p{0.163\linewidth}>{\raggedright\arraybackslash}p{0.275\linewidth}>{\raggedright\arraybackslash}p{0.472\linewidth}}
\toprule
\textbf{Documento} & \textbf{Autor y año} & \textbf{Fórmula o tema que lo identifica} \\ \midrule \endhead
\textit{Rerum novarum} & León XIII, 1891 & Cuestión obrera; primera encíclica social \\ \addlinespace[2pt]
\textit{Quadragesimo anno} & Pío XI, 1931 & Principio de subsidiariedad \\ \addlinespace[2pt]
\textit{Mater et magistra} & Juan XXIII, 1961 & Socialización; cuestión social mundial \\ \addlinespace[2pt]
\textit{Pacem in terris} & Juan XXIII, 1963 & Derechos humanos; ``a todos los hombres de buena voluntad'' \\ \addlinespace[2pt]
\textit{Gaudium et spes} & Vaticano II, 1965 & Iglesia solidaria del género humano; definición de bien común \\ \addlinespace[2pt]
\textit{Dignitatis humanae} & Vaticano II, 1965 & Libertad religiosa \\ \addlinespace[2pt]
\textit{Populorum progressio} & Pablo VI, 1967 & ``El desarrollo es el nuevo nombre de la paz''; todo el hombre y todos los hombres \\ \addlinespace[2pt]
\textit{Octogesima adveniens} & Pablo VI, 1971 & Insuficiencia de las ideologías; nuevos problemas \\ \addlinespace[2pt]
\textit{Evangelii nuntiandi} & Pablo VI, 1975 & Evangelización y promoción humana \\ \addlinespace[2pt]
\textit{Laborem exercens} & Juan Pablo II, 1981 & El trabajo, clave de la cuestión social; dimensión subjetiva \\ \addlinespace[2pt]
\textit{Sollicitudo rei socialis} & Juan Pablo II, 1987 & Solidaridad como virtud; estructuras de pecado; la DSI es teología moral \\ \addlinespace[2pt]
\textit{Centesimus annus} & Juan Pablo II, 1991 & Juicio sobre el capitalismo; democracia y valores \\ \addlinespace[2pt]
Compendio de la DSI & Pontificio Consejo Justicia y Paz, 2004 & ``Humanismo integral y solidario'' \\ \addlinespace[2pt]
\textit{Deus caritas est} & Benedicto XVI, 2005 & Justicia y caridad \\ \addlinespace[2pt]
\textit{Caritas in veritate} & Benedicto XVI, 2009 & Cuestión antropológica; ideología tecnocrática \\ \addlinespace[2pt]
\textit{Evangelii gaudium} & Francisco, 2013 & Iglesia en salida; economía que mata; cuatro principios \\ \addlinespace[2pt]
\textit{Laudato si'} & Francisco, 2015 & Casa común; paradigma tecnocrático; ecología integral \\ \addlinespace[2pt]
\textit{Fratelli tutti} & Francisco, 2020 & Fraternidad y amistad social; mejor política \\ \addlinespace[2pt]
\textit{Dignitas infinita} & Dicasterio para la Doctrina de la Fe, 2024 & Cuatro sentidos de dignidad; violaciones graves \\ \addlinespace[2pt]
\textit{Dilexit nos} & Francisco, 2024 & Volver al corazón \\ \addlinespace[2pt]
\textit{Dilexi te} & León XIV, 2025 & Amor a los pobres; los pobres como sujetos \\ \addlinespace[2pt]
\bottomrule\end{longtable}}

\end{document}
